\input{preamble}

% OK, start here.
%
\begin{document}

\title{Espaces algébriques décents}


\maketitle

\phantomsection
\label{section-phantom}

\tableofcontents

\section{Introduction}
\label{section-introduction}

\noindent
Dans ce chapitre, nous étudions les propriétés « locales » des espaces
algébriques généraux, c'est-à-dire de ceux qui ne sont pas quasi-séparés.
Les espaces algébriques quasi-séparés sont étudiés dans \cite{Kn}.
Il apparaît que des phénomènes essentiellement nouveaux se produisent,
notamment en ce qui concerne les points et leurs spécialisations, pour les
espaces algébriques plus généraux. D'autre part, dans la plupart des résultats
fondamentaux sur les espaces algébriques, il n'est pas nécessaire de se
préoccuper de ces phénomènes. C'est pourquoi nous avons choisi de réunir cette
matière dans un chapitre séparé, après l'exposé habituel de la théorie.



\section{Conventions}
\label{section-conventions}

\noindent
Nous supposons une fois pour toutes que tous les schémas sont contenus dans
un grand site fppf $\Sch_{fppf}$. De plus, tout anneau $A$ considéré est tel
que $\Spec(A)$ est (isomorphe à) un objet de ce grand site.

\medskip\noindent
Soit $S$ un schéma et soit $X$ un espace algébrique sur $S$.
Dans ce chapitre et le suivant, nous noterons $X \times_S X$ le produit de $X$
par lui-même (dans la catégorie des espaces algébriques sur $S$), au lieu de
$X \times X$.



\section{Fibres universellement bornées}
\label{section-universally-bounded}

\noindent
Nous examinons brièvement ce que signifie, pour un morphisme d'un schéma vers
un espace algébrique, le fait d'avoir des fibres universellement bornées. Voir
Morphismes, section \ref{morphisms-section-universally-bounded}
pour les définitions et résultats analogues concernant les morphismes de
schémas.

\begin{definition}
\label{definition-universally-bounded}
Soit $S$ un schéma. Soit $X$ un espace algébrique sur $S$, et soit
$U$ un schéma sur $S$. Soit $f : U \to X$ un morphisme sur $S$.
Nous dirons que {\it les fibres de $f$ sont universellement
bornées}\footnote{Cette terminologie n'est probablement pas standard.}
s'il existe un entier $n$ tel que, pour tout corps $k$ et tout morphisme
$\Spec(k) \to X$, le produit fibré $\Spec(k) \times_X U$ soit un schéma fini
sur $k$, dont le degré sur $k$ est $\leq n$.
\end{definition}

\noindent
Cette définition a un sens, car le produit fibré
$\Spec(k) \times_X U$ est un schéma. De plus, si $X$ est un schéma,
on retrouve la notion de
Morphismes, définition \ref{morphisms-definition-universally-bounded}
en vertu de
Morphismes, lemme \ref{morphisms-lemma-characterize-universally-bounded}.

\begin{lemma}
\label{lemma-composition-universally-bounded}
Soit $S$ un schéma. Soit $X$ un espace algébrique sur $S$.
Soit $V \to U$ un morphisme de schémas sur $S$, et soit
$U \to X$ un morphisme de $U$ vers $X$. Si les fibres de
$V \to U$ et de $U \to X$ sont universellement bornées, alors celles de
$V \to X$ le sont aussi.
\end{lemma}

\begin{proof}
Soit $n$ un entier qui convient à $V \to U$, et soit $m$ un entier qui
convient à $U \to X$ dans la
Définition \ref{definition-universally-bounded}.
Soit $\Spec(k) \to X$ un morphisme, où $k$ est un corps.
Considérons les morphismes
$$
\Spec(k) \times_X V
\longrightarrow
\Spec(k) \times_X U
\longrightarrow
\Spec(k).
$$
Par hypothèse, le schéma $\Spec(k) \times_X U$ est fini de degré au plus $m$
sur $k$, et $n$ borne le degré des fibres du premier morphisme. Par
Morphismes, lemme \ref{morphisms-lemma-composition-universally-bounded},
nous en déduisons que $\Spec(k) \times_X V$ est fini sur $k$ de degré au plus
$nm$.
\end{proof}

\begin{lemma}
\label{lemma-base-change-universally-bounded}
Soit $S$ un schéma.
Soit $Y \to X$ un morphisme représentable d'espaces algébriques sur $S$.
Soit $U \to X$ un morphisme d'un schéma vers $X$.
Si les fibres de $U \to X$ sont universellement bornées, alors celles de
$U \times_X Y \to Y$ le sont aussi.
\end{lemma}

\begin{proof}
Cela résulte immédiatement de la définition et des propriétés des produits
fibrés. (Notons que $U \times_X Y$ est un schéma puisque nous avons supposé
$Y \to X$ représentable, si bien que la définition s'applique.)
\end{proof}

\begin{lemma}
\label{lemma-descent-universally-bounded}
Soit $S$ un schéma. Soit $g : Y \to X$ un morphisme représentable d'espaces
algébriques sur $S$. Soit $f : U \to X$ un morphisme d'un schéma vers $X$.
Soit $f' : U \times_X Y \to Y$ le changement de base de $f$.
Si
$$
\Im(|f| : |U| \to |X|) \subset \Im(|g| : |Y| \to |X|)
$$
et si $f'$ a des fibres universellement bornées, alors $f$ a des fibres
universellement bornées.
\end{lemma}

\begin{proof}
Soit $n \geq 0$ un entier qui borne les degrés des produits fibrés
$\Spec(k) \times_Y (U \times_X Y)$ au sens de la
Définition \ref{definition-universally-bounded} appliquée au morphisme $f'$.
Nous affirmons que $n$ convient aussi à $f$. En effet, supposons que
$x : \Spec(k) \to X$ soit un morphisme depuis le spectre d'un corps.
Alors, soit $\Spec(k) \times_X U$ est vide (et il n'y a rien à démontrer),
soit $x$ appartient à l'image de $|f|$. D'après
Propriétés des espaces,
lemme \ref{spaces-properties-lemma-points-cartesian}
et l'hypothèse du lemme, cela signifie qu'il existe une extension de corps
$k'/k$ et un diagramme commutatif
$$
\xymatrix{
\Spec(k') \ar[r] \ar[d] & Y \ar[d] \\
\Spec(k) \ar[r] & X
}
$$
On obtient alors
$$
\Spec(k') \times_Y (U \times_X Y) =
\Spec(k') \times_{\Spec(k)} (\Spec(k) \times_X U)
$$
Comme le schéma $\Spec(k') \times_Y (U \times_X Y)$ est supposé fini de degré
$\leq n$ sur $k'$, il en résulte que $\Spec(k) \times_X U$ est lui aussi fini
de degré $\leq n$ sur $k$, comme voulu. (Nous omettons quelques détails.)
\end{proof}

\begin{lemma}
\label{lemma-universally-bounded-permanence}
Soit $S$ un schéma. Soit $X$ un espace algébrique sur $S$.
Considérons un diagramme commutatif
$$
\xymatrix{
U \ar[rd]_g \ar[rr]_f & & V \ar[ld]^h \\
& X &
}
$$
où $U$ et $V$ sont des schémas. Si $g$ a des fibres universellement bornées
et si $f$ est surjectif et plat, alors $h$ a lui aussi des fibres
universellement bornées.
\end{lemma}

\begin{proof}
Supposons que $g$ ait des fibres universellement bornées et que $f$ soit
surjectif et plat. Soit $n \geq 0$ un entier qui borne les degrés des schémas
$\Spec(k) \times_X U$ au sens de la
Définition \ref{definition-universally-bounded}.
Nous affirmons que $n$ convient aussi à $h$.
Soit $\Spec(k) \to X$ un morphisme du spectre d'un corps vers $X$.
Considérons le morphisme de schémas
$$
\Spec(k) \times_X U \longrightarrow \Spec(k) \times_X V
$$
Il est plat et surjectif. Par hypothèse, le schéma de gauche est fini de degré
$\leq n$ sur $\Spec(k)$. Il résulte de
Morphismes, lemme \ref{morphisms-lemma-universally-bounded-permanence}
que le degré du schéma de droite est lui aussi majoré par $n$, comme voulu.
\end{proof}

\begin{lemma}
\label{lemma-universally-bounded-finite-fibres}
Soit $S$ un schéma.
Soit $X$ un espace algébrique sur $S$, et soit $U$ un schéma sur $S$.
Soit $\varphi : U \to X$ un morphisme sur $S$.
Si les fibres de $\varphi$ sont universellement bornées, il existe un entier
$n$ tel que chaque fibre de $|U| \to |X|$ compte au plus $n$ éléments.
\end{lemma}

\begin{proof}
L'entier $n$ de la Définition \ref{definition-universally-bounded} convient.
En effet, choisissons $x \in |X|$. Représentons $x$ par un morphisme
$x : \Spec(k) \to X$. Nous obtenons alors un diagramme commutatif
$$
\xymatrix{
\Spec(k) \times_X U \ar[r] \ar[d] & U \ar[d] \\
\Spec(k) \ar[r]^x & X
}
$$
qui montre (au moyen de
Propriétés des espaces,
lemme \ref{spaces-properties-lemma-points-cartesian})
que l'image réciproque de $x$ dans $|U|$ est l'image de la flèche horizontale
supérieure. Puisque $\Spec(k) \times_X U$ est fini de degré $\leq n$ sur $k$,
il possède au plus $n$ points.
\end{proof}











\section{Conditions de finitude et points}
\label{section-points-monomorphisms}

\noindent
Dans cette section, nous approfondissons la question de savoir quand des
points peuvent être représentés par des monomorphismes de spectres de corps
dans l'espace.

\begin{remark}
\label{remark-recall}
Avant de démontrer le lemme suivant, rappelons quelques faits sur les
morphismes étales de schémas :
\begin{enumerate}
\item Un morphisme étale est plat; les générisations se relèvent donc le long
d'un morphisme étale
(Morphismes, lemmes \ref{morphisms-lemma-etale-flat}
et \ref{morphisms-lemma-generalizations-lift-flat}).
\item Un morphisme étale est non ramifié; un morphisme non ramifié est
localement quasi-fini; ses fibres sont donc discrètes
(Morphismes, lemmes \ref{morphisms-lemma-flat-unramified-etale},
\ref{morphisms-lemma-unramified-quasi-finite}, et
\ref{morphisms-lemma-quasi-finite-at-point-characterize}).
\item Un morphisme étale quasi-compact est quasi-fini et possède en particulier
des fibres finies
(Morphismes, lemmes \ref{morphisms-lemma-quasi-finite-locally-quasi-compact} et
\ref{morphisms-lemma-quasi-finite}).
\item Un schéma étale sur un corps $k$ est une réunion disjointe de spectres
d'extensions finies séparables de $k$
(Morphismes, lemme \ref{morphisms-lemma-etale-over-field}).
\end{enumerate}
Pour une étude générale des morphismes étales, voir
\'Morphismes étales, section \ref{etale-section-etale-morphisms}.
\end{remark}

\begin{lemma}
\label{lemma-U-finite-above-x}
Soit $S$ un schéma. Soit $X$ un espace algébrique sur $S$.
Soit $x \in |X|$. Les conditions suivantes sont équivalentes :
\begin{enumerate}
\item il existe une famille de schémas $U_i$ et de morphismes étales
$\varphi_i : U_i \to X$ telle que
$\coprod \varphi_i : \coprod U_i \to X$ soit surjectif,
et telle que, pour tout $i$, la fibre de
$|U_i| \to |X|$ au-dessus de $x$ soit finie, et
\item pour tout schéma affine $U$ et tout morphisme étale $\varphi : U \to X$,
la fibre de $|U| \to |X|$ au-dessus de $x$ est finie.
\end{enumerate}
\end{lemma}

\begin{proof}
L'implication (2) $\Rightarrow$ (1) est immédiate.
Soit $\varphi_i : U_i \to X$ une famille de morphismes étales comme en (1).
Soit $\varphi : U \to X$ un morphisme étale d'un schéma affine vers $X$.
Considérons les diagrammes de produits fibrés
$$
\xymatrix{
U \times_X U_i \ar[r]_-{p_i} \ar[d]_{q_i} & U_i \ar[d]^{\varphi_i} \\
U \ar[r]^\varphi & X
}
\quad \quad
\xymatrix{
\coprod U \times_X U_i \ar[r]_-{\coprod p_i} \ar[d]_{\coprod q_i} &
\coprod U_i \ar[d]^{\coprod \varphi_i} \\
U \ar[r]^\varphi & X
}
$$
Comme $q_i$ est étale, il est ouvert (voir la Remarque \ref{remark-recall}).
De plus, le morphisme $\coprod q_i$ est surjectif.
Il existe donc un nombre fini d'indices $i_1, \ldots, i_n$ et des ouverts
quasi-compacts $W_{i_j} \subset U \times_X U_{i_j}$ dont les images recouvrent
$U$.
Le morphisme $p_i$ est étale, donc localement quasi-fini (voir la remarque sur
les morphismes étales ci-dessus). Nous pouvons donc appliquer
Morphismes, lemme
\ref{morphisms-lemma-locally-quasi-finite-qc-source-universally-bounded}
pour voir que les fibres de
$p_{i_j}|_{W_{i_j}} : W_{i_j} \to U_{i_j}$ sont finies.
D'après
Propriétés des espaces,
lemme \ref{spaces-properties-lemma-points-cartesian}
et l'hypothèse sur $\varphi_i$, nous en concluons que la fibre de $\varphi$
au-dessus de $x$ est finie. Autrement dit, (2) est vérifiée.
\end{proof}

\begin{lemma}
\label{lemma-R-finite-above-x}
Soit $S$ un schéma. Soit $X$ un espace algébrique sur $S$.
Soit $x \in |X|$. Les conditions suivantes sont équivalentes :
\begin{enumerate}
\item il existe un schéma $U$, un morphisme étale
$\varphi : U \to X$, et des points $u, u' \in U$ dont l'image est
$x$ tels que, si l'on pose $R = U \times_X U$, la fibre de
$$
|R| \to |U| \times_{|X|} |U|
$$
au-dessus de $(u, u')$ soit finie,
\item pour tout schéma $U$, tout morphisme étale $\varphi : U \to X$ et
tous points $u, u' \in U$ dont l'image est
$x$, si l'on pose $R = U \times_X U$, la fibre de
$$
|R| \to |U| \times_{|X|} |U|
$$
au-dessus de $(u, u')$ est finie,
\item il existe un morphisme $\Spec(k) \to X$, où $k$ est un corps,
dans la classe d'équivalence de $x$, tel que les projections
$\Spec(k) \times_X \Spec(k) \to \Spec(k)$ soient
étales et quasi-compactes, et
\item il existe un monomorphisme $\Spec(k) \to X$, où $k$ est un corps,
dans la classe d'équivalence de $x$.
\end{enumerate}
\end{lemma}

\begin{proof}
Supposons (1), c'est-à-dire soit $\varphi : U \to X$ un morphisme étale
d'un schéma vers $X$, et soient $u, u'$ des points de $U$ situés au-dessus de
$x$ tels que la fibre de $|R| \to |U| \times_{|X|} |U|$ au-dessus de $(u, u')$
soit un ensemble fini. Dans cette démonstration, nous considérons un point
$u = \Spec(\kappa(u))$ comme un schéma. Notons que $u \to U$, $u' \to U$ sont
des monomorphismes (voir
Schémas, lemme \ref{schemes-lemma-injective-points-surjective-stalks}); ainsi,
$u \times_X u' \to R = U \times_X U$ est un monomorphisme.
Dans ce langage, l'hypothèse signifie exactement que
$u \times_X u'$ est un schéma dont l'espace topologique sous-jacent possède
un nombre fini de points.
Soit $\psi : W \to X$ un morphisme étale d'un schéma vers $X$.
Soient $w, w' \in W$ des points de $W$ dont l'image est $x$.
Nous devons montrer que $w \times_X w'$ est un schéma dont l'espace topologique
sous-jacent possède un nombre fini de points.
Considérons le diagramme de produit fibré
$$
\xymatrix{
W \times_X U \ar[r]_p \ar[d]_q & U \ar[d]^\varphi \\
W \ar[r]^\psi & X
}
$$
Comme $x$ est l'image de $u$ et de $u'$, nous pouvons choisir des points
$\tilde w, \tilde w'$ de $W \times_X U$ tels que $q(\tilde w) = w$,
$q(\tilde w') = w'$, $u = p(\tilde w)$ et $u' = p(\tilde w')$; voir
Propriétés des espaces,
lemme \ref{spaces-properties-lemma-points-cartesian}.
Comme $p$ et $q$ sont étales, les extensions de corps
$\kappa(w) \subset \kappa(\tilde w) \supset \kappa(u)$ et
$\kappa(w') \subset \kappa(\tilde w') \supset \kappa(u')$ sont
finies séparables; voir la Remarque \ref{remark-recall}.
Nous obtenons alors un diagramme commutatif
$$
\xymatrix{
w \times_X w' \ar[d] &
\tilde w \times_X \tilde w' \ar[l] \ar[d] \ar[r] &
u \times_X u' \ar[d] \\
w \times_X w' &
\tilde w \times_S \tilde w' \ar[l] \ar[r] &
u \times_S u'
}
$$
où les carrés sont des carrés de produits fibrés. Les morphismes
horizontaux inférieurs sont étales et quasi-compacts, car tout schéma de la
forme $\Spec(k) \times_S \Spec(k')$ est affine, et d'après nos observations
ci-dessus sur les extensions de corps.
Nous voyons donc que les flèches horizontales supérieures sont étales et
quasi-compactes, et qu'elles ont par conséquent des fibres finies.
Nous avons vu plus haut que $|u \times_X u'|$ est fini; nous en concluons que
$|w \times_X w'|$ est fini. Autrement dit, (2) est vérifiée.

\medskip\noindent
Supposons (2). Soit $U \to X$ un morphisme étale d'un schéma $U$
tel que $x$ appartienne à l'image de $|U| \to |X|$. Soit $u \in U$ un
point dont l'image est $x$. Nous avons alors vu au paragraphe précédent que
$u = \Spec(\kappa(u)) \to X$ a la propriété que
$u \times_X u$ possède un espace topologique sous-jacent fini. D'autre part,
les projections $u \times_X u \to u$ sont les composés
$$
u \times_X u \longrightarrow
u \times_X U \longrightarrow
u \times_X X = u,
$$
c'est-à-dire les composés d'un monomorphisme (le changement de base du
monomorphisme $u \to U$) et d'un morphisme étale (le changement de base du
morphisme étale $U \to X$). Ainsi, $u \times_X U$ est une réunion disjointe
de spectres de corps qui sont des extensions finies séparables de $\kappa(u)$
(voir la Remarque \ref{remark-recall}). Puisque $u \times_X u$ est fini, son
image dans $u \times_X U$ est une réunion disjointe finie de spectres de corps
qui sont des extensions finies séparables de $\kappa(u)$. Par
Schémas, lemme \ref{schemes-lemma-mono-towards-spec-field},
nous en concluons que $u \times_X u$ est une réunion disjointe finie de
spectres de corps qui sont des extensions finies séparables de $\kappa(u)$.
Autrement dit, nous voyons que $u \times_X u \to u$ est quasi-compact et
étale. Cela signifie que (3) est vérifiée.

\medskip\noindent
Démontrons que (3) implique (4). Soit $\Spec(k) \to X$ un morphisme
du spectre d'un corps vers $X$, dans la classe d'équivalence de $x$,
tel que les deux projections
$t, s : R = \Spec(k) \times_X \Spec(k)  \to \Spec(k)$
soient quasi-compactes et étales.
Cela signifie en particulier que $R$ est une relation d'équivalence étale
sur $\Spec(k)$.
Par Espaces, théorème \ref{spaces-theorem-presentation},
nous savons que le faisceau quotient
$X' = \Spec(k)/R$ est un espace algébrique. Par
Groupoïdes, lemme \ref{groupoids-lemma-quotient-groupoid-restrict},
le morphisme $X' \to X$ est un monomorphisme.
Comme $s, t$ sont quasi-compacts, nous voyons que $R$ est quasi-compact;
Propriétés des espaces,
lemme \ref{spaces-properties-lemma-point-like-spaces}
s'applique donc à $X'$, et nous voyons que
$X' = \Spec(k')$ pour un certain corps $k'$. Nous obtenons ainsi une
factorisation
$$
\Spec(k) \longrightarrow
\Spec(k') \longrightarrow X
$$
qui montre que $\Spec(k') \to X$ est un monomorphisme dont l'image est
$x \in |X|$. Autrement dit, (4) est vérifiée.

\medskip\noindent
Enfin, démontrons que (4) implique (1). Soit $\Spec(k) \to X$
un monomorphisme, où $k$ est un corps, dans la classe d'équivalence de $x$.
Soit $U \to X$ un morphisme étale surjectif d'un schéma $U$ vers $X$.
Soit $u \in U$ un point au-dessus de $x$. Puisque $\Spec(k) \times_X u$
est non vide et que $\Spec(k) \times_X u \to u$ est un monomorphisme,
nous en concluons que $\Spec(k) \times_X u = u$ (voir
Schémas, lemme \ref{schemes-lemma-mono-towards-spec-field}).
Le morphisme $u \to U \to X$ se factorise donc par $\Spec(k) \to X$; voici
le diagramme
$$
\xymatrix{
u \ar[r] \ar[d] & U \ar[d] \\
\Spec(k) \ar[r] & X
}
$$
Comme la flèche verticale de droite est étale, il en résulte que
$\kappa(u)/k$ est une extension finie séparable. Par conséquent,
$$
u \times_X u = u \times_{\Spec(k)} u
$$
est un schéma fini, et le résultat découle de l'analyse de la propriété (1)
faite au premier paragraphe de cette démonstration.
\end{proof}

\begin{lemma}
\label{lemma-weak-UR-finite-above-x}
Soit $S$ un schéma. Soit $X$ un espace algébrique sur $S$.
Soit $x \in |X|$.
Soit $U$ un schéma et soit $\varphi : U \to X$ un morphisme étale.
Les conditions suivantes sont équivalentes :
\begin{enumerate}
\item $x$ appartient à l'image de $|U| \to |X|$ et, si l'on pose
$R = U \times_X U$, les fibres des deux applications
$$
|U| \longrightarrow |X|
\quad\text{et}\quad
|R| \longrightarrow |X|
$$
au-dessus de $x$ sont finies,
\item il existe un monomorphisme $\Spec(k) \to X$, où $k$ est un corps,
dans la classe d'équivalence de $x$, et
le produit fibré $\Spec(k) \times_X U$ est
un schéma fini non vide sur $k$.
\end{enumerate}
\end{lemma}

\begin{proof}
Supposons (1). Cette condition implique clairement la première condition du
Lemme \ref{lemma-R-finite-above-x}; nous obtenons donc un monomorphisme
$\Spec(k) \to X$ dans la classe de $x$. En prenant le produit fibré,
nous voyons que $\Spec(k) \times_X U \to \Spec(k)$ est un schéma
étale sur $\Spec(k)$ ayant un nombre fini de points, donc un schéma fini
non vide sur $k$, c'est-à-dire que (2) est vérifiée.

\medskip\noindent
Supposons (2). Par hypothèse, $x$ appartient à l'image de
$|U| \to |X|$. La finitude de la fibre de
$|U| \to |X|$ au-dessus de $x$ est claire, puisque cette fibre est égale à
$|\Spec(k) \times_X U|$ d'après
Propriétés des espaces,
lemme \ref{spaces-properties-lemma-points-cartesian}.
La finitude de la fibre de $|R| \to |X|$ au-dessus de $x$ est également
claire, puisqu'elle est égale à l'ensemble sous-jacent au schéma
$$
(\Spec(k) \times_X U) \times_{\Spec(k)} (\Spec(k) \times_X U)
$$
qui est fini sur $k$. Ainsi, (1) est vérifiée.
\end{proof}

\begin{lemma}
\label{lemma-UR-finite-above-x}
Soit $S$ un schéma. Soit $X$ un espace algébrique sur $S$.
Soit $x \in |X|$. Les conditions suivantes sont équivalentes :
\begin{enumerate}
\item pour tout schéma affine $U$ et tout morphisme étale
$\varphi : U \to X$, si l'on pose $R = U \times_X U$, les fibres des deux
applications
$$
|U| \longrightarrow |X|
\quad\text{et}\quad
|R| \longrightarrow |X|
$$
au-dessus de $x$ sont finies,
\item il existe des schémas $U_i$ et des morphismes étales
$U_i \to X$ tels que $\coprod U_i \to X$ soit surjectif et que, pour tout
$i$, si l'on pose $R_i = U_i \times_X U_i$, les fibres des deux applications
$$
|U_i| \longrightarrow |X|
\quad\text{et}\quad
|R_i| \longrightarrow |X|
$$
au-dessus de $x$ soient finies,
\item il existe un monomorphisme $\Spec(k) \to X$, où $k$ est un corps,
dans la classe d'équivalence de $x$, et, pour tout schéma affine $U$ et tout
morphisme étale $U \to X$, le produit fibré $\Spec(k) \times_X U$ est
un schéma fini sur $k$,
\item il existe un monomorphisme quasi-compact $\Spec(k) \to X$,
où $k$ est un corps dans la classe d'équivalence de $x$,
\item il existe un morphisme quasi-compact $\Spec(k) \to X$,
où $k$ est un corps dans la classe d'équivalence de $x$, et
\item tout morphisme $\Spec(k) \to X$, où $k$ est un corps dans la
classe d'équivalence de $x$, est quasi-compact.
\end{enumerate}
\end{lemma}

\begin{proof}
L'équivalence de (1) et (3) résulte de l'application du
Lemme \ref{lemma-weak-UR-finite-above-x}
à tout morphisme étale $U \to X$ avec $U$ affine.
Il est clair que (3) implique (2).
Supposons que $U_i \to X$ et $R_i$ soient comme en (2). Le
Lemme \ref{lemma-U-finite-above-x}
nous permet de conclure que, pour tout schéma affine $U$ et tout morphisme
étale $U \to X$, la fibre de $|U| \to |X|$ au-dessus de $x$ est finie.
Écrivons cette fibre $\{u_1, \ldots, u_n\}$. Comme la condition (1) du
Lemme \ref{lemma-R-finite-above-x}
s'applique à $U_i \to X$ pour un certain $i$ tel que $x$ appartienne à
l'image de $|U_i| \to |X|$, nous voyons que la fibre de
$|R = U \times_X U| \to |U| \times_{|X|} |U|$
est finie au-dessus de $(u_a, u_b)$, où $a, b \in \{1, \ldots, n\}$.
La fibre de $|R| \to |X|$ au-dessus de $x$ est donc finie.
Nous voyons ainsi que (1) est vérifiée. À ce stade, nous savons que
(1), (2) et (3) sont équivalentes.

\medskip\noindent
Si (4) est vérifiée, alors, pour tout schéma affine $U$ et tout morphisme
étale $U \to X$, le schéma $\Spec(k) \times_X U$ est d'une part
étale sur $k$ (donc une réunion disjointe de spectres d'extensions finies
séparables de $k$, d'après la
Remarque \ref{remark-recall})
et d'autre part quasi-compact sur $U$ (donc quasi-compact).
Nous voyons ainsi que (3) est vérifiée.
Réciproquement, si $U_i \to X$ est comme en (2) et si $\Spec(k) \to X$
est un monomorphisme comme en (3), alors
$$
\coprod \Spec(k) \times_X U_i
\longrightarrow
\coprod U_i
$$
est quasi-compact (car, au-dessus de chaque $U_i$, nous voyons que
$\Spec(k) \times_X U_i$ est une réunion disjointe finie de spectres de corps).
Ainsi, $\Spec(k) \to X$ est quasi-compact d'après
Morphismes d'espaces, lemme \ref{spaces-morphisms-lemma-quasi-compact-local}.

\medskip\noindent
Il est immédiat que (4) implique (5). Réciproquement, soit
$\Spec(k) \to X$ un morphisme quasi-compact dans la classe d'équivalence de
$x$. Soit $U \to X$ un morphisme étale avec $U$ affine. Considérons le
produit fibré
$$
\xymatrix{
F \ar[r] \ar[d] & U \ar[d] \\
\Spec(k) \ar[r] & X
}
$$
Alors $F \to U$ est quasi-compact, donc $F$ est quasi-compact.
D'autre part, $F \to \Spec(k)$ est étale; ainsi, $F$ est une
réunion disjointe finie de spectres d'extensions finies séparables de $k$
(Remarque \ref{remark-recall}). Comme l'image de $|F| \to |U|$
est la fibre de $|U| \to |X|$ au-dessus de $x$ (Propriétés des espaces, lemme
\ref{spaces-properties-lemma-points-cartesian}), nous en concluons que
la fibre de $|U| \to |X|$ au-dessus de $x$ est finie. Le schéma
$F \times_{\Spec(k)} F$ est aussi une réunion finie de spectres de corps,
car il est lui aussi quasi-compact et étale sur $\Spec(k)$.
Il existe un monomorphisme
$F \times_X F \to F \times_{\Spec(k)} F$; par conséquent, $F \times_X F$ est
une réunion disjointe finie de spectres de corps
(Schémas, lemme \ref{schemes-lemma-mono-towards-spec-field}).
L'image de $F \times_X F \to U \times_X U = R$ est donc finie.
Comme cette image est la fibre de $|R| \to |X|$ au-dessus de $x$ d'après
Propriétés des espaces, lemme \ref{spaces-properties-lemma-points-cartesian},
nous en concluons que (1) est vérifiée. À ce stade, nous savons que
(1) -- (5) sont équivalentes.

\medskip\noindent
Il est clair que (6) implique (5). Réciproquement, supposons que
$\Spec(k) \to X$ soit comme en (4) et soit $\Spec(k') \to X$ un autre
morphisme, où $k'$ est un corps dans la classe d'équivalence de $x$. Par
Propriétés des espaces, lemme
\ref{spaces-properties-lemma-equivalence-class-point-monomorphism},
nous avons une factorisation $\Spec(k') \to \Spec(k) \to X$ du morphisme
donné. C'est un composé de morphismes quasi-compacts, donc un morphisme
quasi-compact (Morphismes d'espaces,
lemme \ref{spaces-morphisms-lemma-composition-quasi-compact}), comme voulu.
\end{proof}

\begin{lemma}
\label{lemma-U-universally-bounded}
Soit $S$ un schéma. Soit $X$ un espace algébrique sur $S$.
Les conditions suivantes sont équivalentes :
\begin{enumerate}
\item il existe des schémas $U_i$ et des morphismes étales
$U_i \to X$ tels que $\coprod U_i \to X$ soit surjectif et que
chaque $U_i \to X$ ait des fibres universellement bornées, et
\item pour tout schéma affine $U$ et tout morphisme étale $\varphi : U \to X$,
les fibres de $U \to X$ sont universellement bornées.
\end{enumerate}
\end{lemma}

\begin{proof}
L'implication (2) $\Rightarrow$ (1) est immédiate.
Supposons (1). Soit $(\varphi_i : U_i \to X)_{i \in I}$ une famille de
morphismes étales de schémas vers $X$, qui recouvre $X$, telle que
chaque $\varphi_i$ ait des fibres universellement bornées.
Soit $\psi : U \to X$ un morphisme étale d'un schéma affine vers $X$.
Pour tout $i$, considérons le diagramme de produit fibré
$$
\xymatrix{
U \times_X U_i \ar[r]_{p_i} \ar[d]_{q_i} & U_i \ar[d]^{\varphi_i} \\
U \ar[r]^\psi & X
}
$$
Comme $q_i$ est étale, il est ouvert (voir la Remarque \ref{remark-recall}).
De plus, nous avons $U = \bigcup \Im(q_i)$, puisque la famille
$(\varphi_i)_{i \in I}$ est surjective. Comme $U$ est affine, donc
quasi-compact, nous pouvons trouver un nombre fini d'indices
$i_1, \ldots, i_n \in I$ et des ouverts quasi-compacts
$W_j \subset U \times_X U_{i_j}$ tels que
$U = \bigcup q_{i_j}(W_j)$.
Le morphisme $p_{i_j}$ est étale, donc localement quasi-fini
(voir la remarque sur les morphismes étales ci-dessus). Nous pouvons donc
appliquer Morphismes, lemme
\ref{morphisms-lemma-locally-quasi-finite-qc-source-universally-bounded}
pour voir que les fibres de $p_{i_j}|_{W_j} : W_j \to U_{i_j}$ sont
universellement bornées. Par le
Lemme \ref{lemma-composition-universally-bounded},
nous voyons donc que les fibres de $W_j \to X$ sont universellement bornées.
Ainsi, $\coprod_{j = 1, \ldots, n} W_j \to X$ a lui aussi des fibres
universellement bornées. Comme $\coprod_{j = 1, \ldots, n} W_j \to X$ se
factorise par le morphisme étale surjectif
$\coprod q_{i_j}|_{W_j} : \coprod_{j = 1, \ldots, n} W_j \to U$, nous
voyons que les fibres de $U \to X$ sont universellement bornées d'après le
Lemme \ref{lemma-universally-bounded-permanence}.
Autrement dit, (2) est vérifiée.
\end{proof}

\begin{lemma}
\label{lemma-characterize-very-reasonable}
Soit $S$ un schéma.
Soit $X$ un espace algébrique sur $S$.
Les conditions suivantes sont équivalentes :
\begin{enumerate}
\item il existe un recouvrement de Zariski $X = \bigcup X_i$ et, pour
chaque $i$, un schéma $U_i$ ainsi qu'un morphisme étale surjectif
quasi-compact $U_i \to X_i$, et
\item il existe des schémas $U_i$ et des morphismes étales $U_i \to X$
tels que les projections $U_i \times_X U_i \to U_i$ soient quasi-compactes
et que $\coprod U_i \to X$ soit surjectif.
\end{enumerate}
\end{lemma}

\begin{proof}
Si (1) est vérifiée, alors les morphismes $U_i \to X_i \to X$ sont étales
(combiner Morphismes, lemme \ref{morphisms-lemma-composition-etale}
et
Espaces, lemmes
\ref{spaces-lemma-composition-representable-transformations-property} et
\ref{spaces-lemma-morphism-schemes-gives-representable-transformation-property}
).
De plus, puisque $U_i \times_X U_i = U_i \times_{X_i} U_i$,
les deux projections $U_i \times_X U_i \to U_i$ sont quasi-compactes.

\medskip\noindent
Si (2) est vérifiée, soit $X_i \subset X$ le sous-espace ouvert correspondant
à l'image de l'application ouverte $|U_i| \to |X|$; voir
Propriétés des espaces,
lemme \ref{spaces-properties-lemma-etale-image-open}.
Les morphismes $U_i \to X_i$ sont surjectifs.
Ainsi, $U_i \to X_i$ est étale surjectif, et les projections
$U_i \times_{X_i} U_i \to U_i$ sont quasi-compactes, car
$U_i \times_{X_i} U_i = U_i \times_X U_i$. Par conséquent, d'après
Espaces, lemme \ref{spaces-lemma-representable-morphisms-spaces-property},
les morphismes $U_i \to X_i$ sont quasi-compacts.
\end{proof}








\section{Conditions sur les espaces algébriques}
\label{section-conditions}

\noindent
Dans cette section, nous étudions les relations entre diverses conditions
naturelles sur les espaces algébriques rencontrées plus haut. Veuillez lire
la section \ref{section-reasonable-decent}
pour vous faire une idée du sens de ces conditions.

\begin{lemma}
\label{lemma-bounded-fibres}
Soit $S$ un schéma. Soit $X$ un espace algébrique sur $S$.
Considérons les conditions suivantes sur $X$ :
\begin{itemize}
\item[] $(\alpha)$ Pour tout $x \in |X|$, les conditions équivalentes du
Lemme \ref{lemma-U-finite-above-x}
sont vérifiées.
\item[] $(\beta)$ Pour tout $x \in |X|$, les conditions équivalentes du
Lemme \ref{lemma-R-finite-above-x}
sont vérifiées.
\item[] $(\gamma)$ Pour tout $x \in |X|$, les conditions équivalentes du
Lemme \ref{lemma-UR-finite-above-x}
sont vérifiées.
\item[] $(\delta)$ Les conditions équivalentes du
Lemme \ref{lemma-U-universally-bounded}
sont vérifiées.
\item[] $(\epsilon)$ Les conditions équivalentes du
Lemme \ref{lemma-characterize-very-reasonable}
sont vérifiées.
\item[] $(\zeta)$ L'espace $X$ est localement quasi-séparé pour la topologie
de Zariski.
\item[] $(\eta)$ L'espace $X$ est quasi-séparé.
\item[] $(\theta)$ L'espace $X$ est représentable, c'est-à-dire que $X$ est un
schéma.
\item[] $(\iota)$ L'espace $X$ est un schéma quasi-séparé.
\end{itemize}
Nous avons
$$
\xymatrix{
& (\theta) \ar@{=>}[rd] & & & &  \\
(\iota) \ar@{=>}[ru] \ar@{=>}[rd] & &
(\zeta) \ar@{=>}[r] &
(\epsilon) \ar@{=>}[r] &
(\delta) \ar@{=>}[r] &
(\gamma) \ar@{<=>}[r] & (\alpha) + (\beta) \\
& (\eta) \ar@{=>}[ru] & & & &
}
$$
\end{lemma}

\begin{proof}
L'implication $(\gamma) \Leftrightarrow (\alpha) + (\beta)$ est immédiate.
Les implications du losange de gauche résultent clairement des définitions.

\medskip\noindent
Supposons $(\zeta)$, c'est-à-dire que $X$ soit localement quasi-séparé pour
la topologie de Zariski. Alors $(\epsilon)$ est vérifiée d'après
Propriétés des espaces, lemme
\ref{spaces-properties-lemma-quasi-separated-quasi-compact-pieces}.

\medskip\noindent
Supposons $(\epsilon)$. D'après le
Lemme \ref{lemma-characterize-very-reasonable},
il existe un recouvrement ouvert de Zariski $X = \bigcup X_i$ tel que, pour
chaque $i$, il existe un schéma $U_i$ et un morphisme étale surjectif
quasi-compact $U_i \to X_i$. Fixons un $i$ et un sous-schéma ouvert affine
$W \subset U_i$. Il suffit de montrer que $W \to X$ a des fibres
universellement bornées, puisque la famille de tous ces morphismes
$W \to X$ recouvre alors $X$.
Pour cela, considérons le diagramme
$$
\xymatrix{
W \times_X U_i \ar[r]_-p \ar[d]_q & U_i \ar[d] \\
W \ar[r] & X
}
$$
Puisque $W \to X$ se factorise par $X_i$, nous voyons que
$W \times_X U_i = W \times_{X_i} U_i$; par conséquent, $q$ est quasi-compact.
Comme $W$ est affine, il en résulte que le schéma $W \times_X U_i$
est quasi-compact. Nous pouvons donc appliquer
Morphismes, lemme
\ref{morphisms-lemma-locally-quasi-finite-qc-source-universally-bounded},
et nous concluons que $p$ a des fibres universellement bornées. Du
Lemme \ref{lemma-descent-universally-bounded},
nous concluons que $W \to X$ a lui aussi des fibres universellement bornées.

\medskip\noindent
Supposons $(\delta)$. Soit $U$ un schéma affine, et soit $U \to X$ un morphisme
étale. Par hypothèse, les fibres du morphisme $U \to X$ sont universellement
bornées. Il en va donc de même des fibres des deux projections
$R = U \times_X U \to U$; voir le
Lemme \ref{lemma-base-change-universally-bounded}.
Et, d'après le
Lemme \ref{lemma-composition-universally-bounded},
les fibres de $R \to X$ sont elles aussi universellement bornées.
Par conséquent, pour tout $x \in X$, les fibres de $|U| \to |X|$ et de
$|R| \to |X|$ au-dessus de $x$ sont finies; voir le
Lemme \ref{lemma-universally-bounded-finite-fibres}.
Autrement dit, les conditions équivalentes du
Lemme \ref{lemma-UR-finite-above-x}
sont vérifiées. Cela prouve $(\delta) \Rightarrow (\gamma)$.
\end{proof}

\begin{lemma}
\label{lemma-properties-local}
Soit $S$ un schéma.
Soit $\mathcal{P}$ l'une des propriétés
$(\alpha)$, $(\beta)$, $(\gamma)$, $(\delta)$, $(\epsilon)$, $(\zeta)$ ou
$(\theta)$ des espaces algébriques énumérées dans le
Lemme \ref{lemma-bounded-fibres}.
Alors, si $X$ est un espace algébrique sur $S$ et si
$X = \bigcup X_i$ est un recouvrement ouvert de Zariski tel que chaque
$X_i$ possède $\mathcal{P}$, alors $X$ possède $\mathcal{P}$.
\end{lemma}

\begin{proof}
Soit $X$ un espace algébrique sur $S$, et soit $X = \bigcup X_i$ un
recouvrement ouvert de Zariski tel que chaque $X_i$ possède $\mathcal{P}$.

\medskip\noindent
Le cas $\mathcal{P} = (\alpha)$. La condition $(\alpha)$ pour $X_i$
signifie que, pour tout $x \in |X_i|$, tout schéma affine $U$ et tout
morphisme étale $\varphi : U \to X_i$, la fibre de
$\varphi : |U| \to |X_i|$ au-dessus de $x$ est finie. Considérons
$x \in X$, un schéma affine $U$ et un morphisme étale $U \to X$.
Puisque $X = \bigcup X_i$ est un recouvrement ouvert de Zariski, il existe
un recouvrement ouvert affine fini $U = U_1 \cup \ldots \cup U_n$ tel que
chaque $U_j \to X$ se factorise par un certain $X_{i_j}$. Par hypothèse,
les fibres de $|U_j | \to |X_{i_j}|$ au-dessus de $x$ sont finies pour
$j = 1, \ldots, n$. Cela signifie clairement que la fibre de
$|U| \to |X|$ au-dessus de $x$ est finie.
Ceci prouve le résultat pour $(\alpha)$.

\medskip\noindent
Le cas $\mathcal{P} = (\beta)$. La condition $(\beta)$ pour $X_i$ signifie
que tout $x \in |X_i|$ est représenté par un monomorphisme du spectre
d'un corps vers $X_i$. Il en va donc de même pour $X$, puisque
$X_i \to X$ est un monomorphisme et que $X = \bigcup X_i$.

\medskip\noindent
Le cas $\mathcal{P} = (\gamma)$.
Notons que $(\gamma) = (\alpha) + (\beta)$ d'après le
Lemme \ref{lemma-bounded-fibres}; le lemme pour $(\gamma)$ résulte donc
des cas traités ci-dessus.

\medskip\noindent
Le cas $\mathcal{P} = (\delta)$. La condition $(\delta)$ pour $X_i$ signifie
qu'il existe des schémas $U_{ij}$ et des morphismes étales
$U_{ij} \to X_i$ à fibres universellement bornées qui recouvrent $X_i$.
Ces schémas fournissent aussi un morphisme étale surjectif
$\coprod U_{ij} \to X$, et $U_{ij} \to X$ a encore des fibres
universellement bornées.

\medskip\noindent
Le cas $\mathcal{P} = (\epsilon)$. La condition $(\epsilon)$ pour $X_i$
signifie que l'on peut trouver un ensemble $J_i$ et des morphismes
$\varphi_{ij} : U_{ij} \to X_i$ tels que chaque $\varphi_{ij}$ soit étale,
que les deux projections $U_{ij} \times_{X_i} U_{ij} \to U_{ij}$ soient
quasi-compactes et que $\coprod_{j \in J_i} U_{ij} \to X_i$ soit surjectif.
Dans ce cas, les composés $U_{ij} \to X_i \to X$ sont étales
(combiner
Morphismes, lemmes
\ref{morphisms-lemma-composition-etale} et
\ref{morphisms-lemma-open-immersion-etale}
et
Espaces, lemmes
\ref{spaces-lemma-composition-representable-transformations-property} et
\ref{spaces-lemma-morphism-schemes-gives-representable-transformation-property}
).
Comme $X_i \subset X$ est un sous-espace, nous voyons que
$U_{ij} \times_{X_i} U_{ij} = U_{ij} \times_X U_{ij}$; la condition sur
les produits fibrés est donc préservée. Et il est clair que
$\coprod_{i, j} U_{ij} \to X$ est surjectif. Ainsi, $X$
satisfait $(\epsilon)$.

\medskip\noindent
Le cas $\mathcal{P} = (\zeta)$. La condition $(\zeta)$ pour $X_i$ signifie
que $X_i$ est localement quasi-séparé pour la topologie de Zariski.
Il est immédiat que $X$ est alors localement quasi-séparé pour la
topologie de Zariski.

\medskip\noindent
Pour $(\theta)$, voir
Propriétés des espaces,
lemme \ref{spaces-properties-lemma-subscheme}.
\end{proof}

\begin{lemma}
\label{lemma-representable-properties}
Soit $S$ un schéma. Soit $\mathcal{P}$ l'une des propriétés
$(\beta)$, $(\gamma)$, $(\delta)$, $(\epsilon)$ ou
$(\theta)$ des espaces algébriques énumérées dans le
Lemme \ref{lemma-bounded-fibres}.
Soient $X$, $Y$ des espaces algébriques sur $S$.
Soit $X \to Y$ un morphisme représentable.
Si $Y$ possède la propriété $\mathcal{P}$, il en est de même de $X$.
\end{lemma}

\begin{proof}
Supposons que $f : X \to Y$ soit un morphisme représentable d'espaces
algébriques et que $Y$ possède $\mathcal{P}$. Soit $x \in |X|$, et posons
$y = f(x) \in |Y|$.

\medskip\noindent
Le cas $\mathcal{P} = (\beta)$. La condition $(\beta)$ pour $Y$ signifie
qu'il existe un monomorphisme $\Spec(k) \to Y$ représentant $y$.
Le produit fibré $X_y = \Spec(k) \times_Y X$ est un schéma, et
$x$ correspond à un point de $X_y$, c'est-à-dire à un monomorphisme
$\Spec(k') \to X_y$. Comme $X_y \to X$ est lui aussi un monomorphisme,
nous voyons que $x$ est représenté par le monomorphisme
$\Spec(k') \to X_y \to X$.
Autrement dit, $(\beta)$ est vérifiée pour $X$.

\medskip\noindent
Le cas $\mathcal{P} = (\gamma)$. Puisque $(\gamma) \Rightarrow (\beta)$,
nous avons vu au paragraphe précédent que $y$ et $x$ peuvent être
représentés par des monomorphismes comme dans le diagramme suivant
$$
\xymatrix{
\Spec(k') \ar[r]_-x \ar[d] & X \ar[d] \\
\Spec(k) \ar[r]^-y & Y
}
$$
De plus, par définition de la propriété $(\gamma)$ au moyen du
Lemme \ref{lemma-UR-finite-above-x} (2),
il existe des schémas $V_i$ et des morphismes étales $V_i \to Y$ tels que
$\coprod V_i \to Y$ soit surjectif et que, pour chaque $i$, en posant
$R_i = V_i \times_Y V_i$, les fibres des deux applications
$$
|V_i| \longrightarrow |Y|
\quad\text{et}\quad
|R_i| \longrightarrow |Y|
$$
au-dessus de $y$ soient finies. Cela signifie que les schémas
$(V_i)_y$ et $(R_i)_y$ sont des schémas finis sur $y = \Spec(k)$.
Comme $X \to Y$ est représentable, les produits fibrés
$U_i = V_i \times_Y X$ sont des schémas. Les morphismes $U_i \to X$
sont étales, et $\coprod U_i \to X$ est surjectif. Enfin, pour chaque $i$,
nous avons
$$
(U_i)_x =
(V_i \times_Y X)_x =
(V_i)_y \times_{\Spec(k)} \Spec(k')
$$
et
$$
(U_i \times_X U_i)_x =
\left((V_i \times_Y X) \times_X (V_i \times_Y X)\right)_x =
(R_i)_y \times_{\Spec(k)} \Spec(k')
$$
de sorte que ces schémas sont finis sur $k'$, comme changements de base des
schémas finis $(V_i)_y$ et $(R_i)_y$. Il s'ensuit que $(\gamma)$ est
vérifiée pour $X$, de nouveau par la seconde condition du
Lemme \ref{lemma-UR-finite-above-x}.

\medskip\noindent
Le cas $\mathcal{P} = (\delta)$. Soit $V \to Y$ un morphisme étale, où
$V$ est un schéma affine. Puisque $Y$ possède la propriété $(\delta)$,
ce morphisme a des fibres universellement bornées. D'après le
Lemme \ref{lemma-base-change-universally-bounded},
le changement de base $V \times_Y X \to X$ a lui aussi des fibres
universellement bornées. La première partie du
Lemme \ref{lemma-U-universally-bounded}
s'applique donc, et nous voyons que $X$ possède lui aussi la propriété
$(\delta)$.

\medskip\noindent
Le cas $\mathcal{P} = (\epsilon)$. Nous utiliserons à plusieurs reprises
Espaces, lemme
\ref{spaces-lemma-base-change-representable-transformations-property}.
Soient $V_i \to Y$ comme dans le
Lemme \ref{lemma-characterize-very-reasonable} (2).
Posons $U_i = X \times_Y V_i$. Les morphismes $U_i \to X$ sont étales,
et $\coprod U_i \to X$ est surjectif. Puisque
$U_i \times_X U_i = X \times_Y (V_i \times_Y V_i)$, nous voyons que les
projections $U_i \times_X U_i \to U_i$ sont les changements de base des
projections $V_i \times_Y V_i \to V_i$, et sont donc elles aussi
quasi-compactes. Ainsi, $X$ satisfait la condition (2) du
Lemme \ref{lemma-characterize-very-reasonable}.

\medskip\noindent
Le cas $\mathcal{P} = (\theta)$. Dans ce cas, le résultat est
Catégories, lemme \ref{categories-lemma-representable-over-representable}.
\end{proof}












\section{Espaces algébriques raisonnables et décents}
\label{section-reasonable-decent}

\noindent
Dans le
Lemme \ref{lemma-bounded-fibres},
nous avons rencontré plusieurs conditions sur les espaces algébriques liées
au comportement des morphismes étales de schémas affines vers $X$
et à l'existence de recouvrements étales particuliers de $X$ par des
schémas. Nous récapitulons ici les différents types de conditions :
$$
\boxed{
\begin{matrix}
(\alpha) & \text{les fibres des morphismes étales issus d'affines
sont finies} \\
(\beta) & \text{les points proviennent de monomorphismes
de spectres de corps} \\
(\gamma) & \text{les points proviennent de monomorphismes quasi-compacts de
spectres de corps} \\
(\delta) & \text{les fibres des morphismes étales issus d'affines sont
universellement bornées} \\
(\epsilon) & \text{recouvrement par des morphismes étales issus de schémas,
quasi-compacts sur leur image}
\end{matrix}
}
$$

\medskip\noindent
Les conditions de la définition suivante ne sont pas exactement des conditions
sur la diagonale de $X$, mais ce sont, en un certain sens, des conditions de
séparation sur $X$.

\begin{definition}
\label{definition-very-reasonable}
Soit $S$ un schéma.
Soit $X$ un espace algébrique sur $S$.
\begin{enumerate}
\item Nous disons que $X$ est {\it décent} si, pour tout point $x \in X$, les
conditions équivalentes du
Lemme \ref{lemma-UR-finite-above-x}
sont vérifiées, autrement dit si la propriété $(\gamma)$ du
Lemme \ref{lemma-bounded-fibres}
est vérifiée.
\item Nous disons que $X$ est {\it raisonnable} si les conditions équivalentes
du
Lemme \ref{lemma-U-universally-bounded}
sont vérifiées, autrement dit si la propriété $(\delta)$ du
Lemme \ref{lemma-bounded-fibres}
est vérifiée.
\item Nous disons que $X$ est {\it très raisonnable} si les conditions
équivalentes du
Lemme \ref{lemma-characterize-very-reasonable}
sont vérifiées, c'est-à-dire si la propriété $(\epsilon)$ du
Lemme \ref{lemma-bounded-fibres}
est vérifiée.
\end{enumerate}
\end{definition}

\noindent
Nous avons les implications suivantes entre ces conditions sur les espaces
algébriques :
$$
\xymatrix{
\text{représentable} \ar@{=>}[rd] & & & \\
 & \text{très raisonnable} \ar@{=>}[r] &
\text{raisonnable} \ar@{=>}[r] &
\text{décent} \\
\text{quasi-séparé} \ar@{=>}[ru] & & &
}
$$
La notion d'espace algébrique très raisonnable est obsolète.
Elle a été introduite parce que cette hypothèse était nécessaire pour
démontrer certains résultats qui le sont maintenant pour la classe des
espaces décents.
La classe des espaces décents est la plus grande classe d'espaces $X$ pour
lesquels on dispose d'une bonne relation entre la topologie de $|X|$ et les
propriétés de $X$ lui-même.

\begin{example}
\label{example-not-decent}
L'espace algébrique $\mathbf{A}^1_{\mathbf{Q}}/\mathbf{Z}$ construit dans
Espaces, exemple \ref{spaces-example-affine-line-translation}
n'est pas décent, car son ``point générique'' ne peut pas être représenté par
un monomorphisme du spectre d'un corps.
\end{example}

\begin{remark}
\label{remark-reasonable}
Les espaces algébriques raisonnables sont techniquement plus faciles à manier
que les espaces algébriques très raisonnables. Par exemple, si $X \to Y$ est
un morphisme quasi-compact, étale et surjectif d'espaces algébriques et si
$X$ est raisonnable, alors $Y$ l'est aussi; voir le
Lemme \ref{lemma-descent-conditions}.
Mais nous ne savons pas si cela reste vrai pour la propriété
``très raisonnable''. Nous donnons plus bas une autre propriété technique
que possèdent les espaces algébriques raisonnables.
\end{remark}

\begin{lemma}
\label{lemma-fun-property-reasonable}
Soit $S$ un schéma.
Soit $X$ un espace algébrique raisonnable quasi-compact sur $S$.
Il existe alors un système inductif filtrant d'espaces algébriques
quasi-compacts et quasi-séparés $X_i$ tel que $X = \colim_i X_i$
(la colimite étant prise dans la catégorie des faisceaux). De plus, on peut
faire en sorte que
\begin{enumerate}
\item pour tout schéma quasi-compact $T$ sur $S$, on ait
$\colim X_i(T) = X(T)$,
\item les morphismes de transition $X_i \to X_{i'}$ du système et les
coprojections $X_i \to X$ soient surjectifs et étales, et
\item si $X$ est un schéma, les espaces algébriques $X_i$ soient des schémas,
et les morphismes de transition $X_i \to X_{i'}$ et les coprojections
$X_i \to X$ soient des isomorphismes locaux.
\end{enumerate}
\end{lemma}

\begin{proof}
Esquissons la démonstration. D'après
Propriétés des espaces, lemme
\ref{spaces-properties-lemma-quasi-compact-affine-cover},
nous avons $X = U/R$, avec $U$ affine.
Dans ce cas, être raisonnable signifie que $U \to X$ est
universellement borné. Il existe donc un entier $N$ tel que les ``fibres'' de
$U \to X$ soient de degré au plus $N$; voir la
Définition \ref{definition-universally-bounded}.
Désignons par $s, t : R \to U$ et $c : R \times_{s, U, t} R \to R$ les
applications structurales du groupoïde.

\medskip\noindent
Affirmation : pour tout ouvert quasi-compact $A \subset R$, il existe un ouvert
$R' \subset R$ tel que
\begin{enumerate}
\item $A \subset R'$,
\item $R'$ soit quasi-compact, et
\item $(U, R', s|_{R'}, t|_{R'}, c|_{R' \times_{s, U, t} R'})$ soit
un groupoïde en schémas.
\end{enumerate}
Notons que le morphisme $e : U \to R$ est ouvert, puisqu'il est une section du
morphisme étale $s : R \to U$; voir
\'Morphismes étales, Proposition \ref{etale-proposition-properties-sections}.
De plus, $U$ est affine, donc quasi-compact. Nous pouvons ainsi remplacer $A$
par $A \cup e(U) \subset R$ et supposer que $A$ contient $e(U)$. Définissons
ensuite par récurrence $A^1 = A$ et
$$
A^n = c(A^{n - 1} \times_{s, U, t} A) \subset R
$$
pour $n \geq 2$. Un raisonnement par récurrence montre que $A^n$ est
quasi-compact pour tout $n \geq 2$, comme image du produit fibré quasi-compact
$A^{n - 1} \times_{s, U, t} A$. Si $k$ est un corps algébriquement clos sur
$S$ et si nous considérons les points à valeurs dans $k$, alors
$$
A^n(k) = \left\{(u, u') \in U(k) \times U(k)
:
\begin{matrix}
\text{il existe } u = u_1, u_2, \ldots, u_{n + 1} = u' \in U(k)
\text{ tels que} \\
(u_i , u_{i + 1}) \in A \text{ pour tout }i = 1, \ldots, n.
\end{matrix}
\right\}
$$
Mais, comme les fibres de $U(k) \to X(k)$ sont de cardinal au plus $N$, dès
que $n > N$, un terme se répète dans la suite ci-dessus et nous pouvons la
raccourcir. Cela prouve que $A^N = A^n$ pour tout $n \geq N$.
Il en résulte que $R' = A^N$ fournit le groupoïde en schémas
$(U, R', s|_{R'}, t|_{R'}, c|_{R' \times_{s, U, t} R'})$, ce qui démontre
l'affirmation.

\medskip\noindent
Considérons le morphisme de faisceaux sur $(\Sch/S)_{fppf}$
$$
\colim_{R' \subset R} U/R' \longrightarrow U/R
$$
où $R' \subset R$ parcourt les sous-schémas ouverts quasi-compacts de $R$
qui définissent des relations d'équivalence étales comme ci-dessus.
Chacun des quotients $U/R'$ est un espace algébrique
(voir Espaces, théorème \ref{spaces-theorem-presentation}).
Comme $R'$ est quasi-compact et $U$ affine, le morphisme
$R' \to U \times_{\Spec(\mathbf{Z})} U$ est quasi-compact;
par conséquent, $U/R'$ est quasi-séparé. Enfin, si $T$ est un schéma
quasi-compact, alors
$$
\colim_{R' \subset R} U(T)/R'(T) \longrightarrow U(T)/R(T)
$$
est une bijection, car tout morphisme de $T$ dans $R$ se factorise par l'une
des sous-relations ouvertes $R'$ grâce à l'affirmation ci-dessus. Cela
implique clairement que la colimite des faisceaux $U/R'$ est $U/R$. Autrement
dit, l'espace algébrique $X = U/R$ est la colimite des espaces algébriques
quasi-séparés $U/R'$.

\medskip\noindent
Les propriétés (1) et (2) résultent de la discussion précédente.
Si $X$ est un schéma, choisir pour $U$ une réunion disjointe finie d'ouverts
affines de $X$ donne (3).
Nous omettons les détails.
\end{proof}

\begin{lemma}
\label{lemma-representable-named-properties}
Soit $S$ un schéma. Soient $X$, $Y$ des espaces algébriques sur $S$.
Soit $X \to Y$ un morphisme représentable.
Si $Y$ est décent (resp.\ raisonnable), alors $X$ l'est aussi.
\end{lemma}

\begin{proof}
C'est une reformulation du Lemme \ref{lemma-representable-properties}.
\end{proof}

\begin{lemma}
\label{lemma-etale-named-properties}
Soit $S$ un schéma. Soit $X \to Y$ un morphisme étale d'espaces
algébriques sur $S$. Si $Y$ est décent, resp.\ raisonnable,
alors $X$ l'est aussi.
\end{lemma}

\begin{proof}
Soit $U$ un schéma affine et $U \to X$ un morphisme étale.
Posons $R = U \times_X U$ et $R' = U \times_Y U$. Notons que
$R \to R'$ est un monomorphisme.

\medskip\noindent
Soit $x \in |X|$. Pour montrer que $X$ est décent, nous devons montrer que
les fibres de $|U| \to |X|$ et de $|R| \to |X|$ au-dessus de $x$ sont finies.
Mais si $Y$ est décent, les fibres de $|U| \to |Y|$ et de
$|R'| \to |Y|$ sont finies. D'où le résultat pour ``décent''.

\medskip\noindent
Pour montrer que $X$ est raisonnable, nous devons montrer que les fibres de
$U \to X$ sont universellement bornées. Or, si $Y$ est raisonnable, les fibres
de $U \to Y$ sont universellement bornées, ce qui implique immédiatement la
même propriété pour les fibres de $U \to X$.
D'où le résultat pour ``raisonnable''.
\end{proof}







\section{Points et spécialisations}
\label{section-specializations}

\noindent
Il existe un morphisme étale d'espaces algébriques $f : X \to Y$ et une
spécialisation non triviale entre des points d'une fibre de
$|f| : |X| \to |Y|$; voir
Exemples, lemme \ref{examples-lemma-specializations-fibre-etale}.
Si la source du morphisme est un schéma, on peut exclure ce phénomène en
imposant la condition ($\alpha$) à $Y$.

\begin{lemma}
\label{lemma-no-specializations-map-to-same-point}
Soit $S$ un schéma.
Soit $X$ un espace algébrique sur $S$.
Soit $U \to X$ un morphisme étale d'un schéma vers $X$.
Supposons que $u, u' \in |U|$ aient pour image le même point $x$ de $|X|$ et
que $u' \leadsto u$. Si le couple $(X, x)$ satisfait les conditions
équivalentes du
Lemme \ref{lemma-U-finite-above-x},
alors $u = u'$.
\end{lemma}

\begin{proof}
Supposons que le couple $(X, x)$ satisfasse les conditions équivalentes du
Lemme \ref{lemma-U-finite-above-x}.
Soit $U$ un schéma, soit $U \to X$ un morphisme étale, et supposons que
$u, u' \in |U|$ aient pour image $x$ dans $|X|$ et que
$u' \leadsto u$. Nous pouvons remplacer, et remplaçons, $U$ par un voisinage
affine de $u$. Soient $t, s : R = U \times_X U \to U$
les projections étales.

\medskip\noindent
Choisissons un point $r \in R$ tel que $t(r) = u$ et $s(r) = u'$.
C'est possible d'après
Propriétés des espaces,
lemme \ref{spaces-properties-lemma-points-presentation}.
Comme les générisations se relèvent le long du morphisme étale $t$
(Remarque \ref{remark-recall}), nous pouvons trouver une spécialisation
$r' \leadsto r$ telle que $t(r') = u'$. Posons $u'' = s(r')$. Alors
$u'' \leadsto u'$.
Nous pouvons donc recommencer et trouver $r'' \leadsto r'$ tel que
$t(r'') = u''$. Posons $u''' = s(r'')$, et ainsi de suite.
Voici le diagramme :
$$
\xymatrix{
& r'' \ar[rd]^s \ar[ld]_t \ar@{~>}[d] & \\
u'' \ar@{~>}[d] & r' \ar[rd]^s \ar[ld]_t \ar@{~>}[d] & u''' \ar@{~>}[d] \\
u' \ar@{~>}[d] & r \ar[rd]^s \ar[ld]_t & u'' \ar@{~>}[d] \\
u & & u'
}
$$
Dans la Remarque \ref{remark-recall}, nous avons vu qu'il n'y a pas de
spécialisations entre les points des fibres du morphisme étale $s$. Ainsi,
si $u^{(n + 1)} = u^{(n)}$ pour un certain $n$, alors aussi
$r^{(n)} = r^{(n - 1)}$, puis, en appliquant $t$,
$u^{(n)} = u^{(n - 1)}$. Toute la tour est alors constante; en particulier,
$u = u'$. Nous voyons donc que, si $u \not = u'$, toutes les spécialisations
sont strictes et que $\{u, u', u'', \ldots\}$ est un ensemble infini de points
de $U$ dont l'image est le point $x$ de $|X|$. Comme nous avons choisi $U$
affine, cela contredit la seconde partie du
Lemme \ref{lemma-U-finite-above-x}, comme voulu.
\end{proof}

\begin{lemma}
\label{lemma-no-specializations-map-to-same-point-Noetherian}
Soit $S$ un schéma. Soit $X$ un espace algébrique sur $S$.
Soit $U \to X$ un morphisme étale d'un schéma vers $X$.
Supposons que $u, u' \in |U|$ aient pour image le même point $x$ de $|X|$ et
que $u' \leadsto u$. Si $X$ est localement noethérien, alors $u = u'$.
\end{lemma}

\begin{proof}
La discussion de Schémas, section \ref{schemes-section-points}
montre que $\mathcal{O}_{U, u'}$ est un localisé de
l'anneau local noethérien $\mathcal{O}_{U, u}$.
D'après Propriétés des espaces, lemme
\ref{spaces-properties-lemma-pre-dimension-local-ring},
nous avons $\dim(\mathcal{O}_{U, u}) = \dim(\mathcal{O}_{U, u'})$.
La théorie de la dimension des anneaux locaux noethériens permet de conclure
que $u = u'$.
\end{proof}

\begin{lemma}
\label{lemma-specialization}
Soit $S$ un schéma.
Soit $X$ un espace algébrique sur $S$.
Soient $x, x' \in |X|$ et supposons que $x' \leadsto x$, c'est-à-dire que $x$
est une spécialisation de $x'$.
Supposons que le couple $(X, x')$ satisfasse les conditions équivalentes
du Lemme \ref{lemma-UR-finite-above-x}.
Alors, pour tout morphisme étale $\varphi : U \to X$ d'un schéma $U$ et tout
$u \in U$ tel que $\varphi(u) = x$, il existe un point $u'\in U$ tel que
$u' \leadsto u$ et $\varphi(u') = x'$.
\end{lemma}

\begin{proof}
Nous pouvons remplacer $U$ par un voisinage ouvert affine de $u$.
Nous pouvons donc supposer que $U$ est affine. Comme $x$ appartient à l'image
du morphisme ouvert $|U| \to |X|$, il en est de même de $x'$. Nous pouvons
ainsi remplacer $X$ par le sous-espace ouvert de Zariski correspondant à
l'image de $|U| \to |X|$; voir
Propriétés des espaces,
lemme \ref{spaces-properties-lemma-etale-image-open}.
Autrement dit, nous pouvons supposer que
$U \to X$ est surjectif et étale.
Soient $s, t : R = U \times_X U \to U$ les projections.
Comme $(X, x')$ satisfait par hypothèse les conditions équivalentes du
Lemme \ref{lemma-UR-finite-above-x},
les fibres de $|U| \to |X|$ et de $|R| \to |X|$
au-dessus de $x'$ sont finies. Écrivons $\{u'_1, \ldots, u'_n\} \subset U$ et
$\{r'_1, \ldots, r'_m\} \subset R$ pour les images inverses respectives
de $\{x'\}$.
Considérons les fermés
$$
T = \overline{\{u'_1\}} \cup \ldots \cup \overline{\{u'_n\}} \subset |U|,
\quad
T' = \overline{\{r'_1\}} \cup \ldots \cup \overline{\{r'_m\}} \subset |R|.
$$
Il est clair que $s(T') \subset T$. Comme $R$ est une relation d'équivalence,
nous avons aussi $t(T') = s(T')$, puisque l'ensemble $\{r_j'\}$ est invariant
par la symétrie de $R$, par construction. Soit $w \in T$ un point quelconque.
Alors $u'_i \leadsto w$ pour un certain $i$. Choisissons $r \in R$
tel que $s(r) = w$. Comme les générisations se relèvent le long de
$s : R \to U$ (voir la Remarque \ref{remark-recall}), nous pouvons trouver
$r' \leadsto r$ tel que $s(r') = u_i'$. Alors $r' = r'_j$ pour un certain $j$,
et nous en déduisons que $w \in s(T')$. Ainsi,
$T = s(T') = t(T')$ est une partie saturée pour la relation d'équivalence
$|R|$, fermée dans $|U|$. Cela signifie que $T$ est l'image inverse d'une
partie fermée (!) $T'' = \varphi(T)$ de $|X|$; voir
Propriétés des espaces,
lemmes \ref{spaces-properties-lemma-points-presentation} et
\ref{spaces-properties-lemma-topology-points}.
Par conséquent, $T'' = \overline{\{x'\}}$.
Ainsi, $T$ contient un point $u_1$ d'image $x$, puisque $x \in T''$.
Autrement dit, pour un certain $i$, il existe une spécialisation
$u'_i \leadsto u_1$ dont l'image est la spécialisation donnée
$x' \leadsto x$.

\medskip\noindent
Pour achever la preuve, choisissons un point $r \in R$ tel que
$s(r) = u$ et $t(r) = u_1$ (à l'aide de
Propriétés des espaces,
lemme \ref{spaces-properties-lemma-points-cartesian}).
Comme les générisations se relèvent le long de $t$ et que
$u'_i \leadsto u_1$, nous pouvons trouver une spécialisation
$r' \leadsto r$ telle que $t(r') = u'_i$.
Posons $u' = s(r')$. Alors $u' \leadsto u$ et $\varphi(u') = x'$, comme voulu.
\end{proof}

\begin{lemma}
\label{lemma-generalizations-lift-flat}
Soit $S$ un schéma. Soit $f : Y \to X$ un morphisme plat d'espaces algébriques
sur $S$. Soient $x, x' \in |X|$ et supposons que $x' \leadsto x$, c'est-à-dire
que $x$ est une spécialisation de $x'$. Supposons que le couple $(X, x')$
satisfasse les conditions équivalentes du
Lemme \ref{lemma-UR-finite-above-x} (par exemple si
$X$ est décent, $X$ est quasi-séparé ou $X$ est représentable).
Alors, pour tout $y \in |Y|$ tel que $f(y) = x$, il existe un point
$y' \in |Y|$ tel que $y' \leadsto y$ et $f(y') = x'$.
\end{lemma}

\begin{proof}
(L'assertion entre parenthèses résulte de la définition des espaces décents
et des implications entre les différentes conditions de séparation
mentionnées dans la section \ref{section-reasonable-decent}.)
Choisissons un schéma $V$ et un morphisme étale surjectif $V \to Y$.
Choisissons $v \in V$ d'image $y$. Il suffit alors de démontrer le lemme pour
$V \to X$. Nous pouvons donc supposer que $Y$ est un schéma.
Choisissons un schéma $U$ et un morphisme étale surjectif $U \to X$.
Choisissons $u \in U$ d'image $x$. D'après le
Lemme \ref{lemma-specialization}, nous pouvons choisir $u' \leadsto u$
d'image $x'$. D'après Propriétés des espaces,
lemme \ref{spaces-properties-lemma-points-cartesian},
nous pouvons choisir $z \in U \times_X Y$ d'image $y$ et $u$.
Nous nous ramenons ainsi au cas du morphisme plat de schémas
$U \times_X Y \to U$. La conclusion résulte alors de
Morphismes, lemme \ref{morphisms-lemma-generalizations-lift-flat}.
\end{proof}






\section{Stratification des espaces algébriques par des schémas}
\label{section-stratifications}

\noindent
Dans cette section, nous montrons qu'un espace algébrique quasi-compact et
quasi-séparé admet une stratification finie par des sous-espaces localement
fermés dont chacun est un schéma, et telle que le recollement des strates
s'effectue par des carrés distingués élémentaires. Nous établissons d'abord
un résultat un peu plus faible pour les espaces algébriques raisonnables.

\begin{lemma}
\label{lemma-quasi-compact-reasonable-stratification}
Soit $S$ un schéma. Soit $W \to X$ un morphisme d'un schéma $W$
vers un espace algébrique $X$, plat, localement de présentation finie,
séparé, localement quasi-fini et à fibres universellement bornées. Il existe
des sous-espaces fermés réduits
$$
\emptyset = Z_{-1} \subset Z_0 \subset Z_1 \subset Z_2 \subset
\ldots \subset Z_n = X
$$
tels que, si $X_r = Z_r \setminus Z_{r - 1}$, la stratification
$X = \coprod_{r = 0, \ldots, n} X_r$ est caractérisée par la propriété
universelle suivante : étant donné $g : T \to X$, la projection
$W \times_X T \to T$ est finie localement libre de degré $r$ si et seulement si
$g(|T|) \subset |X_r|$.
\end{lemma}

\begin{proof}
Soit $n$ un entier qui majore les degrés des fibres de $W \to X$.
Choisissons un schéma $U$ et un morphisme étale surjectif $U \to X$.
Appliquons Compléments sur les morphismes, lemme
\ref{more-morphisms-lemma-stratify-flat-fp-lqf-universally-bounded}
à $W \times_X U \to U$. Nous obtenons des parties fermées
$$
\emptyset = Y_{-1} \subset Y_0 \subset Y_1 \subset Y_2 \subset
\ldots \subset Y_n = U
$$
caractérisées par la propriété de l'énoncé pour le morphisme
$W \times_X U \to U$. Manifestement, la formation de ces parties fermées
commute au changement de base. Posant $R = U \times_X U$, et notant
$s, t : R \to U$ les projections, nous concluons que
$$
s^{-1}(Y_r) = t^{-1}(Y_r)
$$
en tant que parties fermées de $R$. Autrement dit, les parties fermées
$Y_r \subset U$ sont saturées pour la relation $R$. Cela signifie que
$|Y_r|$ est l'image inverse d'une partie fermée $Z_r \subset |X|$.
Notons encore $Z_r \subset X$ le sous-espace algébrique muni de la structure
réduite induite; voir Propriétés des espaces, définition
\ref{spaces-properties-definition-reduced-induced-space}.

\medskip\noindent
Soit $g : T \to X$ un morphisme d'espaces algébriques. Choisissons un schéma
$V$ et un morphisme étale surjectif $V \to T$. Pour démontrer la dernière
assertion du lemme, il suffit de la démontrer pour la composée
$V \to X$ (d'après notre définition des morphismes finis localement libres;
voir Morphismes d'espaces, section
\ref{spaces-morphisms-section-finite-locally-free}).
De même, le morphisme de schémas $W \times_X V \to V$ est fini
localement libre de degré $r$ si et seulement si le morphisme de schémas
$$
W \times_X (U \times_X V)
\longrightarrow
U \times_X V
$$
est fini localement libre de degré $r$ (voir
Descente, lemme \ref{descent-lemma-descending-property-finite-locally-free}).
Par construction, cela se produit si et seulement si
$|U \times_X V| \to |U|$ a son image contenue dans $|Y_r|$, ce qui équivaut
à ce que $|V| \to |X|$ ait son image contenue dans $|Z_r|$.
\end{proof}

\begin{lemma}
\label{lemma-stratify-flat-fp-lqf}
Soit $S$ un schéma. Soit $W \to X$ un morphisme d'un schéma $W$
vers un espace algébrique $X$, plat, localement de présentation finie,
séparé et localement quasi-fini. Il existe alors des sous-espaces ouverts
$$
X = X_0 \supset X_1 \supset X_2 \supset \ldots
$$
tels qu'un morphisme $\Spec(k) \to X$, où $k$ est un corps,
se factorise par $X_d$ si et seulement si
$W \times_X \Spec(k)$ est de degré $\geq d$ sur $k$.
\end{lemma}

\begin{proof}
Choisissons un schéma $U$ et un morphisme étale surjectif $U \to X$.
Appliquons Compléments sur les morphismes, lemme
\ref{more-morphisms-lemma-stratify-flat-fp-lqf}
à $W \times_X U \to U$. Nous obtenons des sous-schémas ouverts
$$
U = U_0 \supset U_1 \supset U_2 \supset \ldots
$$
caractérisés par la propriété de l'énoncé pour le morphisme
$W \times_X U \to U$. Manifestement, la formation de ces sous-schémas ouverts
commute au changement de base. Posons $R = U \times_X U$ et notons
$s, t : R \to U$ les projections. Nous concluons que
$$
s^{-1}(U_d) = t^{-1}(U_d)
$$
en tant que sous-schémas ouverts de $R$. Autrement dit, les sous-schémas
ouverts $U_d \subset U$ sont saturés pour la relation $R$. Cela signifie que
$U_d$ est l'image inverse d'un sous-espace ouvert $X_d \subset X$
(Propriétés des espaces, lemme
\ref{spaces-properties-lemma-subspaces-presentation}).
\end{proof}

\begin{lemma}
\label{lemma-filter-quasi-compact}
Soit $S$ un schéma. Soit $X$ un espace algébrique quasi-compact sur $S$.
Il existe des sous-espaces ouverts
$$
\ldots \subset U_4 \subset U_3 \subset U_2 \subset U_1 = X
$$
ayant les propriétés suivantes :
\begin{enumerate}
\item si l'on pose $T_p = U_p \setminus U_{p + 1}$ (muni de la structure
réduite induite), il existe un schéma séparé $V_p$ et un morphisme étale
surjectif $f_p : V_p \to U_p$ tels que $f_p^{-1}(T_p) \to T_p$ soit un
isomorphisme,
\item si $x \in |X|$ peut être représenté par un morphisme quasi-compact
$\Spec(k) \to X$ dont la source est le spectre d'un corps, alors
$x \in T_p$ pour un certain $p$.
\end{enumerate}
\end{lemma}

\begin{proof}
D'après Propriétés des espaces, lemme
\ref{spaces-properties-lemma-quasi-compact-affine-cover},
nous pouvons choisir un schéma affine $U$ et un morphisme étale surjectif
$U \to X$. Pour $p \geq 0$, posons
$$
W_p = U \times_X \ldots \times_X U \setminus
\text{toutes les diagonales}
$$
où le produit fibré comporte $p$ facteurs. Puisque $U$ est séparé,
le morphisme $U \to X$ est séparé et tous les produits fibrés
$U \times_X \ldots \times_X U$ sont des schémas séparés. Comme $U \to X$
est séparé, la diagonale $U \to U \times_X U$ est une immersion fermée.
Comme $U \to X$ est étale, la diagonale $U \to U \times_X U$ est une
immersion ouverte; voir Morphismes d'espaces, lemmes
\ref{spaces-morphisms-lemma-etale-unramified} et
\ref{spaces-morphisms-lemma-diagonal-unramified-morphism}.
De même, tous les morphismes diagonaux sont des immersions ouvertes et
fermées, et $W_p$ est un sous-schéma ouvert et fermé de
$U \times_X \ldots \times_X U$. De plus, le morphisme
$$
U \times_X \ldots \times_X U \longrightarrow
U \times_{\Spec(\mathbf{Z})} \ldots \times_{\Spec(\mathbf{Z})} U
$$
est localement quasi-fini et séparé (Morphismes d'espaces,
lemme \ref{spaces-morphisms-lemma-fibre-product-after-map}),
et son but est un schéma affine. Toute partie finie de
$U \times_X \ldots \times_X U$ est donc contenue dans un ouvert affine;
voir Compléments sur les morphismes, lemme
\ref{more-morphisms-lemma-separated-locally-quasi-finite-over-affine}.
Il en va donc de même pour $W_p$.
Le groupe symétrique $S_p$ opère librement sur $W_p$ au-dessus de $X$
(puisque nous avons retranché le lieu des points fixes de
$U \times_X \ldots \times_X U$). D'après ce qui précède et
Propriétés des espaces, Proposition
\ref{spaces-properties-proposition-finite-flat-equivalence-global},
le quotient $V_p = W_p/S_p$ est un schéma. Comme l'action de $S_p$ sur
$W_p$ est au-dessus de $X$, il existe un morphisme $V_p \to X$.
Puisque $W_p \to X$ est étale et que $W_p \to V_p$ est étale surjectif,
il s'ensuit que $V_p \to X$ est lui aussi étale; voir
Propriétés des espaces, lemme \ref{spaces-properties-lemma-etale-local}.
Notons que $V_p$ est un schéma séparé d'après
Propriétés des espaces, lemme
\ref{spaces-properties-lemma-quotient-separated}.

\medskip\noindent
Soit $U_p \subset X$ le sous-espace ouvert image de $V_p \to X$.
Par construction, un morphisme $\Spec(k) \to X$, où $k$ est
algébriquement clos, se factorise par $U_p$ si et seulement si
$U \times_X \Spec(k)$ possède $\geq p$ points; comme d'habitude, notons que
$U \times_X \Spec(k)$ est, en tant que schéma, une réunion disjointe
(éventuellement infinie) de copies de $\Spec(k)$; voir
remarque \ref{remark-recall}. Il s'ensuit que les $U_p$ forment une filtration
de $X$ comme dans l'énoncé du lemme. De plus, notre morphisme
$\Spec(k) \to X$ se factorise par $T_p$ si et seulement si
$U \times_X \Spec(k)$ possède exactement $p$ points. Dans ce cas, nous
voyons que $V_p \times_X \Spec(k)$ possède exactement un point.
Posons $Z_p = f_p^{-1}(T_p) \subset V_p$. C'est un sous-schéma fermé de
$V_p$. Alors $Z_p \to T_p$ est un morphisme étale d'espaces algébriques qui
induit une bijection sur les points à valeurs dans $k$ pour tout corps
algébriquement clos $k$. Précisément, cela implique que $Z_p \to T_p$ est
universellement injectif, donc une immersion ouverte d'après
Morphismes d'espaces, lemme
\ref{spaces-morphisms-lemma-etale-universally-injective-open},
et par conséquent un isomorphisme. Cela démontre (1).

\medskip\noindent
Soit $x : \Spec(k) \to X$ un morphisme quasi-compact, où $k$ est un corps.
Alors la composée $\Spec(\overline{k}) \to \Spec(k) \to X$ est elle aussi
quasi-compacte (Morphismes d'espaces, lemme
\ref{spaces-morphisms-lemma-composition-quasi-compact}).
Dans ce cas, le schéma $U \times_X \Spec(\overline{k})$ est quasi-compact.
Comme nous avons vu ci-dessus qu'il est une réunion disjointe de copies de
$\Spec(\overline{k})$, nous constatons qu'il a un nombre fini de points.
Si ce nombre est $p$, alors $x \in T_p$, comme voulu, ce qui achève la preuve.
\end{proof}

\begin{lemma}
\label{lemma-filter-reasonable}
Soit $S$ un schéma. Soit $X$ un espace algébrique quasi-compact et raisonnable
sur $S$. Il existe un entier $n$ et des sous-espaces ouverts
$$
\emptyset = U_{n + 1} \subset
U_n \subset U_{n - 1} \subset \ldots \subset U_1 = X
$$
ayant la propriété suivante : si l'on pose $T_p = U_p \setminus U_{p + 1}$
(muni de la structure réduite induite), il existe un schéma séparé $V_p$ et
un morphisme étale surjectif $f_p : V_p \to U_p$ tels que
$f_p^{-1}(T_p) \to T_p$ soit un isomorphisme.
\end{lemma}

\begin{proof}
La démonstration de ce lemme est identique à celle du
Lemme \ref{lemma-filter-quasi-compact}.
Soit $n$ un entier qui majore les degrés des fibres de $U \to X$;
un tel entier existe puisque $X$ est raisonnable, voir
définition \ref{definition-very-reasonable}.
Nous constatons alors que $U_{n + 1} = \emptyset$, ce qui achève la preuve.
\end{proof}

\begin{lemma}
\label{lemma-stratify-reasonable}
Soit $S$ un schéma. Soit $X$ un espace algébrique quasi-compact et raisonnable
sur $S$. Il existe un entier $n$ et des sous-espaces ouverts
$$
\emptyset = U_{n + 1} \subset
U_n \subset U_{n - 1} \subset \ldots \subset U_1 = X
$$
tels que chaque $T_p = U_p \setminus U_{p + 1}$ (muni de la structure réduite
induite) soit un schéma.
\end{lemma}

\begin{proof}
Conséquence immédiate du Lemme \ref{lemma-filter-reasonable}.
\end{proof}

\noindent
Le résultat suivant est presque identique à
\cite[Proposition 5.7.8]{GruRay}.

\begin{lemma}
\label{lemma-filter-quasi-compact-quasi-separated}
\begin{reference}
Ce résultat est presque identique à \cite[Proposition 5.7.8]{GruRay}.
\end{reference}
Soit $X$ un espace algébrique quasi-compact et quasi-séparé sur
$\Spec(\mathbf{Z})$. Il existe un entier $n$ et des sous-espaces ouverts
$$
\emptyset = U_{n + 1} \subset
U_n \subset U_{n - 1} \subset \ldots \subset U_1 = X
$$
ayant la propriété suivante : si l'on pose $T_p = U_p \setminus U_{p + 1}$
(muni de la structure réduite induite), il existe un schéma $V_p$
quasi-compact et séparé et un morphisme étale surjectif
$f_p : V_p \to U_p$ tels que $f_p^{-1}(T_p) \to T_p$ soit un isomorphisme.
\end{lemma}

\begin{proof}
La démonstration de ce lemme est identique à celle du
Lemme \ref{lemma-filter-quasi-compact}.
Remarquons qu'un espace quasi-séparé est raisonnable; voir
lemme \ref{lemma-bounded-fibres} et
définition \ref{definition-very-reasonable}.
Nous obtenons donc $U_{n + 1} = \emptyset$, comme dans le
Lemme \ref{lemma-filter-reasonable}.
À la fin de l'argument, ajoutons que, puisque $X$ est quasi-séparé,
les schémas $U \times_X \ldots \times_X U$ sont tous quasi-compacts.
Les schémas $W_p$ sont donc quasi-compacts. Par conséquent, les quotients
$V_p = W_p/S_p$ par le groupe symétrique $S_p$ sont des schémas
quasi-compacts.
\end{proof}

\noindent
Le lemme suivant devrait probablement se trouver ailleurs.

\begin{lemma}
\label{lemma-locally-constructible}
Soit $S$ un schéma. Soit $X$ un espace algébrique quasi-séparé sur $S$.
Soit $E \subset |X|$ une partie. Alors $E$ est étale-localement constructible
(Propriétés des espaces, définition
\ref{spaces-properties-definition-locally-constructible})
si et seulement si $E$ est une partie localement constructible de
l'espace topologique $|X|$
(Topologie, définition \ref{topology-definition-constructible}).
\end{lemma}

\begin{proof}
Supposons que $E \subset |X|$ soit une partie localement constructible de
l'espace topologique $|X|$. Soit $f : U \to X$ un morphisme étale, où $U$
est un schéma. Nous devons montrer que $f^{-1}(E)$ est localement constructible
dans $U$. La question est locale sur $U$ et $X$; nous pouvons donc supposer
que $X$ est quasi-compact, que $E \subset |X|$ est constructible et que $U$
est affine. Dans ce cas, $U \to X$ est quasi-compact, donc
$f : |U| \to |X|$ est quasi-compact. Remarquons que les ouverts rétrocompacts
de $|X|$, resp.\ $U$, ne sont autres que les ouverts quasi-compacts de
$|X|$, resp.\ $U$; voir
Topologie, lemme \ref{topology-lemma-topology-quasi-separated-scheme}.
Ainsi, $f^{-1}(E)$ est constructible d'après Topologie, lemme
\ref{topology-lemma-inverse-images-constructibles}.

\medskip\noindent
Réciproquement, supposons que $E$ soit étale-localement constructible.
Nous voulons montrer que $E$ est localement constructible dans l'espace
topologique $|X|$. La question est locale sur $X$; nous pouvons donc supposer
que $X$ est à la fois quasi-compact et quasi-séparé. Montrons que, dans ce cas,
$E$ est constructible dans $|X|$. Choisissons des sous-espaces ouverts
$$
\emptyset = U_{n + 1} \subset
U_n \subset U_{n - 1} \subset \ldots \subset U_1 = X
$$
et des morphismes étales surjectifs $f_p : V_p \to U_p$ induisant des
isomorphismes $f_p^{-1}(T_p) \to T_p = U_p \setminus U_{p + 1}$, où $V_p$
est un schéma quasi-compact et séparé, comme dans le
Lemme \ref{lemma-filter-quasi-compact-quasi-separated}.
Par définition, l'image inverse $E_p \subset V_p$ de $E$ est localement
constructible dans $V_p$. Alors $E_p$ est constructible dans $V_p$ d'après
Propriétés, lemme
\ref{properties-lemma-constructible-quasi-compact-quasi-separated}.
Ainsi, $E_p \cap |f_p^{-1}(T_p)| = E \cap |T_p|$ est constructible dans
$|T_p|$ d'après
Topologie, lemme \ref{topology-lemma-intersect-constructible-with-closed}
(remarquons que $V_p \setminus f_p^{-1}(T_p)$ est quasi-compact, car c'est
l'image inverse de l'espace quasi-compact $U_{p + 1}$ par le morphisme
quasi-compact $f_p$). Par conséquent,
$$
E = (|T_n| \cap E) \cup (|T_{n - 1}| \cap E) \cup \ldots \cup
(|T_1| \cap E)
$$
est constructible d'après
Topologie, lemme \ref{topology-lemma-collate-constructible-from-constructible}.
Nous utilisons ici le fait que $|T_p|$ est constructible dans $|X|$, ce qui
résulte clairement de ce qui précède.
\end{proof}





\section{Recouvrement entier par un schéma}
\label{section-integral-cover}

\noindent
Nous démontrons ici que, pour tout espace algébrique $X$ quasi-compact et
quasi-séparé, il existe un schéma $Y$ et un morphisme entier surjectif
$Y \to X$. Après avoir développé une théorie des limites d'espaces
algébriques, nous démontrerons qu'on peut obtenir un morphisme fini; voir
Limites d'espaces, section \ref{spaces-limits-section-finite-cover}.

\begin{lemma}
\label{lemma-extend-integral-morphism}
Soit $S$ un schéma. Soit $j : V \to Y$ une immersion ouverte quasi-compacte
d'espaces algébriques sur $S$. Soit $\pi : Z \to V$ un morphisme entier.
Il existe alors un morphisme entier $\nu : Y' \to Y$ tel que $Z$ soit
$V$-isomorphe à l'image inverse de $V$ dans $Y'$.
\end{lemma}

\begin{proof}
Puisque $j$ et $\pi$ sont tous deux quasi-compacts et séparés, il en va de
même de $j \circ \pi$. Soit $\nu : Y' \to Y$ la normalisation de $Y$ dans
$Z$; voir Morphismes d'espaces, section
\ref{spaces-morphisms-section-normalization-X-in-Y}.
Bien entendu, $\nu$ est entier; voir
Morphismes d'espaces, lemme
\ref{spaces-morphisms-lemma-characterize-normalization}.
La dernière assertion résulte formellement de
Morphismes d'espaces, lemmes
\ref{spaces-morphisms-lemma-properties-normalization} et
\ref{spaces-morphisms-lemma-normalization-in-integral}.
\end{proof}

\begin{lemma}
\label{lemma-there-is-a-scheme-integral-over}
Soit $S$ un schéma. Soit $X$ un espace algébrique quasi-compact et
quasi-séparé sur $S$.
\begin{enumerate}
\item Il existe un morphisme entier surjectif $Y \to X$, où $Y$ est un
schéma,
\item étant donné un morphisme étale surjectif $U \to X$, on peut choisir
$Y \to X$ de telle sorte que, pour tout $y \in Y$, il existe un voisinage
ouvert $V \subset Y$ tel que $V \to X$ se factorise par $U$.
\end{enumerate}
\end{lemma}

\begin{proof}
La partie (1) est le cas particulier de la partie (2) où $U = X$.
Choisissons un morphisme étale surjectif $U' \to U$, où $U'$ est un schéma.
Il est clair que nous pouvons remplacer $U$ par $U'$ et donc supposer que
$U$ est un schéma. Puisque $X$ est quasi-compact, il existe un nombre fini
d'ouverts affines $U_i \subset U$ tels que $U' = \coprod U_i \to X$ soit
surjectif. En remplaçant encore $U$ par $U'$, nous voyons que nous pouvons
supposer $U$ affine. Puisque $X$ est quasi-séparé, donc raisonnable, il existe
un entier $d$ qui majore le degré des fibres géométriques de $U \to X$
(voir Lemme \ref{lemma-bounded-fibres}).
Nous démontrerons le lemme par récurrence sur $d$ pour tout schéma $U$
quasi-compact et séparé muni d'un morphisme étale surjectif vers $X$.
Si $d = 1$, alors $U = X$ et le résultat vaut avec $Y = U$.
Supposons $d > 1$.

\medskip\noindent
Nous appliquons Morphismes d'espaces, lemme
\ref{spaces-morphisms-lemma-quasi-finite-separated-quasi-affine}
et obtenons une factorisation
$$
\xymatrix{
U \ar[rr]_j \ar[rd] & & Y \ar[ld]^\pi \\
& X
}
$$
où $\pi$ est entier et $j$ une immersion ouverte quasi-compacte. Nous
pouvons supposer, et nous le faisons, que $j(U)$ est schématiquement dense
dans $Y$. Alors $U \times_X Y$ est un schéma quasi-compact et séparé
(puisqu'il est entier sur $U$), et nous avons
$$
U \times_X Y = U \amalg W
$$
Ici, le premier terme est l'image de $U \to U \times_X Y$
(qui est fermée d'après
Morphismes d'espaces, lemme \ref{spaces-morphisms-lemma-semi-diagonal},
et ouverte parce que le morphisme en question est étale entre des espaces
algébriques étales sur $Y$), et le second terme est le complémentaire
(ouvert et fermé). L'image $V \subset Y$ de $W$ est un sous-espace ouvert
contenant $Y \setminus U$.

\medskip\noindent
Le morphisme étale $W \to Y$ a des fibres géométriques de cardinalité $< d$.
En effet, c'est clair par inspection pour les points géométriques de
$U \subset Y$. Puisque $|U| \subset |Y|$ est dense, cela vaut pour tous les
points géométriques de $Y$ d'après le
Lemme \ref{lemma-quasi-compact-reasonable-stratification}
(le degré des fibres d'un morphisme étale quasi-compact n'augmente pas par
spécialisation). Nous pouvons donc appliquer l'hypothèse de récurrence à
$W \to V$ et trouver un morphisme entier surjectif $Z \to V$, où $Z$ est un
schéma, qui, localement pour la topologie de Zariski, se factorise par $W$.
Choisissons une factorisation $Z \to Z' \to Y$ dans laquelle $Z' \to Y$ est
entier et $Z \to Z'$ une immersion ouverte
(Lemme \ref{lemma-extend-integral-morphism}).
Après avoir remplacé $Z'$ par l'adhérence schématique de $Z$ dans $Z'$,
nous pouvons supposer que $Z$ est schématiquement dense dans $Z'$.
Après ce remplacement, nous avons $Z' \times_Y V = Z$. Enfin, soit
$T \subset Y$ le complémentaire $Y \setminus V$, muni de la structure de
sous-espace fermé induite. Considérons le morphisme
$$
Z' \amalg T \longrightarrow X
$$
Par construction, c'est un morphisme entier surjectif.
Puisque $T \subset U$, il est clair que le morphisme $T \to X$ se factorise
par $U$. Par ailleurs, soit $z \in Z'$ un point. Si $z \not \in Z$, alors $z$
s'envoie sur un point de $Y \setminus V \subset U$, et nous trouvons un
voisinage ouvert de $z$ sur lequel le morphisme se factorise par $U$.
Si $z \in Z$, nous disposons d'un voisinage ouvert de $z$ dans $Z$
(qui est aussi un voisinage ouvert de $z$ dans $Z'$) sur lequel le morphisme
se factorise par $W \subset U \times_X Y$, donc par $U$.
\end{proof}

\begin{lemma}
\label{lemma-there-is-a-scheme-integral-over-refined}
Soit $S$ un schéma. Soit $X$ un espace algébrique quasi-compact et
quasi-séparé sur $S$ tel que $|X|$ ait un nombre fini de composantes
irréductibles.
\begin{enumerate}
\item Il existe un morphisme entier surjectif $f : Y \to X$, où $Y$ est un
schéma, tel que $f$ soit fini étale au-dessus d'un ouvert dense quasi-compact
$U \subset X$,
\item étant donné un morphisme étale surjectif $V \to X$, on peut choisir
$Y \to X$ de telle sorte que, pour tout $y \in Y$, il existe un voisinage
ouvert $W \subset Y$ tel que $W \to X$ se factorise par $V$.
\end{enumerate}
\end{lemma}

\begin{proof}
La démonstration est, à peu de chose près, la même que celle du
Lemme \ref{lemma-there-is-a-scheme-integral-over}, avec des précisions
techniques supplémentaires pour obtenir l'ouvert dense quasi-compact $U$
(et, malheureusement, des changements de notation pour suivre $U$).

\medskip\noindent
La partie (1) est le cas particulier de la partie (2) où $V = X$.

\medskip\noindent
Démonstration de (2). Choisissons un morphisme étale surjectif $V' \to V$,
où $V'$ est un schéma. Il est clair que nous pouvons remplacer $V$ par $V'$
et donc supposer que $V$ est un schéma. Puisque $X$ est quasi-compact, il
existe un nombre fini d'ouverts affines $V_i \subset V$ tels que
$V' = \coprod V_i \to X$ soit surjectif. En remplaçant encore $V$ par $V'$,
nous voyons que nous pouvons supposer $V$ affine. Puisque $X$ est
quasi-séparé, donc raisonnable, il existe un entier $d$ qui majore le degré
des fibres géométriques de $V \to X$
(voir Lemme \ref{lemma-bounded-fibres}).

\medskip\noindent
Par récurrence sur $d \geq 1$, nous allons établir l'hypothèse de récurrence
$(H_d)$ suivante :
\begin{itemize}
\item pour tout espace algébrique quasi-compact et quasi-séparé $X$ ayant
un nombre fini de composantes irréductibles, pour tout $m \geq 0$, pour tous
schémas quasi-compacts et séparés $V_j$, $j = 1, \ldots, m$, et pour tous
morphismes étales $\varphi_j : V_j \to X$, $j = 1, \ldots, m$, tels que $d$
majore le degré des fibres géométriques de $\varphi_j : V_j\to X$ et que
$\varphi = \coprod \varphi_j : V = \coprod V_j \to X$ soit surjectif,
l'énoncé du lemme vaut pour $\varphi : V \to X$.
\end{itemize}
Si $d = 1$, alors chaque $\varphi_j$ est une immersion ouverte. Ainsi, $X$
est un schéma et le résultat vaut avec $Y = V$.
Supposons $d > 1$, supposons $(H_{d - 1})$, et soient
$m$, $\varphi_j : V_j \to X$, $j = 1, \ldots, m$ comme dans $(H_d)$.

\medskip\noindent
Soient $\eta_1, \ldots, \eta_n \in |X|$ les points génériques des
composantes irréductibles de $|X|$. D'après
Propriétés des espaces, Proposition
\ref{spaces-properties-proposition-locally-quasi-separated-open-dense-scheme},
il existe un sous-schéma ouvert $U \subset X$ tel que
$\eta_1, \ldots, \eta_n \in U$. En remplaçant $U$ par un ouvert plus petit,
nous pouvons supposer $U$ affine et, d'après
Morphismes, lemme \ref{morphisms-lemma-generically-finite},
nous pouvons supposer que chaque $\varphi_j : V_j \to X$ est fini étale
au-dessus de $U$. Bien entendu, $U$ est quasi-compact et dense dans $X$,
et $\varphi_j^{-1}(U)$ est dense dans $V_j$. En particulier, chaque $V_j$
a un nombre fini de composantes irréductibles.

\medskip\noindent
Fixons $j \in \{1, \ldots, m\}$.
Comme dans Morphismes d'espaces, lemme
\ref{spaces-morphisms-lemma-quasi-finite-separated-quasi-affine},
notons $Y_j$ la normalisation de $X$ dans $V_j$. Nous obtenons une
factorisation
$$
\xymatrix{
V_j \ar[rr] \ar[rd]_{\varphi_j} & & Y_j \ar[ld]^{\pi_j} \\
& X
}
$$
où $\pi_j$ est entier et $V_j \to Y_j$ une immersion ouverte
quasi-compacte. Puisque $Y_j$ est la normalisation de $X$ dans $V_j$, les
Morphismes d'espaces, lemmes
\ref{spaces-morphisms-lemma-properties-normalization} et
\ref{spaces-morphisms-lemma-normalization-in-integral}
montrent que $\varphi_j^{-1}(U) \to \pi_j^{-1}(U)$ est un isomorphisme.
Ainsi, $\pi_j$ est fini étale au-dessus de $U$. Remarquons que $V_j$ est
schématiquement dense dans $Y_j$, puisque $Y_j$ est la normalisation de $X$
dans $V_j$ (cela résulte de la caractérisation de la normalisation relative
dans Morphismes d'espaces, lemme
\ref{spaces-morphisms-lemma-characterize-normalization}). Puisque $V_j$ est
quasi-compact, nous voyons que $|V_j| \subset |Y_j|$ est dense; voir
Morphismes d'espaces, section
\ref{spaces-morphisms-section-scheme-theoretic-closure}
(et en particulier Morphismes d'espaces, lemme
\ref{spaces-morphisms-lemma-quasi-compact-immersion}).
Il s'ensuit que $|Y_j|$ a un nombre fini de composantes irréductibles.
Alors $V_j \times_X Y_j$ est un schéma quasi-compact et séparé
(puisqu'il est fini sur $V_j$), et
$$
V_j \times_X Y_j = V_j \amalg W_j
$$
Ici, le premier terme est l'image de $V_j \to V_j \times_X Y_j$
(qui est fermée d'après
Morphismes d'espaces, lemme \ref{spaces-morphisms-lemma-semi-diagonal},
et ouverte parce que le morphisme en question est étale entre des espaces
algébriques étales sur $Y_j$), et le second terme est le complémentaire
(ouvert et fermé).

\medskip\noindent
Le morphisme étale $W_j \to Y_j$ a des fibres géométriques de cardinalité
$< d$. En effet, c'est clair par inspection pour les points géométriques de
$V_j \subset Y_j$. Puisque $|V_j| \subset |Y_j|$ est dense, cela vaut pour
tous les points géométriques de $Y_j$ d'après le
Lemme \ref{lemma-quasi-compact-reasonable-stratification}
(le degré des fibres d'un morphisme étale quasi-compact n'augmente pas par
spécialisation). En appliquant $(H_{d - 1})$ à
$V_j \amalg W_j \to Y_j$, nous trouvons un morphisme entier surjectif
$Y_j' \to Y_j$, avec $Y_j'$ un schéma, qui se factorise localement pour la
topologie de Zariski par $V_j \amalg W_j$ et qui est fini étale au-dessus
d'un ouvert dense quasi-compact $U_j \subset Y_j$. En remplaçant $U$ par un
ouvert plus petit, nous pouvons supposer, et nous le faisons, que
$\pi_j^{-1}(U) \subset U_j$ (nous pouvons choisir, et nous choisissons, le
même $U$ pour tout $j$; nous omettons certains détails).

\medskip\noindent
Nous affirmons que
$$
Y = \coprod\nolimits_{j = 1, \ldots, m} Y'_j \longrightarrow X
$$
est la solution de notre problème. Premièrement, ce morphisme est entier,
car, sur chaque composante, nous avons la composition $Y'_j \to Y_j \to X$ de
morphismes entiers (Morphismes d'espaces, lemme
\ref{spaces-morphisms-lemma-composition-integral}). Deuxièmement, ce morphisme
se factorise localement pour la topologie de Zariski par
$V = \coprod V_j$, car nous avons vu plus haut que chaque
$Y'_j \to Y_j$ se factorise localement pour la topologie de Zariski par
$V_j \amalg W_j = V_j \times_X Y_j$. Enfin, puisque $Y'_j \to Y_j$ et
$Y_j \to X$ sont tous deux finis étales au-dessus de $U$, leur composition
l'est aussi. Ceci achève la démonstration.
\end{proof}



















\section{Lieu schématique}
\label{section-schematic}

\noindent
Dans cette section, nous démontrons qu'un espace algébrique décent possède un
sous-espace ouvert dense qui est un schéma. Nous le démontrons d'abord pour
les espaces algébriques raisonnables.


\begin{proposition}
\label{proposition-reasonable-open-dense-scheme}
Soit $S$ un schéma. Soit $X$ un espace algébrique sur $S$.
Si $X$ est raisonnable, alors il existe un sous-espace ouvert dense de $X$
qui est un schéma.
\end{proposition}

\begin{proof}
D'après Propriétés des espaces,
lemme \ref{spaces-properties-lemma-subscheme},
la question est locale sur $X$. Nous pouvons donc supposer qu'il existe un
schéma affine $U$ et un morphisme étale surjectif $U \to X$
(Propriétés des espaces, lemme
\ref{spaces-properties-lemma-cover-by-union-affines}).
Soit $n$ un entier qui majore les degrés des fibres de $U \to X$;
un tel entier existe puisque $X$ est raisonnable, voir la
Définition \ref{definition-very-reasonable}.
Nous allons montrer par récurrence sur $n$ que, dès que
\begin{enumerate}
\item $U \to X$ est un morphisme étale surjectif dont les fibres sont de
degré $\leq n$, et
\item $U$ est isomorphe à un sous-schéma localement fermé d'un schéma affine,
\end{enumerate}
le lieu schématique est dense dans $X$.

\medskip\noindent
Soit $X_n \subset X$ le sous-espace ouvert complémentaire du sous-espace
fermé $Z_{n - 1} \subset X$ construit dans le
Lemme \ref{lemma-quasi-compact-reasonable-stratification}
à l'aide du morphisme $U \to X$.
Soit $U_n \subset U$ l'image inverse de $X_n$. Alors
$U_n \to X_n$ est fini localement libre de degré $n$.
Par conséquent, $X_n$ est un schéma d'après
Propriétés des espaces, Proposition
\ref{spaces-properties-proposition-finite-flat-equivalence-global}
(et le fait que tout ensemble fini de points de $U_n$ est contenu dans un
ouvert affine de $U_n$; voir
Propriétés, lemme \ref{properties-lemma-ample-finite-set-in-affine}).

\medskip\noindent
Soit $X' \subset X$ le sous-espace ouvert tel que $|X'|$ soit l'intérieur de
$|Z_{n - 1}|$ dans $|X|$ (voir
Topologie, définition \ref{topology-definition-nowhere-dense}).
Soit $U' \subset U$ son image inverse. Alors $U' \to X'$ est étale surjectif
et le degré de ses fibres est majoré par $n - 1$. Par récurrence, nous voyons
que le lieu schématique de $X'$ est un ouvert dense $X'' \subset X'$.
La topologie élémentaire montre alors que $X'' \cup X_n \subset X$ est
ouvert et dense, ce qui conclut la preuve.
\end{proof}

\begin{theorem}[David Rydh]
\label{theorem-decent-open-dense-scheme}
Soit $S$ un schéma. Soit $X$ un espace algébrique sur $S$.
Si $X$ est décent, alors il existe un sous-espace ouvert dense de $X$ qui est
un schéma.
\end{theorem}

\begin{proof}
Supposons que $X$ soit un espace algébrique décent pour lequel le théorème
est faux. D'après Propriétés des espaces, lemme
\ref{spaces-properties-lemma-subscheme}, il existe un plus grand sous-espace
ouvert $X' \subset X$ qui est un schéma. Puisque $X'$ n'est pas dense dans
$X$, il existe un sous-espace ouvert $X'' \subset X$ tel que
$|X''| \cap |X'| = \emptyset$. En remplaçant $X$ par $X''$, nous obtenons un
espace algébrique décent non vide $X$ qui ne contient {\it aucun}
sous-espace ouvert qui soit un schéma.

\medskip\noindent
Choisissons un schéma affine non vide $U$ et un morphisme étale $U \to X$.
Nous pouvons remplacer, et remplaçons, $X$ par le sous-espace ouvert
correspondant à l'image de $|U| \to |X|$. Considérons la suite décroissante
de sous-espaces ouverts
$$
X = X_0 \supset X_1 \supset X_2 \ldots
$$
construite dans le Lemme \ref{lemma-stratify-flat-fp-lqf}
pour le morphisme $U \to X$. Remarquons que $X_0 = X_1$, puisque $U \to X$
est surjectif. Soit $U = U_0 = U_1 \supset U_2 \ldots$ la suite induite de
sous-schémas ouverts de $U$.

\medskip\noindent
Choisissons un ouvert affine non vide $V_1 \subset U_1$ (par exemple
$V_1 = U_1$). Par récurrence, nous allons construire une suite d'ouverts
affines non vides $V_1 \supset V_2 \supset \ldots$ tels que $V_n \subset U_n$.
En effet, une fois $V_1, \ldots, V_{n - 1}$ construits, nous pouvons toujours
choisir $V_n$, sauf si $V_{n - 1} \cap U_n = \emptyset$. Mais si
$V_{n - 1} \cap U_n = \emptyset$, alors le sous-espace ouvert
$X' \subset X$ tel que
$|X'| = \Im(|V_{n - 1}| \to |X|)$ est contenu dans $|X| \setminus |X_n|$.
Par conséquent, $V_{n - 1} \to X'$ est un morphisme étale dont les fibres ont
un degré majoré par $n - 1$. Autrement dit, $X'$ est raisonnable (par
définition); ainsi, $X'$ contient un sous-schéma ouvert non vide d'après la
Proposition \ref{proposition-reasonable-open-dense-scheme}.
C'est une contradiction, qui montre que nous pouvons choisir $V_n$.

\medskip\noindent
D'après Limites, lemme \ref{limits-lemma-limit-nonempty},
la limite $V_\infty = \lim V_n$ est un schéma non vide. Choisissons un
morphisme $\Spec(k) \to V_\infty$. La composée
$\Spec(k) \to V_\infty \to U \to X$ a, par construction, son image contenue
dans tous les $X_d$. Autrement dit, la fibre $U \times_X \Spec(k)$ est de
degré infini, ce qui contredit la définition d'un espace décent. Cette
contradiction achève la démonstration du théorème.
\end{proof}

\begin{lemma}
\label{lemma-when-quotient-scheme-at-point}
Soit $S$ un schéma. Soit $X \to Y$ un morphisme surjectif, fini et
localement libre d'espaces algébriques sur $S$. Pour $y \in |Y|$, les
conditions suivantes sont équivalentes :
\begin{enumerate}
\item $y$ appartient au lieu schématique de $Y$ ;
\item il existe un ouvert affine $U \subset X$ contenant l'image inverse
de $y$.
\end{enumerate}
\end{lemma}

\begin{proof}
Si $y \in Y$ appartient au lieu schématique, il possède un voisinage ouvert
affine $V \subset Y$, et l'image inverse $U$ de $V$ dans $X$ est ouverte et
finie sur $V$, donc affine. Ainsi, (1) implique (2).

\medskip\noindent
Réciproquement, supposons donné $U \subset X$ comme en (2).
Posons $R = X \times_Y X$ et notons $s, t : R \to X$ les projections.
Considérons $Z = R \setminus (s^{-1}(U) \cap t^{-1}(U))$.
C'est une partie fermée de $R$. L'image $t(Z)$ est une partie fermée
de $X$ que l'on peut décrire informellement comme l'ensemble des
points de $X$ qui sont $R$-équivalents à un point de
$X \setminus U$. Ainsi, $U' = X \setminus t(Z)$ est stable par $R$ et
constitue un sous-espace ouvert de $X$ contenu dans $U$ et contenant
la fibre de $X \to Y$ au-dessus de $y$. Comme $X \to Y$ est ouvert
(Morphismes d'espaces, lemme \ref{spaces-morphisms-lemma-fppf-open}),
l'image de $U'$ est un sous-espace ouvert $V' \subset Y$.
Puisque $U'$ est stable par $R$ et que $R = X \times_Y X$, on voit que $U'$
est l'image inverse de $V'$ (utiliser
Propriétés des espaces, lemme \ref{spaces-properties-lemma-points-cartesian}).
En remplaçant $Y$ par $V'$ et $X$ par $U'$, on voit que l'on peut supposer
que $X$ est un schéma isomorphe à un sous-schéma ouvert d'un schéma affine.

\medskip\noindent
Supposons que $X$ soit un schéma isomorphe à un sous-schéma ouvert d'un
schéma affine. Dans ce cas, le faisceau quotient fppf $X/R$ est un schéma,
voir Propriétés des espaces, Proposition
\ref{spaces-properties-proposition-finite-flat-equivalence-global}.
Puisque $Y$ est un faisceau pour la topologie fppf, on obtient un morphisme
canonique $X/R \to Y$ par lequel se factorise $X \to Y$. Comme $X \to Y$ est
surjectif, fini et localement libre, il est surjectif comme morphisme de
faisceaux
(Espaces, lemme \ref{spaces-lemma-surjective-flat-locally-finite-presentation}).
On en conclut que $X/R \to Y$ est surjectif comme morphisme de faisceaux.
D'autre part, puisque $R = X \times_Y X$ comme faisceaux, on en conclut que
$X/R \to Y$ est injectif comme morphisme de faisceaux. Ainsi, $X/R \to Y$
est un isomorphisme, et l'on voit que $Y$ est représentable.
\end{proof}

\noindent
À ce stade, nous disposons de plusieurs méthodes différentes pour démontrer
le lemme suivant.

\begin{lemma}
\label{lemma-finite-etale-cover-dense-open-scheme}
Soit $S$ un schéma. Soit $X$ un espace algébrique sur $S$.
S'il existe un morphisme fini, étale et surjectif
$U \to X$, où $U$ est un schéma, alors il existe un sous-espace ouvert dense
de $X$ qui est un schéma.
\end{lemma}

\begin{proof}[Première démonstration]
Le morphisme $U \to X$ est fini localement libre. Il existe donc une
décomposition de $X$ en sous-espaces ouverts et fermés $X_d \subset X$ telle
que $U \times_X X_d \to X_d$ soit fini localement libre de degré $d$.
On peut donc supposer que $U \to X$ est fini localement libre de degré $d$.
Dans ce cas, soient $U_i \subset U$, $i \in I$, les ouverts affines.
Pour tout $i$, le morphisme $U_i \to X$ est étale et ses fibres sont
universellement bornées (à savoir, majorées par $d$).
Autrement dit, $X$ est raisonnable, et le résultat découle de la
Proposition \ref{proposition-reasonable-open-dense-scheme}.
\end{proof}

\begin{proof}[Deuxième démonstration]
La question est locale sur $X$
(Propriétés des espaces, lemme \ref{spaces-properties-lemma-subscheme}) ;
on peut donc supposer $X$ quasi-compact. Alors $U$ est quasi-compact.
Il existe alors un sous-schéma ouvert dense $W \subset U$ qui est séparé
(Propriétés, lemme
\ref{properties-lemma-quasi-compact-dense-open-separated}).
Posons $Z = U \setminus W$.
Soit $R = U \times_X U$, et soient $s, t : R \to U$ les projections.
Alors $t^{-1}(Z)$ est une partie rare de $R$
(Topologie, lemme \ref{topology-lemma-open-inverse-image-closed-nowhere-dense})
et donc $\Delta = s(t^{-1}(Z))$ est une partie stable par $R$, fermée et rare
dans $U$
(Morphismes, lemme \ref{morphisms-lemma-image-nowhere-dense-finite}).
Soit $u \in U \setminus \Delta$ un point générique d'une composante
irréductible. Puisque ces points sont denses dans $U \setminus \Delta$
et que $\Delta$ est rare, il suffit de montrer que l'image
$x \in X$ de $u$ appartient au lieu schématique de $X$.
Remarquons que $t(s^{-1}(\{u\})) \subset W$ est un
ensemble fini de points génériques de composantes irréductibles de $W$
(comparer avec
Propriétés des espaces, lemme
\ref{spaces-properties-lemma-codimension-0-points}).
D'après Propriétés, lemme \ref{properties-lemma-maximal-points-affine},
on peut trouver un ouvert affine $V \subset W$ tel que
$t(s^{-1}(\{u\})) \subset V$. Puisque $t(s^{-1}(\{u\}))$ est la fibre
de $|U| \to |X|$ au-dessus de $x$, on conclut par le
Lemme \ref{lemma-when-quotient-scheme-at-point}.
\end{proof}

\begin{proof}[Troisième démonstration]
(Cette démonstration est essentiellement identique à la deuxième, mais
utilise moins de renvois.)
Supposons que $X$ soit un espace algébrique, que $U$ soit un schéma et que
$U \to X$ soit un morphisme fini, étale et surjectif. Écrivons
$R = U \times_X U$ et notons, comme d'habitude,
$s, t : R \to U$ les projections. Remarquons que $s, t$ sont surjectifs,
finis et étales. Assertion : la réunion des ouverts affines stables par $R$
et contenus dans $U$ est topologiquement dense dans $U$.

\medskip\noindent
Démonstration de l'assertion. Soit $W \subset U$ un ouvert affine.
Posons $W' = t(s^{-1}(W)) \subset U$. Comme $s^{-1}(W)$ est affine
(donc quasi-compact), on voit que $W' \subset U$ est un ouvert
quasi-compact. D'après
Propriétés, lemme \ref{properties-lemma-quasi-compact-dense-open-separated},
il existe un ouvert dense $W'' \subset W'$ qui est un schéma séparé.
Posons $\Delta' = W' \setminus W''$. C'est une partie fermée rare de $W'$.
Comme $t|_{s^{-1}(W)} : s^{-1}(W) \to W'$ est ouvert (car il est étale),
on voit que l'image inverse
$(t|_{s^{-1}(W)})^{-1}(\Delta') \subset s^{-1}(W)$
est une partie fermée rare (voir
Topologie, lemme \ref{topology-lemma-open-inverse-image-closed-nowhere-dense}).
Il résulte alors de
Morphismes, lemme \ref{morphisms-lemma-image-nowhere-dense-finite}, que
$$
\Delta = s\left((t|_{s^{-1}(W)})^{-1}(\Delta')\right)
$$
est une partie fermée rare de $W$. Choisissons un point $\eta \in W$,
$\eta \not \in \Delta$, qui soit un point générique d'une composante
irréductible de $W$ (et donc de $U$). Par les choix effectués ci-dessus,
l'ensemble fini
$t(s^{-1}(\{\eta\})) = \{\eta_1, \ldots, \eta_n\}$
est contenu dans le schéma séparé $W''$.
Notons que les fibres de $s$ sont des espaces discrets finis et que les
générisations se relèvent le long du morphisme étale $t$, voir
Morphismes, lemmes \ref{morphisms-lemma-etale-flat}
et \ref{morphisms-lemma-generalizations-lift-flat}.
On voit ainsi que chaque $\eta_i$ est un point générique d'une composante
irréductible de $W''$. Par conséquent, d'après
Propriétés, lemme \ref{properties-lemma-maximal-points-affine},
on peut trouver un ouvert affine $V \subset W''$ tel que
$\{\eta_1, \ldots, \eta_n\} \subset V$.
D'après
Groupoïdes, lemme \ref{groupoids-lemma-find-invariant-affine},
cela implique que $\eta$ appartient à un sous-schéma ouvert affine stable
par $R$ de $U$. L'assertion en résulte, puisque $W$ a été choisi comme un
ouvert affine arbitraire de $U$ et que l'ensemble des points génériques
des composantes irréductibles de $W \setminus \Delta$ est dense dans $W$.

\medskip\noindent
L'assertion permet d'achever la démonstration. En effet, si $W \subset U$ est
un ouvert affine stable par $R$, alors la restriction $R_W$ de $R$ à $W$
vérifie $R_W = s^{-1}(W) = t^{-1}(W)$ (voir
Groupoïdes, définition \ref{groupoids-definition-invariant-open}
et la discussion qui suit). En particulier, les morphismes $R_W \to W$
sont eux aussi finis et étales. Il s'ensuit notamment que $R_W$ est affine.
Ainsi, $W/R_W$ est un schéma, d'après
Groupoïdes, Proposition \ref{groupoids-proposition-finite-flat-equivalence}.
D'autre part, $W/R_W$ est un sous-espace ouvert de $X$, d'après
Espaces, lemme \ref{spaces-lemma-finding-opens}.
Par conséquent, l'existence d'un ensemble dense de points contenus dans des
ouverts affines stables par $R$ de $U$ implique bien que le lieu schématique
de $X$ (voir Propriétés des espaces, lemme
\ref{spaces-properties-lemma-subscheme}) est un ouvert dense dans $X$.
\end{proof}













\section{Corps résiduels et anneaux locaux henséliens}
\label{section-residue-fields-henselian-local-rings}

\noindent
Pour un espace algébrique décent, nous pouvons définir le corps résiduel et
l'anneau local hensélien en un point. Par exemple, le lemme suivant montre que
le corps résiduel d'un point d'un espace décent est bien défini.

\begin{lemma}
\label{lemma-decent-points-monomorphism}
Soit $S$ un schéma. Soit $X$ un espace algébrique sur $S$.
Considérons l'application
$$
\{\Spec(k) \to X \text{ monomorphisme où }k\text{ est un corps}\}
\longrightarrow
|X|
$$
Cette application est toujours injective. Si $X$ est décent, elle est
bijective.
\end{lemma}

\begin{proof}
Nous avons vu dans
Propriétés des espaces,
lemme \ref{spaces-properties-lemma-points-monomorphism}
que cette application est injective en général.
D'après le Lemme \ref{lemma-bounded-fibres}, elle est surjective lorsque $X$
est décent (on peut même dire que cela fait partie de la définition d'un
espace décent).
\end{proof}

\noindent
Soit $S$ un schéma. Soit $X$ un espace algébrique sur $S$.
Si un point $x \in |X|$ peut être représenté par un monomorphisme
$\Spec(k) \to X$, alors le corps $k$ est unique à isomorphisme unique près.
Pour un espace algébrique décent, un tel monomorphisme existe pour tout point
d'après le Lemme \ref{lemma-decent-points-monomorphism} ; la définition
suivante a donc un sens.

\begin{definition}
\label{definition-residue-field}
Soit $S$ un schéma. Soit $X$ un espace algébrique décent sur $S$.
Soit $x \in |X|$. Le {\it corps résiduel de $X$ au point $x$}
est l'unique corps $\kappa(x)$ muni d'un
monomorphisme $\Spec(\kappa(x)) \to X$ représentant $x$.
\end{definition}

\noindent
Soit $S$ un schéma. Soit $f : X \to Y$ un morphisme d'espaces algébriques
décents sur $S$. Soit $x \in |X|$ un point.
Posons $y = f(x) \in |Y|$. Alors la composée $\Spec(\kappa(x)) \to Y$
appartient à la classe d'équivalence qui définit $y$ et se factorise donc par
$\Spec(\kappa(y)) \to Y$. Autrement dit, on obtient un diagramme commutatif
$$
\xymatrix{
\Spec(\kappa(x)) \ar[r]_-x \ar[d] & X \ar[d]^f \\
\Spec(\kappa(y)) \ar[r]^-y & Y
}
$$
Le morphisme vertical de gauche correspond à un homomorphisme de corps
$\kappa(y) \to \kappa(x)$. Nous l'appellerons souvent simplement
l'homomorphisme induit par $f$.

\begin{lemma}
\label{lemma-identifies-residue-fields}
Soit $S$ un schéma. Soit $f : X \to Y$ un morphisme d'espaces algébriques
décents sur $S$. Soit $x \in |X|$ un point
d'image $y = f(x) \in |Y|$.
Les conditions suivantes sont équivalentes :
\begin{enumerate}
\item $f$ induit un isomorphisme $\kappa(y) \to \kappa(x)$ ;
\item le morphisme induit $\Spec(\kappa(x)) \to Y$ est un monomorphisme.
\end{enumerate}
\end{lemma}

\begin{proof}
Cela résulte immédiatement de la discussion précédente.
\end{proof}

\noindent
Le lemme suivant montre que l'anneau local hensélien d'un point
d'un espace algébrique décent est bien défini.

\begin{lemma}
\label{lemma-decent-space-elementary-etale-neighbourhood}
Soit $S$ un schéma. Soit $X$ un espace algébrique décent sur $S$.
Pour tout point $x \in |X|$, il existe un morphisme étale
$$
(U, u) \longrightarrow (X, x)
$$
où $U$ est un schéma affine, $u$ est l'unique point de $U$ au-dessus de
$x$, et l'homomorphisme induit $\kappa(x) \to \kappa(u)$ est un isomorphisme.
\end{lemma}

\begin{proof}
On peut supposer que $X$ est quasi-compact en remplaçant $X$ par un ouvert
quasi-compact contenant $x$. Rappelons que $x$ peut être représenté par un
(mono)morphisme quasi-compact provenant du spectre d'un corps (par définition
des espaces décents). Le lemme résulte donc du
Lemme \ref{lemma-filter-quasi-compact}.
\end{proof}

\begin{definition}
\label{definition-elemenary-etale-neighbourhood}
Soit $S$ un schéma. Soit $X$ un espace algébrique sur $S$.
Soit $x \in X$ un point. Un {\it voisinage étale élémentaire}
est un morphisme étale $(U, u) \to (X, x)$ où $U$ est un schéma,
$u \in U$ est un point au-dessus de $x$, et le morphisme
$u = \Spec(\kappa(u)) \to X$ est un monomorphisme.
Un {\it morphisme de voisinages étales élémentaires}
$(U, u) \to (U', u')$ est, par définition, un morphisme $U \to U'$
au-dessus de $X$ qui envoie $u$ sur $u'$.
\end{definition}

\noindent
Si $X$ n'est pas décent, la catégorie des voisinages étales élémentaires
peut être vide.

\begin{lemma}
\label{lemma-elementary-etale-neighbourhoods}
Soit $S$ un schéma. Soit $X$ un espace algébrique décent sur $S$.
Soit $x$ un point de $X$.
La catégorie des voisinages étales élémentaires de $(X, x)$
est cofiltrante (voir
Catégories, définition \ref{categories-definition-codirected}).
\end{lemma}

\begin{proof}
La catégorie est non vide d'après le
Lemme \ref{lemma-decent-space-elementary-etale-neighbourhood}.
Supposons donnés deux voisinages étales élémentaires
$(U_i, u_i) \to (X, x)$.
Considérons alors $U = U_1 \times_X U_2$. Puisque
$\Spec(\kappa(u_i)) \to X$, $i = 1, 2$, sont tous deux des monomorphismes
dans la classe de $x$ (Lemme \ref{lemma-identifies-residue-fields}),
on voit que
$$
u = \Spec(\kappa(u_1)) \times_X \Spec(\kappa(u_2))
$$
est le spectre d'un corps $\kappa(u)$ tel que les homomorphismes induits
$\kappa(u_i) \to \kappa(u)$ sont des isomorphismes. Alors
$u \to U$ est un point de $U$, et l'on voit que $(U, u) \to (X, x)$ est un
voisinage étale élémentaire qui domine $(U_i, u_i)$.
Si $a, b : (U_1, u_1) \to (U_2, u_2)$ sont deux morphismes entre
nos voisinages étales élémentaires, considérons le schéma
$$
U = U_1 \times_{(a, b), (U_2 \times_X U_2), \Delta} U_2
$$
D'après Propriétés des espaces, lemme
\ref{spaces-properties-lemma-etale-permanence},
$U \to X$ est étale. De plus, exactement de la même manière que précédemment,
on voit que $U$ possède un point $u$
tel que $(U, u) \to (X, x)$ est un voisinage
étale élémentaire. Enfin, $U \to U_1$ égalise $a$ et $b$,
ce qui achève la démonstration.
\end{proof}

\begin{definition}
\label{definition-henselian-local-ring}
Soit $S$ un schéma. Soit $X$ un espace algébrique décent sur $S$.
Soit $x \in |X|$. L'{\it anneau local hensélien de $X$ en $x$} est
$$
\mathcal{O}_{X, x}^h = \colim \Gamma(U, \mathcal{O}_U)
$$
où la limite inductive est prise sur les voisinages étales élémentaires
$(U, u) \to (X, x)$.
\end{definition}

\noindent
Voici l'analogue du
Propriétés des espaces, lemme
\ref{spaces-properties-lemma-describe-etale-local-ring}.

\begin{lemma}
\label{lemma-describe-henselian-local-ring}
Soit $S$ un schéma. Soit $X$ un espace algébrique décent sur $S$.
Soit $x \in |X|$. Soit $(U, u) \to (X, x)$ un voisinage
étale élémentaire. Alors
$$
\mathcal{O}_{X, x}^h = \mathcal{O}_{U, u}^h
$$
Autrement dit : l'anneau local hensélien de $X$ en $x$
est égal au hensélisé $\mathcal{O}_{U, u}^h$
de l'anneau local $\mathcal{O}_{U, u}$ de $U$ en $u$.
\end{lemma}

\begin{proof}
Puisque la catégorie des voisinages étales élémentaires de $(X, x)$
est cofiltrante (Lemme \ref{lemma-elementary-etale-neighbourhoods}),
on voit que
la catégorie des voisinages étales élémentaires de $(U, u)$
est initiale dans
la catégorie des voisinages étales élémentaires de $(X, x)$.
L'égalité résulte alors du
Compléments sur les morphismes, lemme \ref{more-morphisms-lemma-describe-henselization}
et du
Catégories, lemme \ref{categories-lemma-cofinal}
(initiale devient cofinale, car la limite inductive
qui définit les anneaux locaux henséliens est prise sur la
catégorie opposée de la catégorie des voisinages
étales élémentaires).
\end{proof}

\begin{lemma}
\label{lemma-henselian-local-ring-strict}
Soit $S$ un schéma. Soit $X$ un espace algébrique décent sur $S$.
Soit $\overline{x}$ un point géométrique de $X$ au-dessus de $x \in |X|$.
L'anneau local étale $\mathcal{O}_{X, \overline{x}}$ de $X$ en $\overline{x}$
(Propriétés des espaces, définition
\ref{spaces-properties-definition-etale-local-rings})
est le hensélisé strict
de l'anneau local hensélien $\mathcal{O}_{X, x}^h$ de $X$ en $x$.
\end{lemma}

\begin{proof}
Cela résulte du Lemme \ref{lemma-describe-henselian-local-ring},
de Propriétés des espaces, lemme
\ref{spaces-properties-lemma-describe-etale-local-ring}
et du fait que $(R^h)^{sh} = R^{sh}$
pour un anneau local $(R, \mathfrak m, \kappa)$ et une
clôture séparable donnée $\kappa^{sep}$ de $\kappa$.
Cette égalité résulte du
Algèbre, lemme \ref{algebra-lemma-uniqueness-henselian}.
\end{proof}

\begin{lemma}
\label{lemma-residue-field-henselian-local-ring}
Soit $S$ un schéma. Soit $X$ un espace algébrique décent sur $S$.
Soit $x \in |X|$. Le corps résiduel de
l'anneau local hensélien de $X$ en $x$
(Définition \ref{definition-henselian-local-ring})
est le corps résiduel de $X$ en $x$
(Définition \ref{definition-residue-field}).
\end{lemma}

\begin{proof}
Choisissons un voisinage étale élémentaire $(U, u) \to (X, x)$.
Alors $\kappa(u) = \kappa(x)$ et
$\mathcal{O}_{X, x}^h = \mathcal{O}_{U, u}^h$
(Lemme \ref{lemma-describe-henselian-local-ring}).
Le corps résiduel de $\mathcal{O}_{U, u}^h$
est $\kappa(u)$ d'après Algèbre, lemme \ref{algebra-lemma-henselization}
(le résultat de ce lemme est la construction/définition
du hensélisé d'un anneau local, voir
Algèbre, définition \ref{algebra-definition-henselization}).
\end{proof}

\begin{remark}
\label{remark-functoriality-henselian-local-ring}
Soit $S$ un schéma. Soit $f : X \to Y$ un morphisme
d'espaces algébriques décents sur $S$. Soit $x \in |X|$
d'image $y \in |Y|$. Choisissons des points géométriques compatibles
$(\overline{x} \to X, \overline{y} = f(\overline{x}) \to Y)$
au-dessus de $(x, y)$.
Choisissons un voisinage étale élémentaire $(V, v) \to (Y, y)$ (ce qui est
possible d'après le
Lemme \ref{lemma-decent-space-elementary-etale-neighbourhood}).
Alors $V \times_Y X$ est un espace algébrique étale sur $X$
qui possède un unique point $x'$ dont les images sont $x$ dans $X$
et $v$ dans $V$.
(Les détails sont omis; on utilise le fait que tous les points peuvent être
représentés par des monomorphismes provenant de spectres de corps.)
Choisissons un voisinage étale élémentaire $(U, u) \to (V \times_Y X, x')$.
Nous obtenons alors le diagramme commutatif suivant
$$
\xymatrix{
\Spec(\mathcal{O}_{X, \overline{x}}) \ar[r] \ar[d] &
\Spec(\mathcal{O}_{X, x}^h) \ar[r] \ar[d] &
\Spec(\mathcal{O}_{U, u}) \ar[r] \ar[d] &
U \ar[r] \ar[d] &
X \ar[d] \\
\Spec(\mathcal{O}_{Y, \overline{y}}) \ar[r] &
\Spec(\mathcal{O}_{Y, y}^h) \ar[r] &
\Spec(\mathcal{O}_{V, v}) \ar[r] &
V \ar[r] &
Y
}
$$
Cela provient des identifications
$\mathcal{O}_{X, \overline{x}} = \mathcal{O}_{U, u}^{sh}$,
$\mathcal{O}_{X, x}^h = \mathcal{O}_{U, u}^h$,
$\mathcal{O}_{Y, \overline{y}} = \mathcal{O}_{V, v}^{sh}$,
$\mathcal{O}_{Y, y}^h = \mathcal{O}_{V, v}^h$
données par le
Lemme \ref{lemma-describe-henselian-local-ring}
et par
Propriétés des espaces, lemme
\ref{spaces-properties-lemma-describe-etale-local-ring},
ainsi que de la fonctorialité de l'hensélisation (stricte)
étudiée dans Algèbre, sections \ref{algebra-section-ind-etale} et
\ref{algebra-section-henselization}.
\end{remark}









\section{Points des espaces décents}
\label{section-points}

\noindent
Dans cette section, nous démontrons quelques propriétés des points des
espaces algébriques décents. Le lemme suivant montre que la spécialisation
des points se comporte bien dans les espaces algébriques décents.
Espaces, exemple \ref{spaces-example-infinite-product}
montre que ceci n'est {\bf pas} vrai en général.

\begin{lemma}
\label{lemma-decent-no-specializations-map-to-same-point}
Soit $S$ un schéma. Soit $X$ un espace algébrique décent sur $S$.
Soit $U \to X$ un morphisme \'etale d'un schéma vers $X$.
Si $u, u' \in |U|$ ont même image dans $|X|$ et si
$u' \leadsto u$, alors $u = u'$.
\end{lemma}

\begin{proof}
Combiner les Lemmes \ref{lemma-bounded-fibres} et
\ref{lemma-no-specializations-map-to-same-point}.
\end{proof}

\begin{lemma}
\label{lemma-decent-specialization}
Soit $S$ un schéma. Soit $X$ un espace algébrique décent sur $S$.
Soient $x, x' \in |X|$ et supposons que $x' \leadsto x$, c'est-à-dire que
$x$ est une spécialisation de $x'$. Alors, pour tout morphisme \'etale
$\varphi : U \to X$ d'un schéma $U$ vers X et tout $u \in U$ tel que
$\varphi(u) = x$, il existe un point $u'\in U$ tel que
$u' \leadsto u$ et $\varphi(u') = x'$.
\end{lemma}

\begin{proof}
Combiner les Lemmes \ref{lemma-bounded-fibres} et
\ref{lemma-specialization}.
\end{proof}

\begin{lemma}
\label{lemma-kolmogorov}
Soit $S$ un schéma. Soit $X$ un espace algébrique décent sur $S$.
Alors $|X|$ est un espace de Kolmogorov (voir
Topologie, définition \ref{topology-definition-generic-point}).
\end{lemma}

\begin{proof}
Soient $x_1, x_2 \in |X|$ tels que $x_1 \leadsto x_2$ et
$x_2 \leadsto x_1$. Nous devons montrer que $x_1 = x_2$. Choisissons un
schéma $U$ et un morphisme \'etale $U \to X$ tels que $x_1, x_2$
appartiennent tous deux à l'image de $|U| \to |X|$. D'après le
Lemme \ref{lemma-decent-specialization}, nous pouvons trouver une spécialisation
$u_1 \leadsto u_2$ dans $U$ qui s'envoie sur $x_1 \leadsto x_2$.
D'après le Lemme \ref{lemma-decent-specialization}, il existe
$u_2' \leadsto u_1$ qui s'envoie sur $x_2 \leadsto x_1$. Ainsi,
$u_2' \leadsto u_2$ est une spécialisation entre des points de $U$ qui
s'envoient sur le même point de $X$, à savoir $x_2$. Cela n'est possible que
si $u_2' = u_2$, voir le
Lemme \ref{lemma-decent-no-specializations-map-to-same-point}. Par conséquent,
$u_1 = u_2$ également, comme voulu.
\end{proof}

\begin{proposition}
\label{proposition-reasonable-sober}
Soit $S$ un schéma. Soit $X$ un espace algébrique décent sur $S$.
Alors l'espace topologique $|X|$ est sobre (voir
Topologie, définition \ref{topology-definition-generic-point}).
\end{proposition}

\begin{proof}
Nous avons vu au Lemme \ref{lemma-kolmogorov} que $|X|$ est un espace de
Kolmogorov. Il reste donc à montrer que toute partie fermée irréductible
$T \subset |X|$ possède un point générique. D'après
Propriétés des espaces,
lemme \ref{spaces-properties-lemma-reduced-closed-subspace},
il existe un sous-espace fermé $Z \subset X$ tel que $|Z| = T$.
Par définition, cela signifie que $Z \to X$ est un morphisme représentable
d'espaces algébriques. Ainsi, $Z$ est un espace algébrique décent
d'après le Lemme \ref{lemma-representable-properties}. Le
Théorème \ref{theorem-decent-open-dense-scheme}
montre qu'il existe un sous-espace ouvert dense $Z' \subset Z$ qui est un
schéma. Cela signifie que $|Z'| \subset T$ est ouvert dense.
L'espace topologique $|Z'|$ est donc irréductible, ce qui signifie que
$Z'$ est un schéma irréductible. D'après
Schémas, lemme \ref{schemes-lemma-scheme-sober},
on conclut que $|Z'|$ est l'adhérence d'un seul point $\eta \in T$,
et donc aussi que $T = \overline{\{\eta\}}$, ce qui achève la preuve.
\end{proof}

\noindent
Pour les espaces algébriques décents, la dimension se comporte comme prévu.

\begin{lemma}
\label{lemma-dimension-decent-space}
Soit $S$ un schéma. La dimension telle qu'elle est définie dans
Propriétés des espaces, section \ref{spaces-properties-section-dimension}
se comporte bien sur les espaces algébriques décents $X$ au-dessus de $S$.
\begin{enumerate}
\item Si $x \in |X|$, alors $\dim_x(|X|) = \dim_x(X)$, et
\item $\dim(|X|) = \dim(X)$.
\end{enumerate}
\end{lemma}

\begin{proof}
Démonstration de (1). Choisissons un schéma $U$ muni d'un point $u \in U$
et un morphisme \'etale $h : U \to X$ qui envoie $u$ sur $x$.
Par définition, la dimension de $X$ en $x$ est $\dim_u(|U|)$.
Nous pouvons donc choisir $U$ tel que $\dim_x(X) = \dim(|U|)$.
Soit $d$ un entier. Si $\dim(U) \geq d$, il existe alors
une suite de spécialisations non triviales
$u_d \leadsto \ldots \leadsto u_0$ dans $U$. En prenant l'image,
nous trouvons une suite correspondante
$h(u_d) \leadsto \ldots \leadsto h(u_0)$
dont chacune des spécialisations est non triviale d'après le
Lemme \ref{lemma-decent-no-specializations-map-to-same-point}.
Ainsi, l'image de $|U|$ dans $|X|$ est de dimension au moins $d$.
Réciproquement, supposons que $x_d \leadsto \ldots \leadsto x_0$ soit une
suite de spécialisations dans $|X|$ telle que $x_0$ appartienne à l'image de
$|U| \to |X|$. Nous pouvons alors la relever en une suite de spécialisations
dans $U$ d'après le Lemme \ref{lemma-decent-specialization}.

\medskip\noindent
La partie (2) est une conséquence immédiate de la partie (1),
de Topologie, lemme \ref{topology-lemma-dimension-supremum-local-dimensions},
et de Propriétés des espaces, section
\ref{spaces-properties-section-dimension}.
\end{proof}

\begin{lemma}
\label{lemma-dimension-local-ring-quasi-finite}
Soit $S$ un schéma. Soit $X \to Y$ un morphisme localement quasi-fini
d'espaces algébriques sur $S$. Soit $x \in |X|$ d'image $y \in |Y|$.
Alors la dimension de l'anneau local de $Y$ en $y$ est $\geq$ celle de
l'anneau local de $X$ en $x$.
\end{lemma}

\begin{proof}
La définition de la dimension de l'anneau local d'un point d'un espace
algébrique est donnée dans Propriétés des espaces, définition
\ref{spaces-properties-definition-dimension-local-ring}.
Choisissons un morphisme \'etale $(V, v) \to (Y, y)$ où $V$ est un schéma.
Choisissons un morphisme \'etale $U \to V \times_Y X$ et un point $u \in U$
qui s'envoie sur $x \in |X|$ et sur $v \in V$. Alors $U \to V$ est localement
quasi-fini et nous devons démontrer que
$$
\dim(\mathcal{O}_{V, v}) \geq \dim(\mathcal{O}_{U, u})
$$
C'est Algèbre, lemme \ref{algebra-lemma-dimension-inequality-quasi-finite}.
\end{proof}

\begin{lemma}
\label{lemma-dimension-quasi-finite}
Soit $S$ un schéma. Soit $X \to Y$ un morphisme localement quasi-fini
d'espaces algébriques sur $S$. Alors $\dim(X) \leq \dim(Y)$.
\end{lemma}

\begin{proof}
Cela résulte du Lemme \ref{lemma-dimension-local-ring-quasi-finite}
et de Propriétés des espaces, lemme \ref{spaces-properties-lemma-dimension}.
\end{proof}

\noindent
Le lemme suivant est très légèrement plus fort que
Propriétés des espaces,
lemme \ref{spaces-properties-lemma-point-like-spaces}.
Nous améliorerons ce lemme au Lemme \ref{lemma-when-field}.

\begin{lemma}
\label{lemma-decent-point-like-spaces}
Soit $S$ un schéma. Soit $k$ un corps. Soit $X$ un espace algébrique
sur $S$ et supposons qu'il existe un morphisme \'etale surjectif
$\Spec(k) \to X$. Si $X$ est décent, alors $X \cong \Spec(k')$,
où $k/k'$ est une extension finie séparable.
\end{lemma}

\begin{proof}
L'hypothèse implique que $|X| = \{x\}$ ne possède qu'un point. Comme
$X$ est décent, il existe un monomorphisme quasi-compact
$\Spec(k') \to X$ dont l'image est $x$. Alors la projection
$U = \Spec(k') \times_X \Spec(k) \to \Spec(k)$
est un monomorphisme, d'où $U = \Spec(k)$, voir
Schémas, lemme \ref{schemes-lemma-mono-towards-spec-field}.
Par conséquent, la projection $\Spec(k) = U \to \Spec(k')$ est \'etale,
ce qui achève la preuve.
\end{proof}






\section{Espaces réduits à point unique}
\label{section-singleton}

\noindent
Un espace {\it à point unique} est un espace algébrique $X$ tel que $|X|$
est réduit à un point. De tels espaces peuvent en fait être plus
intéressants que le seul spectre d'un corps, voir
Espaces, exemple \ref{spaces-example-Qbar}.
Nous développons un peu de formalisme afin de pouvoir en parler.

\begin{lemma}
\label{lemma-flat-cover-by-field}
Soit $S$ un schéma. Soit $Z$ un espace algébrique sur $S$.
Soit $k$ un corps et soit $\Spec(k) \to Z$ un morphisme surjectif et plat.
Alors tout morphisme $\Spec(k') \to Z$, où $k'$ est un corps, est
surjectif et plat.
\end{lemma}

\begin{proof}
Considérons le carré cartésien
$$
\xymatrix{
T \ar[d] \ar[r] & \Spec(k) \ar[d] \\
\Spec(k') \ar[r] & Z
}
$$
Notons que $T \to \Spec(k')$ est plat et surjectif, de sorte que $T$
n'est pas vide. D'autre part, $T \to \Spec(k)$ est plat puisque
$k$ est un corps. Ainsi, $T \to Z$ est plat et surjectif.
Il résulte de
Morphismes d'espaces, lemme \ref{spaces-morphisms-lemma-flat-permanence}
que $\Spec(k') \to Z$ est plat. Ce morphisme est surjectif puisque
l'hypothèse entraîne que $|Z|$ est réduit à un point.
\end{proof}

\begin{lemma}
\label{lemma-unique-point}
Soit $S$ un schéma.
Soit $Z$ un espace algébrique sur $S$. Les conditions suivantes sont
équivalentes :
\begin{enumerate}
\item $Z$ est réduit et $|Z|$ est réduit à un point,
\item il existe un morphisme plat surjectif $\Spec(k) \to Z$
où $k$ est un corps, et
\item il existe un morphisme localement de type fini, surjectif et plat
$\Spec(k) \to Z$, où $k$ est un corps.
\end{enumerate}
\end{lemma}

\begin{proof}
Supposons (1). Soit $W$ un schéma et
soit $W \to Z$ un morphisme \'etale surjectif. Alors $W$ est
un schéma réduit. Soit $\eta \in W$ un point générique d'une composante
irréductible de $W$. Comme $W$ est réduit, on a
$\mathcal{O}_{W, \eta} = \kappa(\eta)$. Il s'ensuit que le morphisme canonique
$\eta = \Spec(\kappa(\eta)) \to W$ est plat. On voit que le composé
$\eta \to Z$ est plat (voir
Morphismes d'espaces, lemme \ref{spaces-morphisms-lemma-composition-flat}).
Il est aussi surjectif puisque $|Z|$ est réduit à un point. Autrement dit,
(2) est satisfaite.

\medskip\noindent
Supposons (2). Soit $W$ un schéma et
soit $W \to Z$ un morphisme \'etale surjectif. Choisissons un corps
$k$ et un morphisme plat surjectif $\Spec(k) \to Z$.
Alors $W \times_Z \Spec(k)$ est un schéma \'etale sur $k$.
Ainsi, $W \times_Z \Spec(k)$ est une réunion disjointe de spectres de corps
(voir la Remarque \ref{remark-recall}), donc est en particulier réduit. Comme
$W \times_Z \Spec(k) \to W$
est plat et surjectif, on conclut que $W$ est réduit
(Descente, lemme \ref{descent-lemma-descend-reduced}).
Autrement dit, (1) est satisfaite.

\medskip\noindent
Il est clair que (3) implique (2). Enfin, supposons (2). Choisissons un schéma
affine non vide $W$ et un morphisme \'etale $W \to Z$. Choisissons un point
fermé $w \in W$ et posons $k = \kappa(w)$. Le composé
$$
\Spec(k) \xrightarrow{w} W \longrightarrow Z
$$
est localement de type fini d'après
Morphismes d'espaces, lemmes
\ref{spaces-morphisms-lemma-composition-finite-type} et
\ref{spaces-morphisms-lemma-etale-locally-finite-type}.
Il est également plat et surjectif d'après le
Lemme \ref{lemma-flat-cover-by-field}.
Ainsi, (3) est satisfaite.
\end{proof}

\noindent
Le lemme suivant isole une classe d'espaces algébriques à point unique
légèrement meilleure que celle du lemme précédent.

\begin{lemma}
\label{lemma-unique-point-better}
Soit $S$ un schéma. Soit $Z$ un espace algébrique sur $S$.
Les conditions suivantes sont équivalentes :
\begin{enumerate}
\item $Z$ est réduit, localement noethérien, et $|Z|$
est réduit à un point, et
\item il existe un morphisme localement de présentation finie, surjectif et
plat $\Spec(k) \to Z$, où $k$ est un corps.
\end{enumerate}
\end{lemma}

\begin{proof}
Supposons (2) satisfaite. D'après le
Lemme \ref{lemma-unique-point},
on voit que $Z$ est réduit et que $|Z|$ est réduit à un point.
Soit $W$ un schéma et soit $W \to Z$ un morphisme \'etale surjectif.
Choisissons un corps $k$ et un morphisme localement de présentation finie,
surjectif et plat $\Spec(k) \to Z$.
Alors $W \times_Z \Spec(k)$ est un schéma
\'etale sur $k$, donc une réunion disjointe de spectres de corps
(voir la Remarque \ref{remark-recall}),
donc localement noethérien. Comme $W \times_Z \Spec(k) \to W$
est plat, surjectif et localement de présentation finie, on voit
que $\{W \times_Z \Spec(k) \to W\}$ est un recouvrement fppf,
et on conclut que $W$ est localement noethérien
(Descente, lemme
\ref{descent-lemma-Noetherian-local-fppf}).
Autrement dit, (1) est satisfaite.

\medskip\noindent
Supposons (1). Choisissons un schéma affine non vide $W$ et un morphisme
\'etale $W \to Z$. Choisissons un point fermé $w \in W$ et posons
$k = \kappa(w)$. Comme $W$ est localement noethérien, le morphisme
$w : \Spec(k) \to W$ est de présentation finie, voir
Morphismes, lemme \ref{morphisms-lemma-closed-immersion-finite-presentation}.
Par conséquent, le composé
$$
\Spec(k) \xrightarrow{w} W \longrightarrow Z
$$
est localement de présentation finie d'après
Morphismes d'espaces, lemmes
\ref{spaces-morphisms-lemma-composition-finite-presentation} et
\ref{spaces-morphisms-lemma-etale-locally-finite-presentation}.
Il est aussi plat et surjectif d'après le
Lemme \ref{lemma-flat-cover-by-field}.
Ainsi, (2) est satisfaite.
\end{proof}

\begin{lemma}
\label{lemma-monomorphism-into-point}
Soit $S$ un schéma.
Soit $Z' \to Z$ un monomorphisme d'espaces algébriques sur $S$.
Supposons qu'il existe un corps $k$ et un morphisme localement de présentation
finie, surjectif et plat $\Spec(k) \to Z$. Alors ou bien $Z'$
est vide, ou bien $Z' = Z$.
\end{lemma}

\begin{proof}
Nous pouvons supposer que $Z'$ est non vide. Dans ce cas, le
produit fibré $T = Z' \times_Z \Spec(k)$
est non vide, voir
Propriétés des espaces, lemme \ref{spaces-properties-lemma-points-cartesian}.
Maintenant, $T$ est un espace algébrique et la projection
$T \to \Spec(k)$ est un monomorphisme. Ainsi, $T = \Spec(k)$, voir
Morphismes d'espaces, lemme
\ref{spaces-morphisms-lemma-monomorphism-toward-field}.
On conclut que $\Spec(k) \to Z$ se factorise par $Z'$.
Mais comme $\Spec(k) \to Z$ est surjectif, plat et localement de présentation
finie, on voit que $\Spec(k) \to Z$ est surjectif comme application de
faisceaux sur $(\Sch/S)_{fppf}$ (voir
Espaces, remarque \ref{spaces-remark-warning}),
et on conclut que $Z' = Z$.
\end{proof}

\noindent
Le lemme suivant affirme qu'à chaque point d'un espace algébrique on
peut associer un espace algébrique canonique, réduit, localement noethérien
et à point unique.

\begin{lemma}
\label{lemma-find-singleton-from-point}
Soit $S$ un schéma. Soit $X$ un espace algébrique sur $S$.
Soit $x \in |X|$. Alors il existe un unique monomorphisme
$Z \to X$ d'espaces algébriques
sur $S$ tel que $Z$ soit un espace algébrique satisfaisant aux conditions
équivalentes du
Lemme \ref{lemma-unique-point-better}
et tel que l'image de $|Z| \to |X|$ soit $\{x\}$.
\end{lemma}

\begin{proof}
Choisissons un schéma $U$ et un morphisme \'etale surjectif $U \to X$.
Posons $R = U \times_X U$ ; alors $X = U/R$ est une présentation (voir
Espaces, section \ref{spaces-section-presentations}).
Posons
$$
U' = \coprod\nolimits_{u \in U\text{ au-dessus de }x} \Spec(\kappa(u)).
$$
Le morphisme canonique $U' \to U$ est un monomorphisme. Posons
$$
R' = U' \times_X U' = R \times_{(U \times_S U)} (U' \times_S U').
$$
Comme $U' \to U$ est un monomorphisme, on voit que les projections
$s', t' : R' \to U'$ se factorisent en un monomorphisme suivi d'un
morphisme \'etale. Ainsi, puisque $U'$ est une réunion disjointe de spectres
de corps, en utilisant
la Remarque \ref{remark-recall},
et
Schémas, lemme \ref{schemes-lemma-mono-towards-spec-field},
on conclut que $R'$ est une réunion disjointe de spectres de corps et que
les morphismes $s', t' : R' \to U'$ sont \'etales. Ainsi,
$Z = U'/R'$ est un espace algébrique d'après
Espaces, théorème \ref{spaces-theorem-presentation}.
Comme $R'$ est la restriction de $R$ induite par $U' \to U$, on voit que
$Z \to X$ est un monomorphisme d'après
Groupoïdes, lemme \ref{groupoids-lemma-quotient-groupoid-restrict}.
Puisque $Z \to X$ est un monomorphisme, on voit que $|Z| \to |X|$ est
injective, voir Morphismes d'espaces, lemme
\ref{spaces-morphisms-lemma-monomorphism-injective-points}.
D'après
Propriétés des espaces, lemme \ref{spaces-properties-lemma-points-cartesian},
on voit que
$$
|U'| = |Z \times_X U'| \to |Z| \times_{|X|} |U'|
$$
est surjective, ce qui implique (par notre choix de $U'$) que
l'image de $|Z| \to |X|$ est $\{x\}$. On conclut que $|Z|$ est réduit
à un point. Enfin, par construction, $U'$ est localement noethérien et réduit,
c'est-à-dire que $Z$ satisfait aux conditions équivalentes du
Lemme \ref{lemma-unique-point-better}.

\medskip\noindent
Prouvons l'unicité de $Z \to X$. Supposons que
$Z' \to X$ soit un second monomorphisme d'espaces algébriques ayant ces
propriétés. Alors les projections
$$
Z' \longleftarrow Z' \times_X Z \longrightarrow Z
$$
sont des monomorphismes. L'espace algébrique du milieu est non vide d'après
Propriétés des espaces,
lemme \ref{spaces-properties-lemma-points-cartesian}.
Par conséquent, les deux projections sont des isomorphismes d'après le
Lemme \ref{lemma-monomorphism-into-point},
ce qui achève la preuve.
\end{proof}

\noindent
Nous introduisons la terminologie suivante, qui préfigure
les gerbes résiduelles que nous introduirons plus tard, voir
Propriétés des champs, définition
\ref{stacks-properties-definition-residual-gerbe}.

\begin{definition}
\label{definition-residual-space}
Soit $S$ un schéma.
Soit $X$ un espace algébrique sur $S$. Soit $x \in |X|$.
L'
{\it espace résiduel de $X$ en $x$}\footnote{Notation non standard.}
est le monomorphisme $Z_x \to X$ construit dans le
Lemme \ref{lemma-find-singleton-from-point}.
\end{definition}

\noindent
En particulier, nous savons que $Z_x$ est un espace algébrique
localement noethérien, réduit et à point unique,
et qu'il existe un corps et un morphisme surjectif, plat et localement
de présentation finie
$$
\Spec(k) \longrightarrow Z_x.
$$
L'espace résiduel est souvent donné par un monomorphisme
provenant du spectre d'un corps.

\begin{lemma}
\label{lemma-residual-space-monomorphism}
Soit $S$ un schéma. Soit $X$ un espace algébrique sur $S$. Soit $x \in |X|$.
L'espace résiduel $Z_x$ de $X$ en $x$ est isomorphe au spectre d'un corps
si et seulement si $x$ peut être représenté par un monomorphisme
$\Spec(k) \to X$, où $k$ est un corps. Si $X$ est décent, cela vaut pour
tout $x \in |X|$.
\end{lemma}

\begin{proof}
Puisque $Z_x \to X$ est un monomorphisme, si $Z_x = \Spec(k)$ pour un corps
$k$, alors $x$ est représenté par le monomorphisme $\Spec(k) = Z_x \to X$.
Réciproquement, si $\Spec(k) \to X$ est un monomorphisme représentant $x$,
alors $Z_x \times_X \Spec(k) \to \Spec(k)$ est un monomorphisme dont la source
est non vide d'après Propriétés des espaces, lemme
\ref{spaces-properties-lemma-points-cartesian}.
Ainsi, $Z_x \times_X \Spec(k) = \Spec(k)$ d'après Morphismes d'espaces, lemme
\ref{spaces-morphisms-lemma-monomorphism-toward-field}.
On obtient donc un monomorphisme $\Spec(k) \to Z_x$. C'est
un isomorphisme d'après le Lemme \ref{lemma-monomorphism-into-point}.
La dernière assertion découle du Lemme \ref{lemma-decent-points-monomorphism}.
\end{proof}

\noindent
L'espace résiduel est un espace algébrique régulier d'après le lemme suivant.

\begin{lemma}
\label{lemma-residual-space-regular}
Un espace algébrique $Z$ réduit, localement noethérien et à point unique
est régulier.
\end{lemma}

\begin{proof}
Soit $Z$ un espace algébrique réduit, localement noethérien et à point unique
sur un schéma $S$. Soit $W \to Z$ un morphisme \'etale surjectif, où $W$
est un schéma. Soit $k$ un corps et soit $\Spec(k) \to Z$
un morphisme surjectif, plat et localement de présentation finie (voir le
Lemme \ref{lemma-unique-point-better}).
Le schéma $T = W \times_Z \Spec(k)$ est
\'etale sur $k$, donc en particulier régulier, voir la
Remarque \ref{remark-recall}.
Comme $T \to W$ est localement de présentation finie, plat et surjectif,
il s'ensuit que $W$ est régulier, voir
Descente, lemme \ref{descent-lemma-descend-regular}.
Par définition, cela signifie que $Z$ est régulier.
\end{proof}

\begin{lemma}
\label{lemma-factor-through-residual-space-Noetherian}
Soit $S$ un schéma. Soit $f : Y \to X$ un morphisme d'espaces
algébriques sur $S$. Soit $x \in |X|$ un point. Supposons que
\begin{enumerate}
\item $|f|(|Y|)$ soit contenu dans $\{x\} \subset |X|$,
\item $Y$ soit réduit, et
\item $X$ soit localement noethérien.
\end{enumerate}
Alors $f$ se factorise par l'espace résiduel $Z_x$ de $X$ en $x$.
\end{lemma}

\begin{proof}
Remarque préliminaire : puisque $Z_x \to X$ est un monomorphisme, il suffit
de trouver un morphisme \'etale surjectif $Y' \to Y$ tel que $Y' \to X$
se factorise par $Z_x$. Remarquons ici que $Y'$ est également réduit.

\medskip\noindent
Soit $U$ un schéma affine et soit $U \to X$ un morphisme \'etale
tel que $x$ appartienne à l'image de $|U| \to |X|$. Comme $X$ est localement
noethérien, $U$ est un schéma affine noethérien. En vertu de (1),
on voit que $Y' = U \times_X Y \to Y$ est un morphisme \'etale surjectif.
Notons $E \subset |U|$ l'ensemble des points d'image $x$.
Il n'existe aucune spécialisation non triviale entre les éléments de $E$, voir
le Lemme \ref{lemma-no-specializations-map-to-same-point-Noetherian}.
Le morphisme $Y' \to U$ envoie $|Y'|$ dans $E$. Par notre construction
de $Z_x$ dans la preuve du Lemme \ref{lemma-find-singleton-from-point},
nous savons que $\coprod_{u \in E} u \to X$ se factorise par $Z_x$.
Il suffit donc de démontrer que $Y' \to U$ se factorise par
$\coprod_{u \in E} u \to U$. Après avoir remplacé $Y'$ par un recouvrement
\'etale par un schéma (ce que permet notre remarque préliminaire), cela résulte
de Morphismes, lemme \ref{morphisms-lemma-factor-through-set-of-points}.
\end{proof}

\begin{lemma}
\label{lemma-factor-through-residual-space-decent}
Soit $S$ un schéma. Soit $f : Y \to X$ un morphisme d'espaces
algébriques sur $S$. Soit $x \in |X|$ un point. Supposons que
\begin{enumerate}
\item $|f|(|Y|)$ soit contenu dans $\{x\} \subset |X|$,
\item $Y$ soit réduit, et
\item $x$ puisse être représenté par un monomorphisme quasi-compact
$x : \Spec(k) \to X$, où $k$ est un corps (par exemple si $X$ est décent).
\end{enumerate}
Alors $f$ se factorise par
l'espace résiduel $Z_x = \Spec(k)$ de $X$ en $x$.
\end{lemma}

\begin{proof}
D'après le Lemme \ref{lemma-residual-space-monomorphism}, on a
$Z_x = \Spec(k)$.

\medskip\noindent
Remarque préliminaire : puisque $\Spec(k) \to X$ est un monomorphisme,
il suffit de trouver
un morphisme \'etale surjectif $Y' \to Y$ tel que $Y' \to X$
se factorise par $Z_x$. Remarquons ici que $Y'$ est également réduit.

\medskip\noindent
Après avoir remplacé $X$ par un voisinage ouvert quasi-compact de $x$,
nous pouvons supposer $X$ quasi-compact. D'après le
Lemme \ref{lemma-filter-quasi-compact}, $x$
est un point de $T \subset U \subset X$, où $T \to U$ (resp.\ $U \to X$)
est une immersion fermée (resp.\ ouverte), et $T$ est un schéma.
D'après Propriétés des espaces, lemme
\ref{spaces-properties-lemma-factor-through-open-subspace},
$f$ se factorise par $U$ ; nous pouvons donc supposer $U = X$. Alors $f$
se factorise par $T$ puisque $Y$ est réduit, voir
Propriétés des espaces, lemme \ref{spaces-properties-lemma-map-into-reduction}.
Nous pouvons donc supposer que $X = T$ est un schéma.
D'après notre remarque préliminaire, nous pouvons aussi supposer que $Y$
est un schéma. Nous sommes ainsi ramenés à Morphismes, lemme
\ref{morphisms-lemma-factor-through-point}.
\end{proof}

\begin{example}
\label{example-does-not-factor-through-residual-gerbe}
Voici un contre-exemple aux
Lemmes \ref{lemma-factor-through-residual-space-Noetherian} et
\ref{lemma-factor-through-residual-space-decent}
lorsque $X$ n'est ni localement noethérien ni décent.
Soit $k$ un corps. Soit $G$ un groupe profini infini.
Soit $Y$ le groupe $G$ vu comme un $k$-schéma en groupes affine de
dimension zéro, c'est-à-dire
$Y = \Spec(\text{applications localement constantes } G \to k)$.
Soit $\Gamma$ le groupe $G$ vu comme un $k$-schéma en groupes discret,
opérant sur $Y$ par translations. Posons $X = Y/\Gamma$.
C'est un espace algébrique réduit à un point, muni de la projection
$q : Y \to X$.
Soit $e \in G$ l'élément neutre (n'importe quel élément conviendrait), et
regardons-le comme un $k$-point de $Y$. On obtient un $k$-point
$x : \Spec(k) \to X$ qui est un monomorphisme, puisqu'il est une section de
$X \to \Spec(k)$.
Nous affirmons que (bien que $Y$ soit affine et réduit et que
$|X| = \{x\}$) le morphisme $q$ ne se factorise par aucun morphisme
$\Spec(K) \to X$, où $K$ est un corps. Sinon, il se factoriserait par
$x$ d'après Propriétés des espaces, lemme
\ref{spaces-properties-lemma-equivalence-class-point-monomorphism}.
Or, le changement de base de $q$ par $x$ est $\Gamma \to \Spec(k)$, la
projection $\Gamma \to Y$ étant l'application d'orbite
$g \mapsto g \cdot e$. Cette dernière n'a pas de section, d'où l'assertion.
\end{example}

\begin{lemma}
\label{lemma-residual-space-closed}
Soit $S$ un schéma. Soit $X$ un espace algébrique sur $S$. Soit
$x \in |X|$, et soit $Z_x \subset X$ son espace résiduel. Supposons que $X$
soit localement noethérien. Alors $x$ est un point fermé de $|X|$ si et
seulement si le morphisme $Z_x \to X$ est une immersion fermée.
\end{lemma}

\begin{proof}
Si $Z_x \to X$ est une immersion fermée, alors $x$ est un point fermé de
$|X|$, voir Morphismes d'espaces, lemme
\ref{spaces-morphisms-lemma-immersion-when-closed}.
Réciproquement, supposons que $x$ soit un point fermé de $|X|$.
Soit $Z \subset X$ le sous-espace fermé réduit tel que $|Z| = \{x\}$
(Propriétés des espaces,
lemme \ref{spaces-properties-lemma-reduced-closed-subspace}).
Alors $Z$ est localement noethérien d'après Morphismes d'espaces, lemmes
\ref{spaces-morphisms-lemma-immersion-locally-finite-type} et
\ref{spaces-morphisms-lemma-locally-finite-type-locally-noetherian}.
Puisque $Z$ est également réduit et que $|Z| = \{x\}$, $Z = Z_x$ est
l'espace résiduel par définition.
\end{proof}











\section{Espaces décents}
\label{section-decent}

\noindent
Dans cette section, nous rassemblons quelques faits utiles sur les espaces
décents.

\begin{lemma}
\label{lemma-locally-Noetherian-decent-quasi-separated}
Tout espace algébrique décent localement noethérien est quasi-séparé.
\end{lemma}

\begin{proof}
En effet, soit $X$ un espace algébrique (sur un schéma de base quelconque,
par exemple sur $\mathbf{Z}$) qui est décent et localement noethérien.
Soient $U \to X$ et $V \to X$ des morphismes étales, où $U$ et $V$ sont
des schémas affines. Nous devons montrer que $W = U \times_X V$ est
quasi-compact (Propriétés des espaces, lemme
\ref{spaces-properties-lemma-characterize-quasi-separated}).
Puisque $X$ est localement noethérien, les schémas $U$, $V$ sont noethériens
et $W$ est localement noethérien. Puisque $X$ est décent, les fibres du
morphisme $W \to U$ sont finies. En effet, on peut représenter tout
$x \in |X|$ par un monomorphisme quasi-compact $\Spec(k) \to X$.
Alors $U_k$ et $V_k$ sont des réunions disjointes finies de spectres
d'extensions finies séparables de $k$ (Remarque \ref{remark-recall}), et
nous voyons que $W_k = U_k \times_{\Spec(k)} V_k$ est fini.
Soit $n$ le degré maximal d'une fibre de $W \to U$ en un point générique
d'une composante irréductible de $U$. Considérons la stratification
$$
U = U_0 \supset U_1 \supset U_2 \supset \ldots
$$
associée à $W \to U$ dans
Compléments sur les morphismes, lemme \ref{more-morphisms-lemma-stratify-flat-fp-lqf}.
Par le choix de $n$ ci-dessus, nous concluons que $U_{n + 1}$ est vide.
Nous voyons donc que les fibres de $W \to U$ sont universellement bornées.
Nous pouvons alors appliquer Compléments sur les morphismes, lemme
\ref{more-morphisms-lemma-stratify-flat-fp-lqf-universally-bounded}
pour trouver une stratification
$$
\emptyset = Z_{-1} \subset Z_0 \subset Z_1 \subset Z_2 \subset
\ldots \subset Z_n = U
$$
par des parties fermées telle que, si $S_r = Z_r \setminus Z_{r - 1}$,
le morphisme $W \times_U S_r \to S_r$ soit fini localement libre.
Puisque $U$ est noethérien, les schémas $S_r$ sont noethériens, donc les
schémas $W \times_U S_r$ sont noethériens, et par conséquent
$W = \coprod W \times_U S_r$ est quasi-compact, comme voulu.
\end{proof}

\begin{lemma}
\label{lemma-when-field}
Soit $S$ un schéma. Soit $X$ un espace algébrique décent sur $S$.
\begin{enumerate}
\item Si $|X|$ est réduit à un point, alors $X$ est un schéma.
\item Si $|X|$ est réduit à un point et que $X$ est réduit, alors
$X \cong \Spec(k)$ pour un corps $k$.
\end{enumerate}
\end{lemma}

\begin{proof}
Supposons que $|X|$ soit réduit à un point. Il résulte immédiatement du
Théorème \ref{theorem-decent-open-dense-scheme} que $X$ est un schéma,
mais nous pouvons aussi raisonner directement comme suit.
Choisissons un schéma affine $U$ et un morphisme étale surjectif $U \to X$.
Posons $R = U \times_X U$. Alors $U$ et $R$ ont un nombre fini de points
d'après le Lemme \ref{lemma-UR-finite-above-x} (et la définition d'un espace
décent). Tous ces points sont fermés dans $U$ et $R$ d'après le
Lemme \ref{lemma-decent-no-specializations-map-to-same-point}.
Il s'ensuit que $U$ et $R$ sont des schémas affines.
Nous pouvons restreindre $U$ à un espace réduit à un point. Alors $U$ est
le spectre d'un anneau local hensélien, voir
Algèbre, lemme \ref{algebra-lemma-local-dimension-zero-henselian}.
Les projections $R \to U$ sont étales, donc finies et étales puisque
$U$ est le spectre d'un anneau local hensélien de dimension $0$, voir
Algèbre, lemme \ref{algebra-lemma-characterize-henselian}.
Il s'ensuit que $X$ est un schéma d'après
Groupoïdes, Proposition \ref{groupoids-proposition-finite-flat-equivalence}.

\medskip\noindent
L'assertion (2) découle de (1) et du fait qu'un schéma réduit dont l'espace
sous-jacent est réduit à un point est le spectre d'un corps.
\end{proof}

\begin{remark}
\label{remark-one-point-decent-scheme}
Nous verrons dans
Limites d'espaces, lemme \ref{spaces-limits-lemma-reduction-scheme}
qu'un espace algébrique dont la réduction est un schéma est un schéma.
\end{remark}

\begin{lemma}
\label{lemma-algebraic-residue-field-extension-closed-point}
Soit $S$ un schéma. Soit $X$ un espace algébrique décent sur $S$.
Considérons un diagramme commutatif
$$
\xymatrix{
\Spec(k) \ar[rr] \ar[rd] & & X \ar[ld] \\
& S
}
$$
Supposons que le point image $s \in S$ de $\Spec(k) \to S$ soit un point
fermé et que $\kappa(s) \subset k$ soit une extension algébrique.
Alors le point image $x$ de $\Spec(k) \to X$ est un point fermé de $|X|$.
\end{lemma}

\begin{proof}
Supposons que $x \leadsto x'$ pour un $x' \in |X|$. Choisissons un
morphisme étale $U \to X$, où $U$ est un schéma, et un point $u' \in U$
qui s'envoie sur $x'$. Choisissons une spécialisation $u \leadsto u'$ dans
$U$, où $u$ s'envoie sur $x$ dans $X$, voir le
Lemme \ref{lemma-decent-specialization}.
Alors $u$ est l'image d'un point $w$ du schéma
$W = \Spec(k) \times_X U$. Puisque la projection $W \to \Spec(k)$ est étale,
nous voyons que $\kappa(w) \supset k$ est une extension finie. Ainsi
$\kappa(w) \supset \kappa(s)$ est algébrique. Ainsi
$\kappa(u) \supset \kappa(s)$ est algébrique. Donc $u$ est un point fermé de
$U$ d'après Morphismes, lemme
\ref{morphisms-lemma-algebraic-residue-field-extension-closed-point-fibre}.
Ainsi $u = u'$, d'où $x = x'$.
\end{proof}

\begin{lemma}
\label{lemma-finite-residue-field-extension-finite}
Soit $S$ un schéma. Soit $X$ un espace algébrique décent sur $S$.
Considérons un diagramme commutatif
$$
\xymatrix{
\Spec(k) \ar[rr] \ar[rd] & & X \ar[ld] \\
& S
}
$$
Supposons que le point image $s \in S$ de $\Spec(k) \to S$ soit un point
fermé et que l'extension de corps $k/\kappa(s)$ soit finie.
Alors $\Spec(k) \to X$ est un morphisme fini. Si $\kappa(s) = k$,
alors $\Spec(k) \to X$ est une immersion fermée.
\end{lemma}

\begin{proof}
D'après le Lemme \ref{lemma-algebraic-residue-field-extension-closed-point},
le point image $x \in |X|$ est fermé. Soit $Z \subset X$ le sous-espace
fermé réduit tel que $|Z| = \{x\}$ (Propriétés des espaces,
lemme \ref{spaces-properties-lemma-reduced-closed-subspace}).
Remarquons que $Z$ est un espace algébrique décent d'après le
Lemme \ref{lemma-representable-named-properties}.
D'après le Lemme \ref{lemma-when-field}, nous voyons que
$Z = \Spec(k')$ pour un corps $k'$. Bien sûr,
$k \supset k' \supset \kappa(s)$. Alors $\Spec(k) \to Z$ est un morphisme
fini de schémas, et $Z \to X$ est un morphisme fini puisqu'il s'agit d'une
immersion fermée. Par conséquent, $\Spec(k) \to X$ est fini (Morphismes
d'espaces, lemme \ref{spaces-morphisms-lemma-composition-integral}).
Si $k = \kappa(s)$, alors $\Spec(k) = Z$ et $\Spec(k) \to X$ est une
immersion fermée.
\end{proof}

\begin{lemma}
\label{lemma-decent-space-closed-point}
Soit $S$ un schéma. Supposons que $X$ soit un espace algébrique décent sur
$S$. Soit $x \in |X|$ un point fermé. Alors $x$ peut être représenté par une
immersion fermée $i : \Spec(k) \to X$ provenant du spectre d'un corps.
\end{lemma}

\begin{proof}
Nous savons que $x$ peut être représenté par un monomorphisme quasi-compact
$i : \Spec(k) \to X$, où $k$ est un corps
(Définition \ref{definition-very-reasonable}).
Soit $U \to X$ un morphisme étale, où $U$ est un schéma affine.
Comme $x$ est fermé et que $X$ est décent, la fibre $F$ de
$|U| \to |X|$ au-dessus de $x$ est formée de points fermés
(Lemme \ref{lemma-decent-no-specializations-map-to-same-point}).
Puisque $i$ est un monomorphisme, il en va de même de
$U_k = U \times_X \Spec(k) \to U$. En particulier, l'application
$|U_k| \to F$ est injective. Comme $U_k$ est quasi-compact et étale sur un
corps, nous voyons que $U_k$ est une réunion disjointe finie de spectres de
corps (Remarque \ref{remark-recall}). Écrivons
$U_k = \Spec(k_1) \amalg \ldots \amalg \Spec(k_r)$.
Puisque $\Spec(k_i) \to U$ est un monomorphisme, nous voyons que son image
$u_i$ a pour corps résiduel $\kappa(u_i) = k_i$.
Comme $u_i \in F$ est un point fermé, nous concluons que le morphisme
$\Spec(k_i) \to U$ est une immersion fermée. Puisque les $u_i$ sont deux à
deux distincts, $U_k \to U$ est une immersion fermée. Par conséquent, $i$
est une immersion fermée (Morphismes d'espaces, lemme
\ref{spaces-morphisms-lemma-closed-immersion-local}). Cela achève la preuve.
\end{proof}


\section{Espaces localement séparés}
\label{section-locally-separated}

\noindent
Il se trouve qu'un espace algébrique localement séparé est décent.

\begin{lemma}
\label{lemma-infinite-number}
Soit $A$ un anneau. Soit $k$ un corps. Soit $\mathfrak p_n$, $n \geq 1$,
une suite d'idéaux premiers deux à deux distincts de $A$. De plus, pour chaque
$n$, soit $k \to \kappa(\mathfrak p_n)$ un plongement. Alors l'adhérence
de l'image de
$$
\coprod\nolimits_{n \not = m}
\Spec(\kappa(\mathfrak p_n) \otimes_k \kappa(\mathfrak p_m))
\longrightarrow
\Spec(A \otimes A)
$$
rencontre la diagonale.
\end{lemma}

\begin{proof}
Posons $k_n = \kappa(\mathfrak p_n)$. Nous pouvons supposer que
$A = \prod k_n$. Notons $x_n = \Spec(k_n)$ le point ouvert et fermé
correspondant à $A \to k_n$. Alors $\Spec(A) = Z \amalg \{x_n\}$, où $Z$
est une partie fermée non vide. En effet, $Z = V(e_n; n \geq 1)$, où $e_n$
est l'idempotent de $A$ correspondant au facteur $k_n$, et $Z$ est non vide
car l'idéal engendré par les $e_n$ n'est pas égal à $A$. Nous allons montrer
que l'adhérence de l'image contient $\Delta(Z)$. Le noyau de l'application
$$
(\prod k_n) \otimes_k (\prod k_m)
\longrightarrow
\prod\nolimits_{n \not = m} k_n \otimes_k k_m
$$
est l'idéal engendré par les $e_n \otimes e_n$, $n \geq 1$.
Par conséquent, l'adhérence de l'image du morphisme induit sur les spectres
est $V(e_n \otimes e_n; n \geq 1)$, dont l'intersection avec
$\Delta(\Spec(A))$ est $\Delta(Z)$. Il suffit donc de montrer que
$$
\coprod\nolimits_{n \not = m} \Spec(k_n \otimes_k k_m)
\longrightarrow
\Spec(\prod\nolimits_{n \not = m} k_n \otimes_k k_m)
$$
a une image dense. Cela résulte du fait que la famille d'homomorphismes
d'anneaux
$\prod_{n \not = m} k_n \otimes_k k_m \to k_n \otimes_k k_m$
est conjointement injective.
\end{proof}

\begin{lemma}[David Rydh]
\label{lemma-locally-separated-decent}
Un espace algébrique localement séparé est décent.
\end{lemma}

\begin{proof}
Soit $S$ un schéma et soit $X$ un espace algébrique localement séparé sur
$S$. Nous pouvons supposer que $S = \Spec(\mathbf{Z})$, voir
Propriétés des espaces, Définition
\ref{spaces-properties-definition-separated}.
Les produits fibrés sans indication seront pris au-dessus de $\mathbf{Z}$.
Soit $x \in |X|$. Choisissons un schéma $U$, un morphisme étale
$U \to X$ et un point $u \in U$ qui s'envoie sur $x$ dans $|X|$.
Comme d'habitude, nous identifions $u = \Spec(\kappa(u))$.
Puisque $X$ est localement séparé, le morphisme
$$
u \times_X u \to u \times u
$$
est une immersion (Morphismes d'espaces, lemme
\ref{spaces-morphisms-lemma-fibre-product-after-map}).
Par conséquent, Compléments sur les groupoïdes, lemme
\ref{more-groupoids-lemma-locally-closed-image-is-closed}
nous dit que c'est une immersion fermée (en utilisant
Schémas, lemme \ref{schemes-lemma-immersion-when-closed}).
Puisque $u \times_X u \to u \times_X U$ est un monomorphisme (changement de
base de $u \to U$) et que $u \times_X U \to u$ est étale, nous concluons que
$u \times_X u$ est une réunion disjointe de spectres de corps
(voir la Remarque \ref{remark-recall} et
Schémas, lemme \ref{schemes-lemma-mono-towards-spec-field}).
Puisqu'elle est aussi fermée dans le schéma affine $u \times u$, nous
concluons que $u \times_X u$ est une réunion disjointe finie de spectres de
corps. Ainsi, $x$ peut être représenté par un monomorphisme
$\Spec(k) \to X$, où $k$ est un corps, voir le
Lemme \ref{lemma-R-finite-above-x}.

\medskip\noindent
Ensuite, soit $U = \Spec(A)$ un schéma affine et soit $U \to X$ un morphisme
étale. Pour terminer la preuve, il suffit de montrer que
$F = U \times_X \Spec(k)$ est fini. Écrivons
$F = \coprod_{i \in I} \Spec(k_i)$ comme la réunion disjointe de spectres
d'extensions finies séparables de $k$. Nous devons montrer que $I$ est fini.
Posons $R = U \times_X U$. Puisque $X$ est localement séparé, le morphisme
$j : R \to U \times U$ est une immersion. Soit $U' \subset U \times U$ un
ouvert tel que $j$ se factorise par une immersion fermée $j' : R \to U'$.
Soit $e : U \to R$ l'application diagonale. Comme $e$ est un morphisme entre
schémas étales sur $U$ tel que $\Delta = j \circ e$ soit une immersion fermée,
nous concluons que $R = e(U) \amalg W$ pour un sous-schéma ouvert et fermé
$W \subset R$. Puisque $j'$ est une immersion fermée, nous concluons que
$j'(W) \subset U'$ est fermé et disjoint de $j'(e(U))$. Par conséquent,
$\overline{j(W)} \cap \Delta(U) = \emptyset$ dans $U \times U$.
Remarquons que $W$ contient $\Spec(k_i \otimes_k k_{i'})$ pour tous
$i \not = i'$, $i, i' \in I$. D'après le Lemme
\ref{lemma-infinite-number}, nous concluons que $I$ est fini, comme voulu.
\end{proof}








\section{Critère valuatif}
\label{section-valuative-criterion-universally-closed}

\noindent
Pour un morphisme quasi-compact de source un espace décent, le critère
valuatif est une condition nécessaire pour que le morphisme soit
universellement fermé.


\begin{proposition}
\label{proposition-characterize-universally-closed}
Soit $S$ un schéma. Soit $f : X \to Y$ un morphisme d'espaces algébriques
sur $S$. Supposons que $f$ soit quasi-compact et que $X$ soit décent.
Alors $f$ est
universellement fermé si et seulement si la partie d'existence du critère
valuatif est satisfaite.
\end{proposition}

\begin{proof}
Dans
Morphismes d'espaces,
lemme \ref{spaces-morphisms-lemma-quasi-compact-existence-universally-closed},
nous avons vu l'une des implications.
Pour démontrer l'autre, supposons $f$ universellement fermé. Soit
$$
\xymatrix{
\Spec(K) \ar[r] \ar[d] & X \ar[d] \\
\Spec(A) \ar[r] & Y
}
$$
un diagramme comme dans
Morphismes d'espaces,
définition \ref{spaces-morphisms-definition-valuative-criterion}.
Posons $X_A = \Spec(A) \times_Y X$ ; on a alors
$$
\xymatrix{
\Spec(K) \ar[r] \ar[rd] & X_A \ar[d] \\
 & \Spec(A)
}
$$
D'après
Morphismes d'espaces,
lemme \ref{spaces-morphisms-lemma-base-change-quasi-compact},
$X_A \to \Spec(A)$ est quasi-compact. Comme $X_A \to X$
est représentable, $X_A$ est lui aussi décent ; voir le
Lemme \ref{lemma-representable-properties}.
De plus, puisque $f$ est universellement fermé, $X_A \to \Spec(A)$
est universellement fermé.
On peut donc remplacer, et l'on remplace, $X$ par $X_A$ et $Y$ par $\Spec(A)$.

\medskip\noindent
Soit $x' \in |X|$ la classe d'équivalence du morphisme
$\Spec(K) \to X$. Soit $y \in |Y| = |\Spec(A)|$
le point fermé. Posons $y' = f(x')$ ; c'est le point générique de
$\Spec(A)$. Comme $f$ est universellement fermé,
$f(\overline{\{x'\}})$ contient $\overline{\{y'\}}$, donc
contient $y$. Soit $x \in \overline{\{x'\}}$ un point tel que
$f(x) = y$. Soit $U$ un schéma, et soit $\varphi : U \to X$
un morphisme étale tel qu'il existe un point $u \in U$ avec
$\varphi(u) = x$. D'après le
Lemme \ref{lemma-specialization}
et l'hypothèse selon laquelle $X$ est décent,
il existe une spécialisation $u' \leadsto u$ dans $U$ telle que
$\varphi(u') = x'$.
Cela signifie qu'il existe une extension de corps commune
$K \subset K' \supset \kappa(u')$ telle que le diagramme
$$
\xymatrix{
\Spec(K') \ar[r] \ar[d] & U \ar[d] \\
\Spec(K) \ar[r] \ar[rd] & X \ar[d] \\
 & \Spec(A)
}
$$
soit commutatif. On obtient le diagramme commutatif d'anneaux suivant :
$$
\xymatrix{
K' & \mathcal{O}_{U, u} \ar[l] \\
K \ar[u] & \\
 & A \ar[lu] \ar[uu]
}
$$
D'après
Algèbre, lemme \ref{algebra-lemma-dominate},
on peut trouver un anneau de valuation $A' \subset K'$ qui domine l'image de
$\mathcal{O}_{U, u}$ dans $K'$. Comme, par construction, $\mathcal{O}_{U, u}$
domine $A$, on voit que $A'$ domine également $A$. On obtient donc un
diagramme
semblable au deuxième diagramme de
Morphismes d'espaces,
définition \ref{spaces-morphisms-definition-valuative-criterion},
ce qui démontre la proposition.
\end{proof}







\section{Conditions relatives}
\label{section-relative-conditions}

\noindent
Il s'agit d'une section technique (encore une) consacrée à des conditions
sur les espaces algébriques qui concernent les points. Il est probablement
préférable de sauter cette section en première lecture.

\begin{definition}
\label{definition-relative-conditions}
Soit $S$ un schéma. On dit qu'un espace algébrique $X$ sur $S$
{\it possède la propriété $(\beta)$} si $X$ possède la propriété
correspondante
du Lemme \ref{lemma-bounded-fibres}.
Soit $f : X \to Y$ un morphisme d'espaces algébriques sur $S$.
\begin{enumerate}
\item On dit que $f$ {\it possède la propriété $(\beta)$} si, pour tout schéma
$T$ et tout morphisme $T \to Y$, le produit fibré $T \times_Y X$ possède la
propriété $(\beta)$.
\item On dit que $f$ est {\it décent} si, pour tout schéma $T$ et tout
morphisme $T \to Y$, le produit fibré $T \times_Y X$ est un espace
algébrique décent.
\item On dit que $f$ est {\it raisonnable} si, pour tout schéma $T$ et tout
morphisme $T \to Y$, le produit fibré $T \times_Y X$ est un espace
algébrique raisonnable.
\item On dit que $f$ est {\it très raisonnable} si, pour tout schéma $T$ et
tout morphisme $T \to Y$, le produit fibré $T \times_Y X$ est un espace
algébrique très raisonnable.
\end{enumerate}
\end{definition}

\noindent
Nous renvoyons à la Remarque \ref{remark-very-reasonable} pour une discussion
informelle. On verra que la classe des morphismes très raisonnables n'est pas
très utile, mais que celles des morphismes décents et raisonnables le sont.

\begin{lemma}
\label{lemma-properties-trivial-implications}
Soit $S$ un schéma.
Soit $f : X \to Y$ un morphisme d'espaces algébriques sur $S$.
On a les implications suivantes entre les conditions imposées à $f$ :
$$
\xymatrix{
\text{représentable} \ar@{=>}[rd] & & & & \\
& \text{très raisonnable} \ar@{=>}[r] & \text{raisonnable} \ar@{=>}[r] &
\text{décent} \ar@{=>}[r] & (\beta) \\
\text{quasi-séparé} \ar@{=>}[ru] & & & &
}
$$
\end{lemma}

\begin{proof}
Cela résulte immédiatement des définitions, du
Lemme \ref{lemma-bounded-fibres}
et de
Morphismes d'espaces,
Lemme \ref{spaces-morphisms-lemma-separated-local}.
\end{proof}

\noindent
Voici une autre vérification élémentaire.

\begin{lemma}
\label{lemma-property-for-morphism-out-of-property}
Soit $S$ un schéma. Soit $f : X \to Y$ un morphisme d'espaces
algébriques sur $S$. Si $X$ est décent (resp.\ raisonnable, resp.\ possède la
propriété $(\beta)$ du Lemme \ref{lemma-bounded-fibres}), alors $f$ est décent
(resp.\ raisonnable, resp.\ possède la propriété $(\beta)$).
\end{lemma}

\begin{proof}
Soit $T$ un schéma et soit $T \to Y$ un morphisme. Alors $T \to Y$ est
représentable, donc le changement de base $T \times_Y X \to X$ est
représentable. Par conséquent, si $X$ est décent (ou raisonnable), il en est
de
même de $T \times_Y X$ ; voir le
Lemme \ref{lemma-representable-named-properties}.
De même, pour la propriété $(\beta)$, voir le
Lemme \ref{lemma-representable-properties}.
\end{proof}

\begin{lemma}
\label{lemma-base-change-relative-conditions}
Le fait de posséder la propriété $(\beta)$, d'être décent ou d'être
raisonnable est stable par changement de base arbitraire.
\end{lemma}

\begin{proof}
Cela découle immédiatement de la définition.
\end{proof}

\begin{lemma}
\label{lemma-property-over-property}
Soit $S$ un schéma.
Soit $f : X \to Y$ un morphisme d'espaces algébriques sur $S$.
Soit $\omega \in \{\beta, \text{décent}, \text{raisonnable}\}$.
Supposons que $Y$ possède la propriété $(\omega)$ et que $f : X \to Y$
possède $(\omega)$. Alors $X$ possède $(\omega)$.
\end{lemma}

\begin{proof}
Démontrons le lemme dans le cas où $\omega = \beta$. Il faut alors montrer
que tout $x \in |X|$ est représenté par un monomorphisme du spectre d'un
corps vers $X$. Posons $y = f(x) \in |Y|$. Par hypothèse, il existe un corps
$k$ et un monomorphisme $\Spec(k) \to Y$ représentant $y$.
Alors $x$ correspond à un point $x'$ de $\Spec(k) \times_Y X$.
Par hypothèse, $x'$ est représenté par un monomorphisme
$\Spec(k') \to \Spec(k) \times_Y X$. Il est clair que le composé
$\Spec(k') \to X$ est un monomorphisme représentant $x$.

\medskip\noindent
Démontrons le lemme dans le cas où $\omega = \text{décent}$.
Soient $x \in |X|$ et $y = f(x) \in |Y|$. D'après le paragraphe précédent,
on peut choisir un diagramme
$$
\xymatrix{
\Spec(k') \ar[r]_x \ar[d] & X \ar[d]^f \\
\Spec(k) \ar[r]^y & Y
}
$$
dont les flèches horizontales sont des monomorphismes. Puisque $Y$ est décent,
le morphisme $y$ est quasi-compact. Puisque $f$ est décent, l'espace
algébrique $\Spec(k) \times_Y X$ est décent. Le monomorphisme
$\Spec(k') \to \Spec(k) \times_Y X$ est donc quasi-compact.
Le monomorphisme $x : \Spec(k') \to X$ est alors quasi-compact en tant que
composé de morphismes quasi-compacts (utiliser
Morphismes d'espaces, lemmes
\ref{spaces-morphisms-lemma-base-change-quasi-compact} et
\ref{spaces-morphisms-lemma-composition-quasi-compact}).
Comme le point $x$ était arbitraire, il s'ensuit que $X$ est décent.

\medskip\noindent
Démontrons le lemme dans le cas où $\omega = \text{raisonnable}$.
Choisissons $V \to Y$ étale, avec $V$ un schéma affine.
Choisissons $U \to V \times_Y X$ étale, avec $U$ un schéma affine.
Par hypothèse, les fibres de $V \to Y$ sont universellement bornées.
D'après le
Lemme \ref{lemma-base-change-universally-bounded},
les fibres du morphisme $V \times_Y X \to X$ sont universellement bornées.
L'hypothèse sur $f$ montre que les fibres de $U \to V \times_Y X$ sont
universellement bornées. D'après le
Lemme \ref{lemma-composition-universally-bounded},
les fibres du composé $U \to X$ sont universellement bornées.
Il existe donc suffisamment de morphismes étales $U \to X$ provenant de
schémas dont les fibres sont universellement bornées, et l'on conclut que
$X$ est raisonnable.
\end{proof}

\begin{lemma}
\label{lemma-composition-relative-conditions}
Le fait de posséder la propriété $(\beta)$, d'être décent ou d'être
raisonnable est stable par composition.
\end{lemma}

\begin{proof}
Soit $\omega \in \{\beta, \text{décent}, \text{raisonnable}\}$.
Soient $f : X \to Y$ et $g : Y \to Z$ des morphismes d'espaces algébriques
sur le schéma $S$. Supposons que $f$ et $g$ possèdent tous deux la propriété
$(\omega)$. Il faut alors montrer que, pour tout schéma $T$ et tout morphisme
$T \to Z$, l'espace $T \times_Z X$ possède $(\omega)$. D'après le
Lemme \ref{lemma-base-change-relative-conditions},
on se ramène à l'assertion suivante : si $Y$ est un espace algébrique qui
possède la propriété $(\omega)$ et si $f : X \to Y$ est un morphisme qui
possède $(\omega)$, alors $X$ possède $(\omega)$.
C'est le contenu du Lemme \ref{lemma-property-over-property}.
\end{proof}

\begin{lemma}
\label{lemma-fibre-product-conditions}
Soit $S$ un schéma. Soient $f : X \to Y$ et $g : Z \to Y$ des morphismes
d'espaces algébriques sur $S$. Si $X$ et $Z$ sont décents
(resp.\ raisonnables, resp.\ possèdent la propriété
$(\beta)$ du Lemme \ref{lemma-bounded-fibres}), il en va de même de
$X \times_Y Z$.
\end{lemma}

\begin{proof}
En effet, d'après le Lemme \ref{lemma-property-for-morphism-out-of-property},
le morphisme $X \to Y$ possède la propriété considérée. Son changement de
base $X \times_Y Z \to Z$ possède alors cette propriété d'après le
Lemme \ref{lemma-base-change-relative-conditions}.
Enfin, le Lemme \ref{lemma-property-over-property} implique que
$X \times_Y Z$ possède cette propriété.
\end{proof}

\begin{lemma}
\label{lemma-descent-conditions}
Soit $S$ un schéma.
Soit $f : X \to Y$ un morphisme d'espaces algébriques sur $S$.
Soit $\mathcal{P} \in \{(\beta), \text{décent}, \text{raisonnable}\}$.
Supposons que
\begin{enumerate}
\item $f$ soit quasi-compact,
\item $f$ soit étale,
\item $|f| : |X| \to |Y|$ soit surjective, et
\item l'espace algébrique $X$ possède la propriété $\mathcal{P}$.
\end{enumerate}
Alors $Y$ possède la propriété $\mathcal{P}$.
\end{lemma}

\begin{proof}
Démontrons ceci dans le cas où $\mathcal{P} = (\beta)$. Soit $y \in |Y|$
un point. Il faut montrer que $y$ peut être représenté par un monomorphisme
provenant d'un corps. Choisissons un point $x \in |X|$ tel que $f(x) = y$.
Par hypothèse, on peut représenter $x$ par un monomorphisme
$\Spec(k) \to X$, où $k$ est un corps. D'après le
Lemme \ref{lemma-R-finite-above-x},
il suffit de montrer que les projections
$\Spec(k) \times_Y \Spec(k) \to \Spec(k)$
sont étales et quasi-compactes. On peut factoriser la première projection en
$$
\Spec(k) \times_Y \Spec(k)
\longrightarrow
\Spec(k) \times_Y X
\longrightarrow
\Spec(k)
$$
Le premier morphisme est un monomorphisme, et le second est étale et
quasi-compact. D'après
Propriétés des espaces,
Lemme \ref{spaces-properties-lemma-etale-over-field-scheme},
$\Spec(k) \times_Y X$ est un schéma. C'est donc une réunion disjointe finie
de spectres d'extensions finies séparables de $k$. D'après
Schémas, lemme \ref{schemes-lemma-mono-towards-spec-field},
la première flèche identifie $\Spec(k) \times_Y \Spec(k)$ à une réunion
disjointe finie de spectres d'extensions finies séparables de $k$.
Le morphisme de projection est donc étale et quasi-compact.

\medskip\noindent
Démontrons ceci dans le cas où $\mathcal{P} = \text{décent}$.
Le premier paragraphe de la preuve montre déjà que tout $y \in |Y|$ peut être
représenté par un monomorphisme $y : \Spec(k) \to Y$. Fixons un tel $y$.
Choisissons un schéma affine $U$ et un morphisme étale $U \to X$ tels que
l'image de $|U| \to |Y|$ contienne $y$. D'après le
Lemme \ref{lemma-UR-finite-above-x},
il suffit de montrer que $U_y$ est un schéma fini sur $k$. Le produit fibré
$X_y = \Spec(k) \times_Y X$ est un espace algébrique quasi-compact et étale
sur $k$. D'après
Propriétés des espaces,
Lemme \ref{spaces-properties-lemma-etale-over-field-scheme},
c'est donc un schéma. C'est une réunion disjointe finie de spectres
d'extensions finies séparables de $k$. Écrivons
$X_y = \{x_1, \ldots, x_n\}$, où $x_i$ est donné par
$x_i : \Spec(k_i) \to X$ avec $[k_i : k] < \infty$. Par hypothèse, $X$ est
décent, donc les schémas $U_{x_i} = \Spec(k_i) \times_X U$ sont finis sur
$k_i$. Enfin, $U_y = \coprod U_{x_i}$ comme schéma, d'où l'on conclut que
$U_y$ est fini sur $k$, comme voulu.

\medskip\noindent
Démontrons ceci dans le cas où $\mathcal{P} = \text{raisonnable}$.
Choisissons un schéma affine $V$ et un morphisme étale $V \to Y$.
Il faut montrer que les fibres de $V \to Y$ sont universellement bornées.
L'espace algébrique $V \times_Y X$ est quasi-compact.
On peut donc trouver un schéma affine $W$ et un morphisme étale surjectif
$W \to V \times_Y X$ ; voir
Propriétés des espaces,
Lemme \ref{spaces-properties-lemma-quasi-compact-affine-cover}.
Voici un diagramme, dont les flèches pleines sont données :
$$
\xymatrix{
W \ar[r]  \ar[rd] &
V \times_Y X \ar[r] \ar[d] &
X \ar[d]_f & \Spec(k) \ar@{..>}[l]^x \ar@{..>}[ld]^y \\
 & V \ar[r] & Y
}
$$
Les fibres du morphisme $W \to X$ sont universellement bornées, puisque
l'espace $X$ est raisonnable. Soit $n$ un entier qui majore les degrés des
fibres de $W \to X$. Nous affirmons que le même entier majore les fibres de
$V \to Y$. En effet, soit $y \in |Y|$ un point. Il existe alors un point
$x \in |X|$ tel que $f(x) = y$ (voir ci-dessus).
On peut donc trouver un corps $k$ et des morphismes $x, y$ donnés par les
flèches pointillées du diagramme ci-dessus. En particulier, on obtient un
morphisme étale surjectif
$$
\Spec(k) \times_{x, X} W
\to
\Spec(k) \times_{x, X} (V \times_Y X) = \Spec(k) \times_{y, Y} V
$$
ce qui montre que le degré de $\Spec(k) \times_{y, Y} V$ sur $k$ est
inférieur ou égal au degré de $\Spec(k) \times_{x, X} W$ sur $k$, donc qu'il
est $\leq n$, et cela donne le résultat. (Cette dernière partie de
l'argument est la même que dans la preuve du
Lemme \ref{lemma-descent-universally-bounded}.
Malheureusement, ce lemme n'est pas assez général, car il ne s'applique
qu'aux morphismes représentables.)
\end{proof}

\begin{lemma}
\label{lemma-relative-conditions-local}
Soit $S$ un schéma.
Soit $f : X \to Y$ un morphisme d'espaces algébriques sur $S$.
Soit $\mathcal{P} \in \{(\beta), \text{décent}, \text{raisonnable}, \text{très raisonnable}\}$.
Les conditions suivantes sont équivalentes :
\begin{enumerate}
\item $f$ possède $\mathcal{P}$,
\item pour tout schéma affine $Z$ et tout morphisme $Z \to Y$, le changement
de base $Z \times_Y X \to Z$ de $f$ possède $\mathcal{P}$,
\item pour tout schéma affine $Z$ et tout morphisme $Z \to Y$, l'espace
algébrique $Z \times_Y X$ possède $\mathcal{P}$, et
\item il existe un recouvrement de Zariski $Y = \bigcup Y_i$ tel que chaque
morphisme $f^{-1}(Y_i) \to Y_i$ possède $\mathcal{P}$.
\end{enumerate}
Si $\mathcal{P} \in \{(\beta), \text{décent}, \text{raisonnable}\}$, ces conditions sont
aussi
équivalentes à la suivante :
\begin{enumerate}
\item[(5)] il existe un schéma $V$ et un morphisme étale surjectif
$V \to Y$ tels que le changement de base $V \times_Y X \to V$ possède
$\mathcal{P}$.
\end{enumerate}
\end{lemma}

\begin{proof}
Les implications (1) $\Rightarrow$ (2) $\Rightarrow$ (3) $\Rightarrow$ (4)
sont immédiates.
Pour voir que (3) $\Rightarrow$ (1), procédons comme suit.
Soit $Z \to Y$ un morphisme dont la source est un schéma sur $S$.
Considérons l'espace algébrique $Z \times_Y X$. Si l'on suppose (3), alors,
pour tout ouvert affine $W \subset Z$, le sous-espace ouvert
$W \times_Y X$ de $Z \times_Y X$ possède la propriété $\mathcal{P}$.
D'après le
Lemme \ref{lemma-properties-local},
l'espace $Z \times_Y X$ possède donc la propriété $\mathcal{P}$, ce qui donne
(1). Un argument analogue, que nous omettons, montre que (4) implique (1).

\medskip\noindent
L'implication (1) $\Rightarrow$ (5) est immédiate. Soit $V \to Y$ un morphisme
étale provenant d'un schéma comme dans (5). Soit $Z$ un schéma affine et soit
$Z \to Y$ un morphisme. Considérons le diagramme
$$
\xymatrix{
Z \times_Y V \ar[r]_q \ar[d]_p & V \ar[d] \\
Z \ar[r] & Y
}
$$
Puisque $p$ est étale, donc ouvert, on peut choisir un nombre fini de
sous-schémas ouverts affines $W_i \subset Z \times_Y V$ tels que
$Z = \bigcup p(W_i)$. Considérons le diagramme commutatif
$$
\xymatrix{
V \times_Y X \ar[d] &
(\coprod W_i) \times_Y X \ar[l] \ar[d] \ar[r] &
Z \times_Y X \ar[d] \\
V &
\coprod W_i \ar[l] \ar[r] &
Z
}
$$
On sait que $V \times_Y X$ possède la propriété $\mathcal{P}$. D'après le
Lemme \ref{lemma-representable-properties},
$(\coprod W_i) \times_Y X$ possède la propriété $\mathcal{P}$.
Remarquons que le morphisme
$(\coprod W_i) \times_Y X \to Z \times_Y X$ est étale et quasi-compact,
puisqu'il est le changement de base de $\coprod W_i \to Z$.
Le Lemme \ref{lemma-descent-conditions}
montre donc que $Z \times_Y X$ possède la propriété $\mathcal{P}$.
\end{proof}

\begin{remark}
\label{remark-very-reasonable}
Une description informelle des propriétés $(\beta)$, décent, raisonnable et
très raisonnable a été donnée dans la section \ref{section-reasonable-decent}.
Un morphisme possède l'une de ces propriétés si, de manière très
approximative,
ses fibres possèdent les propriétés correspondantes.
La propriété d'être décent est utile pour démontrer des résultats sur les
spécialisations de points de $|X|$. La propriété d'être raisonnable est un peu
plus forte et se prête assez facilement aux raisonnements techniques.
\end{remark}

\noindent
Voici un lemme promis plus haut qui utilise les morphismes décents.

\begin{lemma}
\label{lemma-re-characterize-universally-closed}
Soit $S$ un schéma.
Soit $f : X \to Y$ un morphisme d'espaces algébriques sur $S$.
Supposons $f$ quasi-compact et décent.
(C'est par exemple le cas si $f$ est représentable ou quasi-séparé ; voir le
Lemme \ref{lemma-properties-trivial-implications}.)
Alors $f$ est universellement fermé si et seulement si la partie d'existence
du critère valuatif est satisfaite.
\end{lemma}

\begin{proof}
Dans
Morphismes d'espaces,
Lemme \ref{spaces-morphisms-lemma-quasi-compact-existence-universally-closed},
nous avons démontré que tout morphisme quasi-compact qui satisfait la partie
d'existence du critère valuatif est universellement fermé.
Pour démontrer la réciproque, supposons que $f$ soit universellement fermé.
Dans la preuve de la
Proposition \ref{proposition-characterize-universally-closed},
nous avons vu qu'il suffit de montrer, pour tout anneau de valuation $A$ et
tout morphisme $\Spec(A) \to Y$, que le changement de base
$f_A : X_A \to \Spec(A)$ satisfait la partie d'existence du critère valuatif.
Par définition, l'espace algébrique $X_A$ possède la propriété $(\gamma)$ ; la
Proposition \ref{proposition-characterize-universally-closed}
s'applique donc au morphisme $f_A$, ce qui achève la preuve.
\end{proof}





\section{Points des fibres}
\label{section-points-fibres}

\noindent
Soit $S$ un schéma. Considérons un diagramme cartésien
\begin{equation}
\label{equation-points-fibres}
\xymatrix{
W \ar[r]_q \ar[d]_p & Z \ar[d]^g \\
X \ar[r]^f & Y
}
\end{equation}
d'espaces algébriques sur $S$. Soient $x \in |X|$ et $z \in |Z|$
des points qui s'envoient sur le même point $y \in |Y|$. On peut se demander :
quand l'ensemble
\begin{equation}
\label{equation-fibre}
F_{x, z} = \{ w \in |W| \text{ tel que }p(w) = x\text{ et }q(w) = z\}
\end{equation}
est-il fini ?

\begin{example}
\label{example-schemes}
Si $X, Y, Z$ sont des schémas, alors l'ensemble $F_{x, z}$ s'identifie à
l'ensemble sous-jacent au spectre de
$\kappa(x) \otimes_{\kappa(y)} \kappa(z)$
(Schémas, lemme \ref{schemes-lemma-points-fibre-product}). On obtient donc un
ensemble fini si l'extension $\kappa(y) \subset \kappa(x)$ est finie ou si
l'extension $\kappa(y) \subset \kappa(z)$ est finie. C'est en particulier
toujours le cas si $g$ est quasi-fini en $z$ (Morphismes, lemme
\ref{morphisms-lemma-residue-field-quasi-finite}).
\end{example}

\begin{example}
\label{example-not-finite}
Soit $K$ un corps de caractéristique $0$ muni d'un automorphisme $\sigma$
d'ordre infini. Posons $Y = \Spec(K)/\mathbf{Z}$ et
$X = \mathbf{A}^1_K/\mathbf{Z}$, où $\mathbf{Z}$ agit sur $K$ par $\sigma$
et sur $\mathbf{A}^1_K = \Spec(K[t])$ par $t \mapsto t + 1$.
Soit $Z = \Spec(K)$. Alors $W = \mathbf{A}^1_K$. Voici le diagramme :
$$
\xymatrix{
\mathbf{A}^1_K \ar[r]_q \ar[d]_p & \Spec(K) \ar[d]^g \\
\mathbf{A}^1_K/\mathbf{Z} \ar[r]^f & \Spec(K)/\mathbf{Z}
}
$$
Prenons pour $x$ le point correspondant à $t = 0$ et pour $z$ l'unique point
de $\Spec(K)$. On voit alors que $F_{x, z} = \mathbf{Z}$ en tant qu'ensemble.
\end{example}

\begin{lemma}
\label{lemma-surjective-on-fibres}
Dans la situation de (\ref{equation-points-fibres}), si $Z' \to Z$ est un
morphisme et si $z' \in |Z'|$ s'envoie sur $z$, alors l'application induite
$F_{x, z'} \to F_{x, z}$ est surjective.
\end{lemma}

\begin{proof}
Posons $W' = X \times_Y Z' = W \times_Z Z'$. Alors l'application
$|W'| \to |W| \times_{|Z|} |Z'|$ est surjective d'après
Propriétés des espaces, lemme
\ref{spaces-properties-lemma-points-cartesian}.
D'où la surjectivité de $F_{x, z'} \to F_{x, z}$.
\end{proof}

\begin{lemma}
\label{lemma-qf-and-qc-finite-fibre}
Dans le diagramme (\ref{equation-points-fibres}), l'ensemble
(\ref{equation-fibre}) est fini si $f$ est de type fini et si $f$ est
quasi-fini en $x$.
\end{lemma}

\begin{proof}
Le morphisme $q$ est quasi-fini en tout $w \in F_{x, z}$ ; voir
Morphismes d'espaces, lemme
\ref{spaces-morphisms-lemma-base-change-quasi-finite-locus}.
Le lemme résulte donc de
Morphismes d'espaces, lemme
\ref{spaces-morphisms-lemma-quasi-finite-at-a-finite-number-of-points}.
\end{proof}

\begin{lemma}
\label{lemma-decent-finite-fibre}
Dans le diagramme (\ref{equation-points-fibres}), l'ensemble
(\ref{equation-fibre}) est fini si $y$ peut être représenté par un
monomorphisme $\Spec(k) \to Y$, où $k$ est un corps, et si $g$ est quasi-fini
en $z$. (Cas particulier : $Y$ est décent et $g$ est étale.)
\end{lemma}

\begin{proof}
En appliquant deux fois le Lemme \ref{lemma-surjective-on-fibres},
on peut remplacer $Z$ par $Z_k = \Spec(k) \times_Y Z$ et
$X$ par $X_k = \Spec(k) \times_Y X$. On peut aussi remplacer, et l'on
remplace, $Y$ par $\Spec(k)$. Remarquons que $Z_k \to \Spec(k)$ est quasi-fini
en $z$ d'après Morphismes d'espaces, lemme
\ref{spaces-morphisms-lemma-base-change-quasi-finite-locus}.
Choisissons un schéma $V$, un point $v \in V$ et un morphisme étale
$V \to Z_k$ qui envoie $v$ sur $z$. Choisissons un schéma $U$, un point
$u \in U$ et un morphisme étale $U \to X_k$ qui envoie $u$ sur $x$.
D'après le Lemme \ref{lemma-surjective-on-fibres},
il suffit encore de montrer que $F_{u, v}$ est fini dans le diagramme
$$
\xymatrix{
U \times_{\Spec(k)} V \ar[r] \ar[d] & V \ar[d] \\
U \ar[r] & \Spec(k)
}
$$
Le morphisme $V \to \Spec(k)$ est quasi-fini en $v$
(cela résulte de la discussion générale de
Morphismes d'espaces, section
\ref{spaces-morphisms-section-local-source-target},
et de la définition de la quasi-finitude en un point).
La finitude résulte alors de l'Exemple \ref{example-schemes}.
La remarque entre parenthèses de l'énoncé découle du fait que les points des
espaces décents sont représentés par des monomorphismes provenant de corps,
et du fait qu'un morphisme étale d'espaces algébriques est localement
quasi-fini.
\end{proof}

\begin{lemma}
\label{lemma-topology-fibre}
\begin{slogan}
Les fibres des points à valeurs dans un corps des espaces algébriques ont la
topologie de Zariski attendue.
\end{slogan}
Soit $S$ un schéma. Soit $f : X \to Y$ un morphisme d'espaces algébriques
sur $S$.
Soit $y \in |Y|$ et supposons que $y$ soit représenté par un monomorphisme
quasi-compact $\Spec(k) \to Y$. Alors $|X_k| \to |X|$ est un homéomorphisme
sur $f^{-1}(\{y\}) \subset |X|$, muni de la topologie induite.
\end{lemma}

\begin{proof}
Nous utiliserons
Propriétés des espaces, lemme \ref{spaces-properties-lemma-etale-open},
ainsi que
Morphismes d'espaces, lemme
\ref{spaces-morphisms-lemma-monomorphism-injective-points},
sans autre mention.
Soit $V \to Y$ un morphisme étale, avec $V$ affine, tel qu'il existe un point
$v \in V$ qui s'envoie sur $y$. Comme $\Spec(k) \to Y$ est quasi-compact,
seul un nombre fini de points de $V$ s'envoient sur $y$
(Lemme \ref{lemma-UR-finite-above-x}). Après avoir rétréci
$V$, on peut supposer que $v$ est le seul. Choisissons un schéma $U$ et un
morphisme étale surjectif $U \to X$.
Considérons le diagramme commutatif
$$
\xymatrix{
U \ar[d] & U_V \ar[l] \ar[d] & U_v \ar[l] \ar[d] \\
X \ar[d] & X_V \ar[l] \ar[d] & X_v \ar[l] \ar[d] \\
Y & V \ar[l] & v \ar[l]
}
$$
Comme $U_v \to U_V$ identifie $U_v$ à une partie de $U_V$ munie de la
topologie induite (Schémas, lemme
\ref{schemes-lemma-fibre-topological}),
et comme $|U_V| \to |X_V|$ et $|U_v| \to |X_v|$ sont surjectives et ouvertes,
on voit que $|X_v| \to |X_V|$ est un homéomorphisme sur son image, munie de la
topologie induite.
D'autre part, l'image inverse de $f^{-1}(\{y\})$ par l'application ouverte
$|X_V| \to |X|$ est égale à $|X_v|$.
On en conclut que $|X_v| \to f^{-1}(\{y\})$ est ouverte.
Le morphisme $X_v \to X$ se factorise par $X_k$, et l'application
$|X_k| \to |X|$ est injective, d'image $f^{-1}(\{y\})$, d'après
Propriétés des espaces, lemme
\ref{spaces-properties-lemma-points-cartesian}. En utilisant
$|X_v| \to |X_k| \to f^{-1}(\{y\})$, le lemme résulte de la surjectivité de
$X_v \to X_k$.
\end{proof}

\begin{lemma}
\label{lemma-conditions-on-point-on-space-over-field}
Soit $X$ un espace algébrique localement de type fini sur un corps $k$.
Soit $x \in |X|$. Considérons les conditions suivantes :
\begin{enumerate}
\item $\dim_x(|X|) = 0$,
\item $x$ est fermé dans $|X|$ et, si $x' \leadsto x$  dans $|X|$, alors
$x' = x$,
\item $x$ est un point isolé de $|X|$,
\item $\dim_x(X) = 0$,
\item $X \to \Spec(k)$ est quasi-fini en $x$.
\end{enumerate}
Alors (2), (3), (4) et (5) sont équivalentes.
Si $X$ est décent, (1) est équivalente aux autres.
\end{lemma}

\begin{proof}
Les assertions (4) et (5) sont par exemple équivalentes d'après
Morphismes d'espaces, lemmes
\ref{spaces-morphisms-lemma-locally-finite-type-quasi-finite-part} et
\ref{spaces-morphisms-lemma-quasi-finite-at-point}.

\medskip\noindent
Soit $U \to X$ un morphisme étale, où $U$ est un schéma affine, et soit
$u \in U$ un point qui s'envoie sur $x$. De plus, si $x$ est un point fermé,
par exemple dans les cas (2) ou (3), on peut supposer, et l'on suppose, que
$u$ est fermé. Remarquons que $\dim_u(U) = \dim_x(X)$ par définition, et que
ce nombre est égal à $\dim(\mathcal{O}_{U, u})$ si $u$ est fermé ; voir
Algèbre, lemme
\ref{algebra-lemma-dimension-closed-point-finite-type-field}.

\medskip\noindent
Si $\dim_x(X) > 0$ et si $u$ est fermé, les arguments précédents permettent de
choisir une spécialisation non triviale $u' \leadsto u$ dans $U$.
Alors le degré de transcendance de $\kappa(u')$ sur $k$ est strictement
supérieur à celui de $\kappa(u)$ sur $k$. Il s'ensuit que les images $x$ et
$x'$ dans $X$ sont distinctes, car les degrés de transcendance de $x/k$ et de
$x'/k$ sont bien définis ; voir Morphismes d'espaces, définition
\ref{spaces-morphisms-definition-dimension-fibre}.
Cela s'applique notamment dans les cas (2) et (3), et l'on conclut que (2) et
(3) impliquent (4).

\medskip\noindent
Réciproquement, si $X \to \Spec(k)$ est localement quasi-fini en $x$, alors
$U \to \Spec(k)$ est localement quasi-fini en $u$, de sorte que $u$ est un
point isolé de $U$
(Morphismes, lemme \ref{morphisms-lemma-quasi-finite-at-point-characterize}).
Comme $|U| \to |X|$ est continue et ouverte, il s'ensuit que (5) implique (2)
et (3).

\medskip\noindent
Supposons que $X$ soit décent et que (1) soit satisfaite. Alors
$\dim_x(X) = \dim_x(|X|)$ d'après le
Lemme \ref{lemma-dimension-decent-space}, ce qui achève la preuve.
\end{proof}

\begin{lemma}
\label{lemma-conditions-on-space-over-field}
Soit $X$ un espace algébrique localement de type fini sur un corps $k$.
Considérons les conditions suivantes :
\begin{enumerate}
\item $|X|$ est un ensemble fini,
\item $|X|$ est un espace discret,
\item $\dim(|X|) = 0$,
\item $\dim(X) = 0$,
\item $X \to \Spec(k)$ est localement quasi-fini,
\end{enumerate}
Alors (2), (3), (4) et (5) sont équivalentes.
Si $X$ est décent, (1) implique les autres.
\end{lemma}

\begin{proof}
Les assertions (4) et (5) sont par exemple équivalentes d'après
Morphismes d'espaces, lemme
\ref{spaces-morphisms-lemma-locally-finite-type-quasi-finite-part}.

\medskip\noindent
Soit $U \to X$ un morphisme étale surjectif, où $U$ est un schéma.

\medskip\noindent
Si $\dim(U) > 0$, choisissons une spécialisation non triviale
$u \leadsto u'$ dans $U$. Le degré de transcendance de $\kappa(u)$ sur $k$
est strictement supérieur à celui de $\kappa(u')$ sur $k$.
Il s'ensuit que les images $x$ et $x'$ dans $X$ sont distinctes, car les
degrés
de transcendance de $x/k$ et de $x'/k$ sont bien définis ; voir
Morphismes d'espaces, définition
\ref{spaces-morphisms-definition-dimension-fibre}.
On conclut que (2) et (3) impliquent (4).

\medskip\noindent
Réciproquement, si $X \to \Spec(k)$ est localement quasi-fini, alors $U$ est
localement noethérien
(Morphismes, lemme \ref{morphisms-lemma-finite-type-noetherian}),
de dimension $0$
(Morphismes, lemme
\ref{morphisms-lemma-locally-quasi-finite-rel-dimension-0}),
et est donc une réunion disjointe de spectres d'anneaux locaux artiniens
(Propriétés, lemme \ref{properties-lemma-locally-Noetherian-dimension-0}).
Ainsi, $U$ est un espace topologique discret et, comme $|U| \to |X|$ est
continue et ouverte, il en va de même de $|X|$.
Autrement dit, (4) implique (2) et (3).

\medskip\noindent
Supposons que $X$ soit décent et que (1) soit satisfaite. On peut alors
choisir $U$ affine ci-dessus. Les fibres de $|U| \to |X|$ sont finies
(cela fait partie
de la propriété qui définit les espaces décents). Ainsi, $U$ est un schéma de
type fini sur $k$ qui possède un nombre fini de points. Ainsi, $U$ est
quasi-fini
sur $k$ (Morphismes, lemme \ref{morphisms-lemma-finite-fibre}),
ce qui signifie par définition que $X \to \Spec(k)$ est localement quasi-fini.
\end{proof}

\begin{lemma}
\label{lemma-conditions-on-point-in-fibre-and-qf}
Soit $S$ un schéma. Soit $f : X \to Y$ un morphisme d'espaces algébriques
sur $S$, localement de type fini. Soit $x \in |X|$, d'image
$y \in |Y|$. Munissons $F = f^{-1}(\{y\})$ de la topologie induite par $|X|$.
Soit $k$ un corps et soit $\Spec(k) \to Y$ un morphisme appartenant à la
classe d'équivalence qui définit $y$. Posons $X_k = \Spec(k) \times_Y X$.
Soit $\tilde x \in |X_k|$ un point qui s'envoie sur $x \in |X|$.
Considérons les conditions suivantes :
\begin{enumerate}
\item
\label{item-fibre-at-x-dim-0}
$\dim_x(F) = 0$,
\item
\label{item-isolated-in-fibre}
$x$ est isolé dans $F$,
\item
\label{item-no-specializations-in-fibre}
$x$ est fermé dans $F$ et, si $x' \leadsto x$ dans $F$, alors $x = x'$,
\item
\label{item-dimension-top-k-fibre}
$\dim_{\tilde x}(|X_k|) = 0$,
\item
\label{item-isolated-in-k-fibre}
$\tilde x$ est isolé dans $|X_k|$,
\item
\label{item-no-specializations-in-k-fibre}
$\tilde x$ est fermé dans $|X_k|$ et, si $\tilde x' \leadsto \tilde x$
dans $|X_k|$, alors $\tilde x = \tilde x'$,
\item
\label{item-k-fibre-at-x-dim-0}
$\dim_{\tilde x}(X_k) = 0$,
\item
\label{item-quasi-finite-at-x}
$f$ est quasi-fini en $x$.
\end{enumerate}
On a alors
$$
\xymatrix{
(\ref{item-dimension-top-k-fibre}) \ar@{=>}[r]_{f\text{ décent}} &
(\ref{item-isolated-in-k-fibre}) \ar@{<=>}[r] &
(\ref{item-no-specializations-in-k-fibre}) \ar@{<=>}[r] &
(\ref{item-k-fibre-at-x-dim-0}) \ar@{<=>}[r] &
(\ref{item-quasi-finite-at-x})
}
$$
Si $Y$ est décent, les conditions (\ref{item-isolated-in-fibre}) et
(\ref{item-no-specializations-in-fibre}) sont équivalentes entre elles ainsi
qu'aux conditions
(\ref{item-isolated-in-k-fibre}),
(\ref{item-no-specializations-in-k-fibre}),
(\ref{item-k-fibre-at-x-dim-0}) et
(\ref{item-quasi-finite-at-x}).
Si $Y$ et $X$ sont décents, toutes les conditions sont équivalentes.
\end{lemma}

\begin{proof}
D'après le Lemme \ref{lemma-conditions-on-point-on-space-over-field}, les
conditions (\ref{item-isolated-in-k-fibre}),
(\ref{item-no-specializations-in-k-fibre}) et
(\ref{item-k-fibre-at-x-dim-0}) sont équivalentes entre elles, ainsi qu'à la
condition que $X_k \to \Spec(k)$ soit quasi-fini en $\tilde x$.
Par conséquent, d'après Morphismes d'espaces, lemme
\ref{spaces-morphisms-lemma-base-change-quasi-finite-locus},
elles sont aussi équivalentes à (\ref{item-quasi-finite-at-x}).
Si $f$ est décent, alors $X_k$ est un espace algébrique décent, et le
Lemme \ref{lemma-conditions-on-point-on-space-over-field}
montre que (\ref{item-dimension-top-k-fibre}) implique
(\ref{item-isolated-in-k-fibre}).

\medskip\noindent
Si $Y$ est décent, on peut choisir un monomorphisme quasi-compact
$\Spec(k') \to Y$ appartenant à la classe d'équivalence de $y$. Dans ce cas,
le Lemme \ref{lemma-topology-fibre}
nous dit que $|X_{k'}| \to F$ est un homéomorphisme.
Combiné aux arguments précédents, cela entraîne les autres assertions du
lemme ; les détails sont omis.
\end{proof}

\begin{lemma}
\label{lemma-conditions-on-fibre-and-qf}
Soit $S$ un schéma. Soit $f : X \to Y$ un morphisme d'espaces algébriques
sur $S$, localement de type fini. Soit $y \in |Y|$. Soit $k$ un corps et soit
$\Spec(k) \to Y$ un morphisme appartenant à la classe d'équivalence qui
définit $y$. Posons $X_k = \Spec(k) \times_Y X$ et munissons
$F = f^{-1}(\{y\})$ de la topologie induite par $|X|$. Considérons les
conditions suivantes :
\begin{enumerate}
\item
\label{item-fibre-finite}
$F$ est fini,
\item
\label{item-fibre-discrete}
$F$ est un espace topologique discret,
\item
\label{item-fibre-no-specializations}
$\dim(F) = 0$,
\item
\label{item-k-fibre-finite}
$|X_k|$ est un ensemble fini,
\item
\label{item-k-fibre-discrete}
$|X_k|$ est un espace discret,
\item
\label{item-k-fibre-no-specializations}
$\dim(|X_k|) = 0$,
\item
\label{item-k-fibre-dim-0}
$\dim(X_k) = 0$,
\item
\label{item-quasi-finite-at-points-fibre}
$f$ est quasi-fini en tout point de $|X|$ situé au-dessus de $y$.
\end{enumerate}
On a alors
$$
\xymatrix{
(\ref{item-fibre-finite}) &
(\ref{item-k-fibre-finite}) \ar@{=>}[l] \ar@{=>}[r]_{f\text{ décent}} &
(\ref{item-k-fibre-discrete}) \ar@{<=>}[r] &
(\ref{item-k-fibre-no-specializations}) \ar@{<=>}[r] &
(\ref{item-k-fibre-dim-0}) \ar@{<=>}[r] &
(\ref{item-quasi-finite-at-points-fibre})
}
$$
Si $Y$ est décent, les conditions (\ref{item-fibre-discrete}) et
(\ref{item-fibre-no-specializations}) sont équivalentes entre elles ainsi
qu'aux conditions (\ref{item-k-fibre-discrete}),
(\ref{item-k-fibre-no-specializations}), (\ref{item-k-fibre-dim-0}) et
(\ref{item-quasi-finite-at-points-fibre}).
Si $Y$ et $X$ sont décents, (\ref{item-fibre-finite}) implique toutes les
autres conditions.
\end{lemma}

\begin{proof}
D'après le Lemme \ref{lemma-conditions-on-space-over-field}, les conditions
(\ref{item-k-fibre-discrete}),
(\ref{item-k-fibre-no-specializations}) et (\ref{item-k-fibre-dim-0})
sont équivalentes entre elles, ainsi qu'à la condition que
$X_k \to \Spec(k)$ soit localement quasi-fini.
Par conséquent, d'après Morphismes d'espaces, lemme
\ref{spaces-morphisms-lemma-base-change-quasi-finite-locus},
elles sont aussi équivalentes à (\ref{item-quasi-finite-at-points-fibre}).
Si $f$ est décent, alors $X_k$ est un espace algébrique décent, et le
Lemme \ref{lemma-conditions-on-space-over-field}
montre que (\ref{item-k-fibre-finite}) implique
(\ref{item-k-fibre-discrete}).

\medskip\noindent
L'application $|X_k| \to F$ est surjective d'après
Propriétés des espaces, lemme
\ref{spaces-properties-lemma-points-cartesian},
et l'on obtient
(\ref{item-k-fibre-finite}) $\Rightarrow$ (\ref{item-fibre-finite}).

\medskip\noindent
Si $Y$ est décent, on peut choisir un monomorphisme quasi-compact
$\Spec(k') \to Y$ appartenant à la classe d'équivalence de $y$. Dans ce cas,
le Lemme \ref{lemma-topology-fibre}
nous dit que $|X_{k'}| \to F$ est un homéomorphisme.
Combiné aux arguments précédents, cela entraîne les autres assertions du
lemme ; les détails sont omis.
\end{proof}







\section{Monomorphismes}
\label{section-monomorphisms}

\noindent
Voici un autre cas où les monomorphismes sont représentables.
On trouvera davantage d'informations dans Morphismes d'espaces, compléments,
section \ref{spaces-more-morphisms-section-monomorphisms}.

\begin{lemma}
\label{lemma-monomorphism-toward-disjoint-union-dim-0-rings}
Soit $S$ un schéma. Soit $Y$ une réunion disjointe de spectres d'anneaux
locaux de dimension zéro sur $S$.
Soit $f : X \to Y$ un monomorphisme d'espaces algébriques sur $S$.
Alors $f$ est représentable, c'est-à-dire que $X$ est un schéma.
\end{lemma}

\begin{proof}
On se ramène immédiatement au cas où $Y = \Spec(A)$ et où
$A$ est un anneau local de dimension zéro, c'est-à-dire où
$\Spec(A) = \{\mathfrak m_A\}$
est réduit à un seul point. Si $X = \emptyset$, il n'y a rien à démontrer.
Sinon, choisissons un schéma affine non vide $U = \Spec(B)$ et un morphisme
étale $U \to X$. Comme $|X|$ est réduit à un seul point (en tant que partie de
$|Y|$ ; voir Morphismes d'espaces, lemme
\ref{spaces-morphisms-lemma-monomorphism-injective-points}),
on voit que $U \to X$ est surjectif. Remarquons que
$U \times_X U = U \times_Y U = \Spec(B \otimes_A B)$.
Les homomorphismes d'anneaux $B \to B \otimes_A B$ sont donc étales.
Puisque
$$
(B \otimes_A B)/\mathfrak m_A(B \otimes_A B)
=
(B/\mathfrak m_AB) \otimes_{A/\mathfrak m_A} (B/\mathfrak m_AB)
$$
on voit que
$B/\mathfrak m_AB \to (B \otimes_A B)/\mathfrak m_A(B \otimes_A B)$
est plat, et même libre de rang égal à la dimension de
$B/\mathfrak m_AB$ comme espace vectoriel sur $A/\mathfrak m_A$. Comme
$B \to B \otimes_A B$ est étale, cela n'est possible que si cette dimension
est finie (voir par exemple
Morphismes, lemmes \ref{morphisms-lemma-etale-universally-bounded} et
\ref{morphisms-lemma-locally-quasi-finite-qc-source-universally-bounded}).
Tout idéal premier de $B$ est situé au-dessus de $\mathfrak m_A$ (l'unique
idéal premier de $A$). Par conséquent,
$\Spec(B) = \Spec(B/\mathfrak m_A B)$ comme espace topologique, et cet espace
est un ensemble discret fini, puisque $B/\mathfrak m_A B$ est un anneau
artinien ; voir
Algèbre, lemmes \ref{algebra-lemma-finite-dimensional-algebra} et
\ref{algebra-lemma-artinian-finite-length}.
Tous les idéaux premiers de $B$ sont donc maximaux, et
$B = B_1 \times \ldots \times B_n$ est un produit d'un nombre fini d'anneaux
locaux de dimension zéro ; voir
Algèbre, lemme \ref{algebra-lemma-product-local}.
Ainsi, $B \to B \otimes_A B$ est fini étale, car tous les anneaux locaux
$B_i$ sont henséliens d'après
Algèbre, lemme \ref{algebra-lemma-local-dimension-zero-henselian}.
Enfin, $X$ est un schéma affine d'après
Groupoïdes, proposition \ref{groupoids-proposition-finite-flat-equivalence}.
\end{proof}








\section{Points génériques}
\label{section-generic-points}

\noindent
Cette section prolonge
Propriétés des espaces, section
\ref{spaces-properties-section-generic-points}.

\begin{lemma}
\label{lemma-decent-generic-points}
Soit $S$ un schéma. Soit $X$ un espace algébrique décent sur $S$.
Soit $x \in |X|$. Les conditions suivantes sont équivalentes :
\begin{enumerate}
\item $x$ est un point générique d'une composante irréductible de $|X|$,
\item pour tout morphisme étale $(Y, y) \to (X, x)$ d'espaces algébriques
pointés, $y$ est un point générique d'une composante irréductible de $|Y|$,
\item pour un certain morphisme étale $(Y, y) \to (X, x)$ d'espaces
algébriques pointés, $y$ est un point générique d'une composante irréductible
de $|Y|$,
\item l'anneau local de $X$ en $x$ est de dimension zéro, et
\item $x$ est un point de codimension $0$ sur $X$.
\end{enumerate}
\end{lemma}

\begin{proof}
Les conditions (4) et (5) sont équivalentes par définition pour tout espace
algébrique ; voir Propriétés des espaces, définition
\ref{spaces-properties-definition-dimension-local-ring}.
Remarquons que tout $Y$ comme en (2) et (3) est décent d'après le
Lemme \ref{lemma-etale-named-properties}.
Il suffit donc de démontrer l'équivalence de (1) et (4) : l'équivalence avec
(2) et (3) en découle, puisque la dimension de l'anneau local de $Y$ en $y$
est égale à celle de l'anneau local de $X$ en $x$.
Soit $f : U \to X$ un morphisme étale provenant d'un schéma affine et soit
$u \in U$ un point qui s'envoie sur $x$.

\medskip\noindent
Supposons (1). Soit $u' \leadsto u$ une spécialisation dans $U$.
Alors $f(u') = f(u) = x$. D'après le
Lemme \ref{lemma-decent-no-specializations-map-to-same-point},
on a $u' = u$. Ainsi, $u$ est un point générique d'une composante irréductible
de $U$. Par conséquent, $\dim(\mathcal{O}_{U, u}) = 0$, et (4) est satisfaite.

\medskip\noindent
Supposons (4). Le point $x$ appartient à une composante irréductible
$T \subset |X|$. Puisque $|X|$ est sobre
(Proposition \ref{proposition-reasonable-sober}),
$T$ possède un point générique $x'$. Bien sûr, $x' \leadsto x$.
On peut relever cette spécialisation en une spécialisation
$u' \leadsto u$ dans $U$
(Lemme \ref{lemma-decent-specialization}). Cela contredit l'hypothèse
$\dim(\mathcal{O}_{U, u}) = 0$, sauf si $u' = u$, c'est-à-dire si $x' = x$.
\end{proof}

\begin{lemma}
\label{lemma-codimension-local-ring}
Soit $S$ un schéma. Soit $X$ un espace algébrique décent sur $S$.
Soit $T \subset |X|$ une partie fermée irréductible. Soit $\xi \in T$
son point générique (Proposition \ref{proposition-reasonable-sober}).
Alors $\text{codim}(T, |X|)$
(Topologie, définition \ref{topology-definition-codimension})
est la dimension de l'anneau local de $X$ en $\xi$
(Propriétés des espaces, définition
\ref{spaces-properties-definition-dimension-local-ring}).
\end{lemma}

\begin{proof}
Choisissons un schéma $U$, un point $u \in U$ et un morphisme étale
$U \to X$ qui envoie $u$ sur $\xi$. Alors toute suite de spécialisations non
triviales $\xi_e \leadsto \ldots \leadsto \xi_0 = \xi$ peut se relever en une
suite $u_e \leadsto \ldots \leadsto u_0 = u$ dans $U$ d'après le
Lemme \ref{lemma-decent-specialization}.
Réciproquement, toute suite de spécialisations non triviales
$u_e \leadsto \ldots \leadsto u_0 = u$ dans $U$ s'envoie sur une suite de
spécialisations non triviales
$\xi_e \leadsto \ldots \leadsto \xi_0 = \xi$ d'après le
Lemme \ref{lemma-decent-no-specializations-map-to-same-point}.
Comme $|X|$ et $U$ sont des espaces topologiques sobres, on conclut que la
codimension de $T$ dans $|X|$ est égale à celle de $\overline{\{u\}}$ dans
$U$. Le lemme se ramène ainsi au cas des schémas, qui est
Propriétés, lemme \ref{properties-lemma-codimension-local-ring}.
\end{proof}

\begin{lemma}
\label{lemma-get-reasonable}
Soit $S$ un schéma. Soit $X$ un espace algébrique sur $S$. Supposons que
\begin{enumerate}
\item tout schéma quasi-compact et étale sur $X$ possède un nombre fini de
composantes irréductibles, et
\item tout $x \in |X|$ de codimension $0$ sur $X$ puisse être représenté
par un monomorphisme $\Spec(k) \to X$.
\end{enumerate}
Alors $X$ est un espace algébrique raisonnable.
\end{lemma}

\begin{proof}
Soit $U$ un schéma affine et soit $a : U \to X$ un morphisme étale.
Il faut montrer que les fibres de $a$ sont universellement bornées. D'après
l'hypothèse (1), le schéma $U$ possède un nombre fini de composantes
irréductibles. Soient $u_1, \ldots, u_n \in U$ les points génériques de ces
composantes irréductibles. Soit $\{x_1, \ldots, x_m\} \subset |X|$ l'image de
$\{u_1, \ldots, u_n\}$. Chaque $x_j$ est un point de codimension $0$.
D'après l'hypothèse (2), on peut choisir un monomorphisme
$\Spec(k_j) \to X$ qui représente $x_j$. D'après Propriétés des espaces, lemme
\ref{spaces-properties-lemma-codimension-0-points}, on a
$$
U \times_X \Spec(k_j) = \coprod\nolimits_{a(u_i) = x_j} \Spec(\kappa(u_i))
$$
C'est un schéma fini sur $\Spec(k_j)$, de degré
$d_j = \sum_{a(u_i) = x_j} [\kappa(u_i) : k_j]$. Posons $n = \max d_j$.

\medskip\noindent
Remarquons que $a$ est séparé
(Propriétés des espaces, lemme
\ref{spaces-properties-lemma-separated-cover}).
Considérons la stratification
$$
X = X_0 \supset X_1 \supset X_2 \supset \ldots
$$
associée à $U \to X$ dans le Lemme \ref{lemma-stratify-flat-fp-lqf}.
Le choix de $n$ ci-dessus montre que $X_{n + 1}$ est vide.
En effet, sinon, $a^{-1}(X_{n + 1})$ serait un ouvert non vide de $U$, et
contiendrait donc l'un des $u_i$. Cela signifierait que $X_{n + 1}$ contient
$x_j = a(u_i)$, ce qui est impossible.
On voit donc que les fibres de $U \to X$ sont universellement bornées
(en fait, par l'entier $n$).
\end{proof}

\begin{lemma}
\label{lemma-finitely-many-irreducible-components}
Soit $S$ un schéma. Soit $X$ un espace algébrique sur $S$.
Les conditions suivantes sont équivalentes :
\begin{enumerate}
\item $X$ est décent et $|X|$ possède un nombre fini de composantes
irréductibles,
\item tout schéma quasi-compact et étale sur $X$ possède un nombre fini de
composantes irréductibles, il existe un nombre fini de $x \in |X|$ de
codimension $0$ sur $X$, et chacun d'eux peut être représenté par un
monomorphisme $\Spec(k) \to X$,
\item il existe un ouvert dense $X' \subset X$ qui est un schéma, tel que
$X'$ possède un nombre fini de composantes irréductibles de points génériques
$\{x'_1, \ldots, x'_m\}$ et que le morphisme $x'_j \to X$ soit quasi-compact
pour $j = 1, \ldots, m$.
\end{enumerate}
De plus, si ces conditions sont satisfaites, alors $X$ est raisonnable et les
points $x'_j \in |X|$ sont les points génériques des composantes irréductibles
de $|X|$.
\end{lemma}

\begin{proof}
Dans la preuve, nous utiliserons sans autre mention Propriétés des espaces,
lemme \ref{spaces-properties-lemma-codimension-0-points}.
Supposons (1). D'après le
Théorème \ref{theorem-decent-open-dense-scheme},
$X$ possède un sous-schéma ouvert dense $X'$.
Comme l'adhérence d'une composante irréductible de $|X'|$ est une composante
irréductible de $|X|$, on voit que $|X'|$ possède un nombre fini de
composantes irréductibles. Ainsi, (3) est satisfaite.

\medskip\noindent
Supposons que $X' \subset X$ soit comme en (3). Soit
$\{x'_1, \ldots, x'_m\}$ l'ensemble des points génériques des composantes
irréductibles de $X'$. Soit $a : U \to X$ un morphisme étale, où $U$ est un
schéma quasi-compact. Pour démontrer (2), il suffit de montrer que $U$ possède
un nombre fini de composantes irréductibles dont les points génériques sont
situés au-dessus de $\{x'_1, \ldots, x'_m\}$. Il suffit de le démontrer pour
les membres d'un recouvrement ouvert affine fini de $U$ ; on peut donc
supposer, et l'on suppose, que $U$ est affine.
Remarquons que $U' = a^{-1}(X') \subset U$ est un ouvert dense.
Puisque $U' \to X'$ est un morphisme étale de schémas, les points génériques
des composantes irréductibles de $U'$ sont les points situés au-dessus de
$\{x'_1, \ldots, x'_m\}$. Comme $x'_j \to X$ est quasi-compact, seul un nombre
fini de points de $U$ sont situés au-dessus de $x'_j$
(Lemme \ref{lemma-UR-finite-above-x}). Ainsi, $U'$ possède un nombre fini de
composantes irréductibles, ce qui implique que les adhérences de ces
composantes irréductibles sont les composantes irréductibles de $U$.
Ainsi, (2) est satisfaite.

\medskip\noindent
Supposons (2). Le Lemme \ref{lemma-get-reasonable} implique alors (1) ainsi
que l'assertion finale. (On utilise aussi qu'un espace algébrique raisonnable
est décent ; voir la discussion qui suit la
Définition \ref{definition-very-reasonable}.)
\end{proof}



\section{Morphismes génériquement finis}
\label{section-generically-finite}

\noindent
Cette section traite, pour les morphismes d'espaces algébriques, les résultats
exposés dans Morphismes, section
\ref{morphisms-section-generically-finite}, et dans
Variétés, section \ref{varieties-section-generically-finite},
pour les morphismes de schémas.

\begin{lemma}
\label{lemma-generically-finite}
Soit $S$ un schéma. Soit $f : X \to Y$ un morphisme d'espaces algébriques
sur $S$. Supposons que $f$ soit quasi-séparé et de type fini.
Soit $y \in |Y|$ un point de codimension $0$ sur $Y$.
Les conditions suivantes sont équivalentes :
\begin{enumerate}
\item l'espace $|X_k|$ est fini, où $\Spec(k) \to Y$ représente $y$,
\item $X \to Y$ est quasi-fini en tout point de $|X|$ situé au-dessus de $y$,
\item il existe un sous-espace ouvert $Y' \subset Y$, avec $y \in |Y'|$,
tel que $Y' \times_Y X \to Y'$ soit fini.
\end{enumerate}
Si $Y$ est décent, ces conditions sont aussi équivalentes à
\begin{enumerate}
\item[(4)] l'ensemble $f^{-1}(\{y\})$ est fini.
\end{enumerate}
\end{lemma}

\begin{proof}
L'équivalence de (1) et (2) résulte du
Lemme \ref{lemma-conditions-on-fibre-and-qf}
(et du fait qu'un morphisme quasi-séparé est décent, d'après le
Lemme \ref{lemma-properties-trivial-implications}).

\medskip\noindent
Supposons satisfaites les conditions équivalentes (1) et (2). Choisissons un
schéma affine $V$ et un morphisme étale $V \to Y$ qui envoie un point
$v \in V$ sur $y$. Alors $v$ est un point générique d'une composante
irréductible de $V$, d'après Propriétés des espaces, lemme
\ref{spaces-properties-lemma-codimension-0-points}.
Choisissons un schéma affine $U$ et un morphisme étale surjectif
$U \to V \times_Y X$. Alors $U \to V$ est de type fini.
Le morphisme $U \to V$ est quasi-fini en tout point situé au-dessus de $v$,
d'après (2). Il s'ensuit que la fibre de $U \to V$ au-dessus de $v$ est finie
(Morphismes, lemme
\ref{morphisms-lemma-quasi-finite-at-a-finite-number-of-points}). D'après
Morphismes, lemme \ref{morphisms-lemma-generically-finite},
après avoir rétréci $V$, on peut supposer que $U \to V$ est fini.
Posons
$$
R = U \times_{V \times_Y X} U.
$$
Puisque $f$ est quasi-séparé, $V \times_Y X$ est quasi-séparé, et $R$ est
donc un schéma quasi-compact. De plus, le morphisme $R \to V$ est quasi-fini,
comme composé du morphisme étale $R \to U$ et du morphisme fini $U \to V$.
On peut donc appliquer de nouveau Morphismes, lemme
\ref{morphisms-lemma-generically-finite} et, après avoir encore rétréci $V$,
supposer que $R \to V$ est fini. Cela implique bien sûr que les deux
projections $R \to U$ sont finies étales. Il s'ensuit que
$U/R = V \times_Y X$ est un schéma affine ; voir Groupoïdes, proposition
\ref{groupoids-proposition-finite-flat-equivalence}.
D'après Morphismes, lemme \ref{morphisms-lemma-image-proper-is-proper},
on conclut que $V \times_Y X \to V$ est propre, puis, d'après
Morphismes, lemme \ref{morphisms-lemma-finite-proper},
que $V \times_Y X \to V$ est fini. Enfin, soit $Y' \subset Y$ le
sous-espace ouvert de $Y$
correspondant à l'image de $|V| \to |Y|$.
D'après Morphismes d'espaces, lemme
\ref{spaces-morphisms-lemma-integral-local}, on conclut que
$Y' \times_Y X \to Y'$ est fini, puisque son changement de base à $V$ est
fini et que $V \to Y'$ est un morphisme étale surjectif.

\medskip\noindent
Si $Y$ est décent et $f$ quasi-séparé, alors $X$ est également décent ;
utiliser les Lemmes \ref{lemma-properties-trivial-implications} et
\ref{lemma-property-over-property}.
Le Lemme \ref{lemma-conditions-on-fibre-and-qf} montre donc que (4) implique
(1) et (2). Réciproquement, (2) implique (4) d'après Morphismes d'espaces,
lemme \ref{spaces-morphisms-lemma-quasi-finite-at-a-finite-number-of-points}.
\end{proof}

\begin{lemma}
\label{lemma-generically-finite-reprise}
Soit $S$ un schéma. Soit $f : X \to Y$ un morphisme d'espaces algébriques
sur $S$. Supposons que $f$ soit quasi-séparé et localement de type fini,
et que $Y$ soit quasi-séparé. Soit $y \in |Y|$ un point de codimension $0$
sur $Y$. Les conditions suivantes sont équivalentes :
\begin{enumerate}
\item l'ensemble $f^{-1}(\{y\})$ est fini,
\item l'espace $|X_k|$ est fini, où $\Spec(k) \to Y$ représente $y$,
\item il existe des sous-espaces ouverts $X' \subset X$ et $Y' \subset Y$
tels que $f(X') \subset Y'$, $y \in |Y'|$ et
$f^{-1}(\{y\}) \subset |X'|$, et tels que
$f|_{X'} : X' \to Y'$ soit fini.
\end{enumerate}
\end{lemma}

\begin{proof}
Les espaces algébriques quasi-séparés étant décents, l'équivalence de (1)
et (2) résulte du Lemme \ref{lemma-conditions-on-fibre-and-qf}.
Pour montrer que (1) et (2) impliquent (3), on peut, et l'on va, remplacer
$Y$ par un ouvert quasi-compact contenant $y$.
Comme $f^{-1}(\{y\})$ est fini, on peut trouver un sous-espace ouvert
quasi-compact $X' \subset X$ contenant la fibre.
La restriction $f|_{X'} : X' \to Y$ est quasi-compacte et quasi-séparée
d'après Morphismes d'espaces, lemme
\ref{spaces-morphisms-lemma-quasi-compact-quasi-separated-permanence}
(c'est ici qu'on utilise que $Y$ est quasi-séparé).
En appliquant le Lemme \ref{lemma-generically-finite} à
$f|_{X'} : X' \to Y$, on obtient (3).
Nous omettons la démonstration de l'implication de (3) vers (2).
\end{proof}

\begin{lemma}
\label{lemma-quasi-finiteness-over-generic-point}
Soit $S$ un schéma. Soit $f : X \to Y$ un morphisme d'espaces algébriques
sur $S$. Supposons que $f$ soit localement de type fini.
Notons $X^0 \subset |X|$, respectivement $Y^0 \subset |Y|$, l'ensemble des
points de codimension $0$ de $X$, respectivement de $Y$. Soit $y \in Y^0$.
Les conditions suivantes sont équivalentes :
\begin{enumerate}
\item $f^{-1}(\{y\}) \subset X^0$,
\item $f$ est quasi-fini en tout point situé au-dessus de $y$,
\item $f$ est quasi-fini en tout $x \in X^0$ situé au-dessus de $y$.
\end{enumerate}
\end{lemma}

\begin{proof}
Soit $V$ un schéma et soit $V \to Y$ un morphisme étale surjectif.
Soit $U$ un schéma et soit $U \to V \times_Y X$ un morphisme étale
surjectif. Alors $f$ est quasi-fini en l'image $x$ d'un point $u \in U$
si et seulement si $U \to V$ est quasi-fini en $u$. De plus, $x \in X^0$
si et seulement si $u$ est le point générique d'une composante irréductible
de $U$ (Propriétés des espaces, lemme
\ref{spaces-properties-lemma-codimension-0-points}).
Le lemme se ramène donc au cas du morphisme $U \to V$, c'est-à-dire à
Morphismes, lemme
\ref{morphisms-lemma-quasi-finiteness-over-generic-point}.
\end{proof}

\begin{lemma}
\label{lemma-finite-over-dense-open}
Soit $S$ un schéma. Soit $f : X \to Y$ un morphisme d'espaces algébriques
sur $S$. Supposons que $f$ soit localement de type fini.
Notons $X^0 \subset |X|$, respectivement $Y^0 \subset |Y|$, l'ensemble des
points de codimension $0$ de $X$, respectivement de $Y$. Supposons que
\begin{enumerate}
\item $Y$ soit décent,
\item $X^0$ et $Y^0$ soient finis et $f^{-1}(Y^0) = X^0$,
\item soit $f$ soit quasi-compact, soit $f$ soit séparé.
\end{enumerate}
Alors il existe un ouvert dense $V \subset Y$ tel que
$f^{-1}(V) \to V$ soit fini.
\end{lemma}

\begin{proof}
D'après les Lemmes \ref{lemma-finitely-many-irreducible-components} et
\ref{lemma-decent-generic-points}, on peut supposer que $Y$ est un schéma
ayant un nombre fini de composantes irréductibles. En rétrécissant encore,
on peut supposer que $Y$ est un schéma affine irréductible de point générique
$y$. La fibre de $f$ au-dessus de $y$ est alors finie.

\medskip\noindent
Supposons que $f$ soit quasi-compact et $Y$ affine irréductible. Alors $X$
est quasi-compact, et l'on peut choisir un schéma affine $U$ et un morphisme
étale surjectif $U \to X$. Le morphisme $U \to Y$ est alors de type fini,
et la fibre de $U \to Y$ au-dessus de $y$ est l'ensemble $U^0$ des points
génériques des composantes irréductibles de $U$ (Propriétés des espaces,
lemme \ref{spaces-properties-lemma-codimension-0-points}).
Ainsi, $U^0$ est fini
(Morphismes, lemme
\ref{morphisms-lemma-quasi-finite-at-a-finite-number-of-points}),
et, après avoir rétréci $Y$, on peut supposer que $U \to Y$ est fini
(Morphismes, lemme \ref{morphisms-lemma-generically-finite}).
Considérons ensuite $R = U \times_X U$. Comme la projection
$s : R \to U$ est étale, on a $R^0 = s^{-1}(U^0)$, qui est situé au-dessus
de $y$. Puisque $R \to U \times_Y U$ est un monomorphisme, on conclut que
$R^0$ est fini, car $U \times_Y U \to Y$ est fini.
De plus, $R$ est séparé
(Propriétés des espaces, lemme \ref{spaces-properties-lemma-separated-cover}).
On peut donc rétrécir encore une fois $Y$ pour se ramener au cas où $R$ est
fini sur $Y$
(Morphismes, lemme \ref{morphisms-lemma-finite-over-dense-open}).
Il s'ensuit que $X = U/R$ est fini sur $Y$, par les mêmes arguments que dans
la démonstration du Lemme \ref{lemma-generically-finite}
(ou bien on peut simplement appliquer ce lemme, puisqu'il en résulte
immédiatement que $X$ est également quasi-séparé).

\medskip\noindent
Supposons que $f$ soit séparé et $Y$ affine irréductible. Choisissons
$V \subset Y$ et $U \subset X$ comme dans le
Lemme \ref{lemma-generically-finite-reprise}.
Puisque $f|_U : U \to V$ est fini, $U \subset f^{-1}(V)$ est à la fois fermé
et ouvert
(Morphismes d'espaces, lemmes
\ref{spaces-morphisms-lemma-universally-closed-permanence} et
\ref{spaces-morphisms-lemma-finite-proper}).
Ainsi, $f^{-1}(V) = U \amalg W$ pour un certain sous-espace ouvert $W$ de
$X$. Mais, comme $U$ contient tous les points de codimension $0$ de $X$,
on conclut que $W = \emptyset$
(Propriétés des espaces, lemme
\ref{spaces-properties-lemma-codimension-0-points-dense}),
comme souhaité.
\end{proof}






\section{Morphismes birationnels}
\label{section-birational}

\noindent
La définition suivante d'un morphisme birationnel d'espaces algébriques
semble être celle qui se rapproche le plus de notre définition
(Morphismes, définition \ref{morphisms-definition-birational})
d'un morphisme birationnel de schémas.

\begin{definition}
\label{definition-birational}
Soit $S$ un schéma. Soient $X$ et $Y$ des espaces algébriques sur $S$.
Supposons que $X$ et $Y$ soient décents et que $|X|$ et $|Y|$ possèdent
un nombre fini de composantes irréductibles. On dit qu'un morphisme
$f : X \to Y$ est {\it birationnel} si
\begin{enumerate}
\item $|f|$ induit une bijection entre l'ensemble des points génériques
des composantes irréductibles de $|X|$ et l'ensemble des points génériques
des composantes irréductibles de $|Y|$, et
\item pour tout point générique $x \in |X|$ d'une composante irréductible,
l'homomorphisme d'anneaux locaux
$\mathcal{O}_{Y, f(x)} \to \mathcal{O}_{X, x}$ est un isomorphisme
(voir la précision ci-dessous).
\end{enumerate}
\end{definition}

\noindent
Précision : puisque $X$ et $Y$ sont décents, les espaces topologiques
$|X|$ et $|Y|$ sont sobres
(Proposition \ref{proposition-reasonable-sober}).
La condition (1) a donc bien un sens. De plus, puisque nous avons supposé que
$|X|$ et $|Y|$ ont un nombre fini de composantes irréductibles, les points
génériques $x_1, \ldots, x_n \in |X|$, respectivement
$y_1, \ldots, y_n \in |Y|$, appartiennent à tout ouvert dense de $|X|$,
respectivement de $|Y|$. En particulier, ils appartiennent au lieu schématique
de $X$, respectivement de $Y$, d'après le
Théorème \ref{theorem-decent-open-dense-scheme}.
On peut donc définir $\mathcal{O}_{X, x_i}$, respectivement
$\mathcal{O}_{Y, y_i}$, comme l'anneau local de ce schéma en $x_i$,
respectivement en $y_i$.

\medskip\noindent
On en conclut que, si le morphisme $f : X \to Y$ est birationnel, il existe
des sous-espaces ouverts denses $X' \subset X$ et $Y' \subset Y$ tels que
\begin{enumerate}
\item $f(X') \subset Y'$,
\item $X'$ et $Y'$ soient représentables, et
\item $f|_{X'} : X' \to Y'$ soit birationnel au sens de Morphismes,
définition \ref{morphisms-definition-birational}.
\end{enumerate}
Nous exigeons toutefois que $X$ et $Y$ soient décents et possèdent un nombre
fini de composantes irréductibles. Le
Lemme \ref{lemma-finitely-many-irreducible-components} donne d'autres
caractérisations des espaces algébriques décents ayant un nombre fini de
composantes irréductibles. Dans la plupart des cas, les morphismes
birationnels sont des isomorphismes au-dessus d'ouverts denses.

\begin{lemma}
\label{lemma-birational-dominant}
Soit $S$ un schéma. Soit $f : X \to Y$ un morphisme d'espaces algébriques
sur $S$, supposés décents et ayant un nombre fini de composantes
irréductibles. Si $f$ est birationnel, alors $f$ est dominant.
\end{lemma}

\begin{proof}
Cela résulte immédiatement des définitions. Voir Morphismes d'espaces,
définition \ref{spaces-morphisms-definition-dominant}.
\end{proof}

\begin{lemma}
\label{lemma-birational-generic-fibres}
Soit $S$ un schéma. Soit $f : X \to Y$ un morphisme birationnel d'espaces
algébriques sur $S$, supposés décents et ayant un nombre fini de composantes
irréductibles. Si $y \in |Y|$ est le point générique d'une composante
irréductible, alors le changement de base
$X \times_Y \Spec(\mathcal{O}_{Y, y}) \to
\Spec(\mathcal{O}_{Y, y})$ est un isomorphisme.
\end{lemma}

\begin{proof}
Soient $X' \subset X$ et $Y' \subset Y$ les sous-espaces ouverts maximaux
qui sont représentables ; voir le
Lemme \ref{lemma-finitely-many-irreducible-components}.
D'après le Lemme \ref{lemma-quasi-finiteness-over-generic-point},
la fibre de $f$ au-dessus de $y$ est constituée de points de codimension $0$
de $X$ et est donc contenue dans $X'$. Ainsi,
$X \times_Y \Spec(\mathcal{O}_{Y, y}) =
X' \times_{Y'} \Spec(\mathcal{O}_{Y', y})$, et le résultat découle de
Morphismes, lemme \ref{morphisms-lemma-birational-generic-fibres}.
\end{proof}

\begin{lemma}
\label{lemma-birational-birational}
Soit $S$ un schéma. Soit $f : X \to Y$ un morphisme birationnel d'espaces
algébriques sur $S$, supposés décents et ayant un nombre fini de composantes
irréductibles. Supposons satisfaite l'une des conditions suivantes :
\begin{enumerate}
\item $f$ est localement de type fini et $Y$ est réduit
(par exemple intègre),
\item $f$ est localement de présentation finie.
\end{enumerate}
Alors il existe des ouverts denses $U \subset X$ et $V \subset Y$ tels que
$f(U) \subset V$ et que $f|_U : U \to V$ soit un isomorphisme.
\end{lemma}

\begin{proof}
D'après le Lemme \ref{lemma-finitely-many-irreducible-components},
on peut supposer que $X$ et $Y$ sont des schémas. Dans ce cas, le résultat est
Morphismes, lemme \ref{morphisms-lemma-birational-birational}.
\end{proof}

\begin{lemma}
\label{lemma-birational-isomorphism-over-dense-open}
Soit $S$ un schéma. Soit $f : X \to Y$ un morphisme birationnel d'espaces
algébriques sur $S$, supposés décents et ayant un nombre fini de composantes
irréductibles. Supposons que
\begin{enumerate}
\item soit $f$ soit quasi-compact, soit $f$ soit séparé, et
\item soit $f$ soit localement de type fini et $Y$ soit réduit, soit $f$
soit localement de présentation finie.
\end{enumerate}
Alors il existe un ouvert dense $V \subset Y$ tel que
$f^{-1}(V) \to V$ soit un isomorphisme.
\end{lemma}

\begin{proof}
D'après le Lemme \ref{lemma-finitely-many-irreducible-components},
on peut supposer que $Y$ est un schéma. D'après le
Lemme \ref{lemma-finite-over-dense-open}, on peut supposer que $f$ est fini.
Alors $X$ est aussi un schéma, et le résultat découle de Morphismes,
lemme \ref{morphisms-lemma-birational-isomorphism-over-dense-open}.
\end{proof}

\begin{lemma}
\label{lemma-birational-etale-localization}
Soit $S$ un schéma. Soit $f : X \to Y$ un morphisme d'espaces algébriques
sur $S$, supposés décents et ayant un nombre fini de composantes
irréductibles. Si $f$ est birationnel et si $V \to Y$ est un morphisme étale
avec $V$ affine, alors $X \times_Y V$ est décent, possède un nombre fini de
composantes irréductibles et $X \times_Y V \to V$ est birationnel.
\end{lemma}

\begin{proof}
L'espace algébrique $U = X \times_Y V$ est décent
(Lemme \ref{lemma-etale-named-properties}).
Les points génériques de $V$ et de $U$ sont les éléments de $|V|$ et de $|U|$
qui sont situés au-dessus des points génériques de $|Y|$ et de $|X|$
(Lemme \ref{lemma-decent-generic-points}).
Comme $Y$ est décent, on en conclut que $V$ ne possède qu'un nombre fini de
points génériques. Soit $\xi \in |X|$ un point générique d'une composante
irréductible. D'après la discussion qui suit la
Définition \ref{definition-birational}, on a un carré cartésien
$$
\xymatrix{
\Spec(\mathcal{O}_{X, \xi}) \ar[d] \ar[r] & X \ar[d] \\
\Spec(\mathcal{O}_{Y, f(\xi)}) \ar[r] & Y
}
$$
dont les morphismes horizontaux sont des monomorphismes qui identifient les
anneaux locaux, et dont la flèche verticale de gauche est un isomorphisme.
Il en résulte que, dans le diagramme
$$
\xymatrix{
\Spec(\mathcal{O}_{X, \xi}) \times_X U \ar[d] \ar[r] & U \ar[d] \\
\Spec(\mathcal{O}_{Y, f(\xi)}) \times_Y V \ar[r] & V
}
$$
la flèche verticale de gauche est un isomorphisme. Les flèches horizontales
ont leur image contenue dans le lieu schématique de $U$ et de $V$ et
identifient les anneaux locaux (nous omettons certains détails). Comme les
images des flèches horizontales sont les points de $|U|$, respectivement de
$|V|$, situés au-dessus de $\xi$, respectivement de $f(\xi)$, on conclut.
\end{proof}

\begin{lemma}
\label{lemma-birational-induced-morphism-normalizations}
Soit $S$ un schéma. Soit $f : X \to Y$ un morphisme birationnel entre des
espaces algébriques sur $S$, supposés décents et ayant un nombre fini de
composantes irréductibles. Alors les normalisations $X^\nu \to X$ et
$Y^\nu \to Y$ existent, et il existe un diagramme commutatif
$$
\xymatrix{
X^\nu \ar[r] \ar[d] &  Y^\nu \ar[d] \\
X \ar[r] & Y
}
$$
d'espaces algébriques sur $S$. Le morphisme $X^\nu \to Y^\nu$ est
birationnel.
\end{lemma}

\begin{proof}
D'après le Lemme \ref{lemma-finitely-many-irreducible-components},
$X$ et $Y$ satisfont aux conditions équivalentes de Morphismes d'espaces,
lemme \ref{spaces-morphisms-lemma-prepare-normalization}, et les
normalisations sont définies. D'après Morphismes d'espaces, lemme
\ref{spaces-morphisms-lemma-normalization-normal}, l'espace algébrique
$X^\nu$ est normal et envoie les points de codimension $0$ sur des points de
codimension $0$. Comme $f$ envoie les points de codimension $0$ sur des
points de codimension $0$ (ils coïncident avec les points génériques sur les
espaces décents d'après le Lemme \ref{lemma-decent-generic-points}),
Morphismes d'espaces, lemme
\ref{spaces-morphisms-lemma-normalization-normal}, fournit une factorisation
du composé $X^\nu \to X \to Y$ par $Y^\nu$.

\medskip\noindent
Remarquons que $X^\nu$ et $Y^\nu$ sont décents, par exemple d'après le
Lemme \ref{lemma-representable-named-properties}.
De plus, les applications $X^\nu \to X$ et $Y^\nu \to Y$ induisent des
bijections sur les composantes irréductibles (voir les références ci-dessus) ;
ainsi, $X^\nu$ et $Y^\nu$ possèdent tous deux un nombre fini de composantes
irréductibles, et l'application $X^\nu \to Y^\nu$ induit une bijection entre
leurs points génériques. Pour montrer que $X^\nu \to Y^\nu$ est birationnel,
il suffit donc de montrer qu'il induit un isomorphisme sur les anneaux locaux
en ces points. Pour cela, on peut remplacer $X$ et $Y$ par des voisinages
ouverts de leurs points génériques ; on peut donc supposer que $X$ et $Y$
sont des schémas affines irréductibles de points génériques $x$ et $y$.
Puisque $f$ est birationnel, l'homomorphisme
$\mathcal{O}_{Y, y} \to \mathcal{O}_{X, x}$ est un isomorphisme.
Soient $x^\nu \in X^\nu$ et $y^\nu \in Y^\nu$ les points situés au-dessus
de $x$ et de $y$. Par construction de la normalisation, on a
$\mathcal{O}_{X^\nu, x^\nu} = \mathcal{O}_{X, x}/\mathfrak m_x$,
et de même sur $Y$. Ainsi, l'homomorphisme
$\mathcal{O}_{Y^\nu, y^\nu} \to \mathcal{O}_{X^\nu, x^\nu}$
est lui aussi un isomorphisme.
\end{proof}

\begin{lemma}
\label{lemma-finite-birational-over-normal}
Soit $S$ un schéma. Soit $f : X \to Y$ un morphisme d'espaces algébriques
sur $S$. Supposons que
\begin{enumerate}
\item $X$ et $Y$ soient décents et aient un nombre fini de composantes
irréductibles,
\item $f$ soit intégral et birationnel,
\item $Y$ soit normal, et
\item $X$ soit réduit.
\end{enumerate}
Alors $f$ est un isomorphisme.
\end{lemma}

\begin{proof}
Soit $V \to Y$ un morphisme étale, avec $V$ affine. Il suffit de montrer que
$U = X \times_Y V \to V$ est un isomorphisme. D'après le
Lemme \ref{lemma-birational-etale-localization} et sa démonstration,
$U$ et $V$ sont décents et possèdent un nombre fini de composantes
irréductibles, et $U \to V$ est birationnel. D'après Propriétés, lemme
\ref{properties-lemma-normal-locally-finite-nr-irreducibles},
$V$ est une réunion disjointe finie de schémas intègres.
On peut donc supposer que $V$ est intègre. Comme $f$ est birationnel, $U$ est
irréductible et réduit, c'est-à-dire intègre
(remarquer que $U$ est un schéma, puisque $f$ est intégral, donc
représentable). On peut ainsi supposer que $X$ et $Y$ sont des schémas
intègres, et le résultat découle du cas des schémas ; voir Morphismes, lemme
\ref{morphisms-lemma-finite-birational-over-normal}.
\end{proof}

\begin{lemma}
\label{lemma-normalization-normal}
Soit $S$ un schéma. Soit $f : X \to Y$ un morphisme intégral birationnel
d'espaces algébriques décents sur $S$, ayant un nombre fini de composantes
irréductibles. Il existe alors une factorisation $Y^\nu \to X \to Y$, et
$Y^\nu \to X$ est la normalisation de $X$.
\end{lemma}

\begin{proof}
Considérons l'application $X^\nu \to Y^\nu$ du
Lemme \ref{lemma-birational-induced-morphism-normalizations}.
Cette application est intégrale d'après Morphismes d'espaces, lemme
\ref{spaces-morphisms-lemma-finite-permanence}.
Elle est donc un isomorphisme d'après le
Lemme \ref{lemma-finite-birational-over-normal}.
\end{proof}




\section{Espaces de Jacobson}
\label{section-jacobson}

\noindent
Nous avons défini la propriété de Jacobson pour les espaces algébriques dans
Propriétés des espaces, remarque
\ref{spaces-properties-remark-list-properties-local-etale-topology}.
Pour les espaces algébriques représentables, elle coïncide avec la propriété
étudiée dans Propriétés, section \ref{properties-section-jacobson}.
Le rapport entre la propriété de Jacobson et le comportement de l'espace
topologique $|X|$ n'est pas évident pour un espace algébrique général $X$.
Cependant, un espace algébrique décent $X$ (par exemple quasi-séparé ou
localement séparé) est de Jacobson si et seulement si $|X|$ est de Jacobson
(voir le Lemme \ref{lemma-decent-Jacobson}).

\begin{lemma}
\label{lemma-Jacobson-universally-Jacobson}
Soit $S$ un schéma. Soit $X$ un espace algébrique de Jacobson sur $S$.
Tout espace algébrique localement de type fini sur $X$ est de Jacobson.
\end{lemma}

\begin{proof}
Soit $U \to X$ un morphisme étale surjectif, où $U$ est un schéma.
Alors $U$ est de Jacobson (par définition), et, pour tout morphisme de schémas
$V \to U$ localement de type fini, $V$ est de Jacobson d'après le résultat
correspondant pour les schémas (Morphismes, lemme
\ref{morphisms-lemma-Jacobson-universally-Jacobson}).
Ainsi, si $Y \to X$ est un morphisme d'espaces algébriques localement de type
fini, en posant $V = U \times_X Y$, on voit par définition que $Y$ est de
Jacobson.
\end{proof}

\begin{lemma}
\label{lemma-Jacobson-ft-points-lift-to-closed}
Soit $S$ un schéma. Soit $X$ un espace algébrique de Jacobson sur $S$.
Soient $x \in X_{\text{ft-pts}}$ et $g : W \to X$ un morphisme localement de
type fini, où $W$ est un schéma. Si $x \in \Im(|g|)$, alors il existe un point
fermé de $W$ qui s'envoie sur $x$.
\end{lemma}

\begin{proof}
Soit $U \to X$ un morphisme étale, où $U$ est un schéma, et soit $u \in U$
un point fermé qui s'envoie sur $x$ ; voir Morphismes d'espaces, lemme
\ref{spaces-morphisms-lemma-identify-finite-type-points}.
Remarquons que $W$, $W \times_X U$ et $U$ sont des schémas de Jacobson d'après
le Lemme \ref{lemma-Jacobson-universally-Jacobson}.
Sur ces schémas, les points de type fini sont donc exactement les points
fermés, d'après Morphismes, lemme
\ref{morphisms-lemma-jacobson-finite-type-points}.
L'image inverse $T \subset W \times_X U$ de $u$ est une partie fermée non vide
(puisque $x$ appartient à l'image de $W \to X$).
D'après Morphismes, lemme \ref{morphisms-lemma-enough-finite-type-points},
il existe un point fermé $t$ de $W \times_X U$ qui s'envoie sur $u$.
Comme $W \times_X U \to W$ est localement de type fini, l'image de $t$ dans
$W$ est fermée d'après Morphismes, lemme
\ref{morphisms-lemma-jacobson-finite-type-points}.
\end{proof}

\begin{lemma}
\label{lemma-decent-Jacobson-ft-pts}
Soit $S$ un schéma. Soit $X$ un espace algébrique décent de Jacobson sur $S$.
Alors $X_{\text{ft-pts}} \subset |X|$ est l'ensemble des points fermés.
\end{lemma}

\begin{proof}
Si $x \in |X|$ est fermé, on peut représenter $x$ par une immersion fermée
$\Spec(k) \to X$ ; voir le Lemme \ref{lemma-decent-space-closed-point}.
Ainsi, $x$ est certainement un point de type fini.

\medskip\noindent
Réciproquement, soit $x \in |X|$ un point de type fini. On sait que $x$ peut
être représenté par un monomorphisme quasi-compact $\Spec(k) \to X$, où $k$
est un corps (Définition \ref{definition-very-reasonable}). D'autre part, par
définition, il existe un morphisme $\Spec(k') \to X$ localement de type fini
qui représente $x$
(Morphismes, définition \ref{morphisms-definition-finite-type-point}).
On obtient une factorisation $\Spec(k') \to \Spec(k) \to X$.
Soit $U \to X$ un morphisme étale quelconque, avec $U$ affine, et considérons
les morphismes
$$
\Spec(k') \times_X U \to \Spec(k) \times_X U \to U.
$$
Le schéma quasi-compact $\Spec(k) \times_X U$ est étale sur $\Spec(k)$ ;
c'est donc une réunion disjointe finie de spectres de corps
(Remarque \ref{remark-recall}). De plus, le premier morphisme est surjectif
et localement de type fini
(Morphismes, lemme \ref{morphisms-lemma-permanence-finite-type}) ;
il est donc surjectif sur les points de type fini
(Morphismes, lemme
\ref{morphisms-lemma-finite-type-points-surjective-morphism}), et le composé
(qui est localement de type fini) envoie les points de type fini sur des
points fermés, puisque $U$ est de Jacobson
(Morphismes, lemme \ref{morphisms-lemma-jacobson-finite-type-points}).
Ainsi, l'image de $\Spec(k) \times_X U \to U$ est un ensemble fini de points
fermés, donc fermé. Comme cela vaut pour tout $U$ affine et tout morphisme
étale $U \to X$, on conclut que $x \in |X|$ est fermé.
\end{proof}

\begin{lemma}
\label{lemma-decent-Jacobson}
Soit $S$ un schéma. Soit $X$ un espace algébrique décent sur $S$.
Alors $X$ est de Jacobson si et seulement si $|X|$ est de Jacobson.
\end{lemma}

\begin{proof}
Supposons que $X$ soit de Jacobson et que $T \subset |X|$ soit une partie
fermée. D'après Morphismes d'espaces, lemme
\ref{spaces-morphisms-lemma-enough-finite-type-points}, on voit que
$T \cap X_{\text{ft-pts}}$ est dense dans $T$.
D'après le Lemme \ref{lemma-decent-Jacobson-ft-pts},
$X_{\text{ft-pts}}$ est l'ensemble des points fermés de $|X|$.
Ainsi, $|X|$ est bien de Jacobson.

\medskip\noindent
Supposons que $|X|$ soit de Jacobson. Soit $f : U \to X$ un morphisme étale,
où $U$ est un schéma affine. Il faut montrer que $U$ est de Jacobson.
Si $x \in |X|$ est fermé, la fibre $F = f^{-1}(\{x\})$ est une partie finie
(par définition de décent) et fermée (par construction de la topologie sur
$|X|$) de $U$. Comme il n'y a pas de spécialisations entre les points de $F$
(Lemme \ref{lemma-decent-no-specializations-map-to-same-point}),
on conclut que tout point de $F$ est fermé dans $U$.
Si $U$ n'est pas de Jacobson, il existe un point non fermé $u \in U$ tel que
$\{u\}$ soit localement fermé (Topologie, lemme
\ref{topology-lemma-non-jacobson-Noetherian-characterize}).
Nous allons montrer que $f(u) \in |X|$ est fermé ; d'après ce qui précède,
$u$ serait alors fermé dans $U$, contradiction qui achèvera la démonstration.
Pour cela, on peut remplacer $U$ par un voisinage ouvert affine de $u$.
On peut donc supposer que $\{u\}$ est fermé dans $U$.
Posons $R = U \times_X U$, avec projections $s, t : R \to U$.
Alors $s^{-1}(\{u\}) = \{r_1, \ldots, r_m\}$ est fini
(par définition des espaces décents). Après avoir remplacé $U$ par un
voisinage ouvert affine plus petit de $u$, on peut supposer que
$t(r_j) = u$ pour $j = 1, \ldots, m$. Il s'ensuit que $\{u\}$ est une partie
fermée $R$-invariante de $U$. Ainsi, $\{f(u)\}$ est une partie localement
fermée de $X$, puisqu'elle est fermée dans l'ouvert $|f|(|U|)$ de $|X|$.
Comme $|X|$ est de Jacobson, on conclut que $f(u)$ est fermé dans $|X|$,
comme souhaité.
\end{proof}

\begin{lemma}
\label{lemma-punctured-spec}
Soit $S$ un schéma. Soit $X$ un espace algébrique décent localement
noethérien sur $S$. Soit $x \in |X|$. Alors
$$
W = \{x' \in |X| : x' \leadsto x,\ x' \not = x\}
$$
est un espace topologique noethérien, spectral, sobre et de Jacobson.
\end{lemma}

\begin{proof}
On peut effectuer le remplacement par tout sous-espace ouvert contenant $x$.
On peut donc supposer que $X$ est quasi-compact.
Alors $|X|$ est un espace topologique noethérien
(Propriétés des espaces, lemme
\ref{spaces-properties-lemma-Noetherian-topology}).
Ainsi, $W$ est un espace topologique noethérien
(Topologie, lemme \ref{topology-lemma-Noetherian}).

\medskip\noindent
En combinant le Lemme \ref{lemma-locally-Noetherian-decent-quasi-separated}
et Propriétés des espaces, lemme
\ref{spaces-properties-lemma-quasi-compact-quasi-separated-spectral},
on voit que $|X|$ est un espace topologique spectral.
D'après Topologie, lemme \ref{topology-lemma-make-spectral-space},
$W \cup \{x\}$ est un espace topologique spectral.
Or $W$ est un ouvert quasi-compact de $W \cup \{x\}$ ; $W$ est donc spectral
d'après Topologie, lemme \ref{topology-lemma-spectral-sub}.

\medskip\noindent
Soit $E \subset W$ une partie fermée irréductible. Si $Z \subset |X|$ est
l'adhérence de $E$, alors $x \in Z$. Il existe un unique point générique
$\eta \in Z$ d'après la Proposition \ref{proposition-reasonable-sober}.
Bien sûr, $\eta \in W$, donc $\eta \in E$. On conclut que $E$ possède un
unique point générique, c'est-à-dire que $W$ est sobre.

\medskip\noindent
Soit $x' \in W$ un point tel que $\{x'\}$ soit localement fermé dans $W$.
Pour achever la démonstration, il faut montrer que $x'$ est un point fermé de
$W$. Sinon, il existe une spécialisation non triviale
$x' \leadsto x'_1$ dans $W$. Soient $U$ un schéma affine, $u \in U$ un point
et $U \to X$ un morphisme étale envoyant $u$ sur $x$. D'après le
Lemme \ref{lemma-decent-specialization}, on peut choisir des spécialisations
$u' \leadsto u'_1 \leadsto u$ qui s'envoient sur
$x' \leadsto x'_1 \leadsto x$.
Soit $\mathfrak p' \subset \mathcal{O}_{U, u}$ l'idéal premier correspondant
à $u'$. L'existence de ces spécialisations implique que
$\dim(\mathcal{O}_{U, u}/\mathfrak p') \geq 2$.
Ainsi, tout ouvert non vide de
$\Spec(\mathcal{O}_{U, u}/\mathfrak p')$ est infini d'après Algèbre, lemme
\ref{algebra-lemma-Noetherian-local-domain-dim-2-infinite-opens}.
D'après le Lemme \ref{lemma-decent-no-specializations-map-to-same-point},
on obtient une application continue
$$
\Spec(\mathcal{O}_{U, u}/\mathfrak p')
\setminus \{\mathfrak m_u/\mathfrak p'\}
\longrightarrow
W.
$$
Comme le point générique du membre de gauche s'envoie sur $x'$, l'image est
contenue dans $\overline{\{x'\}}$. On en conclut que l'image inverse de
$\{x'\}$ par la flèche affichée est un ouvert non vide, donc infini.
Cependant, les fibres de $U \to X$ sont finies puisque $X$ est décent ; la
fibre au-dessus de $x'$ serait donc infinie. Cette contradiction achève la
démonstration.
\end{proof}






\section{Irréductibilité locale}
\label{section-irreducible-local-ring}

\noindent
Nous avons déjà défini le nombre géométrique de branches d'un espace
algébrique en un point dans Propriétés des espaces, section
\ref{spaces-properties-section-irreducible-local-ring}.
Le nombre de branches d'un espace algébrique en un point ne peut être défini
que pour les espaces algébriques décents.

\begin{lemma}
\label{lemma-irreducible-local-ring}
Soit $S$ un schéma. Soit $X$ un espace algébrique décent sur $S$.
Soit $x \in |X|$ un point. Les conditions suivantes sont équivalentes :
\begin{enumerate}
\item pour tout voisinage étale élémentaire $(U, u) \to (X, x)$,
l'anneau local $\mathcal{O}_{U, u}$ possède un unique idéal premier minimal,
\item pour tout voisinage étale élémentaire $(U, u) \to (X, x)$,
il existe une unique composante irréductible de $U$ passant par $u$,
\item pour tout voisinage étale élémentaire $(U, u) \to (X, x)$,
l'anneau local $\mathcal{O}_{U, u}$ est unibranche,
\item l'anneau local hensélien $\mathcal{O}_{X, x}^h$ possède un unique idéal
premier minimal.
\end{enumerate}
\end{lemma}

\begin{proof}
L'équivalence de (1) et (2) résulte du fait que les composantes irréductibles
de $U$ passant par $u$ sont en correspondance bijective ($1$-$1$) avec les
premiers minimaux de l'anneau local de $U$ en $u$. L'anneau
$\mathcal{O}_{X, x}^h$ est l'hensélisé de $\mathcal{O}_{U, u}$ ; voir la
discussion qui suit la Définition \ref{definition-henselian-local-ring}.
En particulier, (3) et (4) sont équivalentes d'après Compléments d'algèbre,
lemme \ref{more-algebra-lemma-unibranch}.
L'équivalence de (2) et (3) découle de Compléments sur les morphismes, lemme
\ref{more-morphisms-lemma-nr-branches}.
\end{proof}

\begin{definition}
\label{definition-unibranch}
Soit $S$ un schéma. Soit $X$ un espace algébrique décent sur $S$.
Soit $x \in |X|$. On dit que $X$ est {\it unibranche en $x$} si les
conditions équivalentes du Lemme \ref{lemma-irreducible-local-ring}
sont satisfaites. On dit que $X$ est {\it unibranche} si $X$ est unibranche
en tout $x \in |X|$.
\end{definition}

\noindent
Cette définition est compatible avec celle donnée pour les schémas
(Propriétés, définition \ref{properties-definition-unibranch}).

\begin{lemma}
\label{lemma-nr-branches-local-ring}
Soit $S$ un schéma. Soit $X$ un espace algébrique décent sur $S$.
Soit $x \in |X|$ un point. Soit $n \in \{1, 2, \ldots\}$ un entier.
Les conditions suivantes sont équivalentes :
\begin{enumerate}
\item pour tout voisinage étale élémentaire $(U, u) \to (X, x)$, le nombre
d'idéaux premiers minimaux de l'anneau local $\mathcal{O}_{U, u}$ est
$\leq n$, et, pour au moins un choix de $(U, u)$, il est égal à $n$,
\item pour tout voisinage étale élémentaire $(U, u) \to (X, x)$, le nombre
de composantes irréductibles de $U$ passant par $u$ est $\leq n$, et, pour
au moins un choix de $(U, u)$, il est égal à $n$,
\item pour tout voisinage étale élémentaire $(U, u) \to (X, x)$, le nombre
de branches de $U$ en $u$ est $\leq n$, et, pour au moins un choix de
$(U, u)$, il est égal à $n$,
\item le nombre d'idéaux premiers minimaux de
$\mathcal{O}_{X, x}^h$ est $n$.
\end{enumerate}
\end{lemma}

\begin{proof}
L'équivalence de (1) et (2) résulte du fait que les composantes irréductibles
de $U$ passant par $u$ sont en correspondance bijective ($1$-$1$) avec les
premiers minimaux de l'anneau local de $U$ en $u$.
L'anneau $\mathcal{O}_{X, x}^h$ est l'hensélisé de $\mathcal{O}_{U, u}$ ;
voir la discussion qui suit la
Définition \ref{definition-henselian-local-ring}.
En particulier, (3) et (4) sont équivalentes d'après Compléments d'algèbre,
lemme \ref{more-algebra-lemma-unibranch}.
L'équivalence de (2) et (3) découle de Compléments sur les morphismes, lemme
\ref{more-morphisms-lemma-nr-branches}.
\end{proof}

\begin{definition}
\label{definition-number-of-branches}
Soit $S$ un schéma. Soit $X$ un espace algébrique décent sur $S$.
Soit $x \in |X|$. Le {\it nombre de branches de $X$ en $x$} est égal soit à
$n \in \mathbf{N}$ si les conditions équivalentes du
Lemme \ref{lemma-nr-branches-local-ring} sont satisfaites, soit à $\infty$.
\end{definition}






\section{Espaces algébriques caténaires}
\label{section-catenary}

\noindent
Cette section étend aux espaces algébriques les résultats de Propriétés,
section \ref{properties-section-catenary}, et de Morphismes, section
\ref{morphisms-section-universally-catenary}.

\begin{definition}
\label{definition-catenary}
Soit $S$ un schéma. Soit $X$ un espace algébrique décent sur $S$.
On dit que $X$ est {\it caténaire} si $|X|$ est caténaire
(Topologie, définition \ref{topology-definition-catenary}).
\end{definition}

\noindent
Si $X$ est représentable, cette définition équivaut à la notion
correspondante pour le schéma qui représente $X$.

\begin{lemma}
\label{lemma-scheme-with-dimension-function}
Soit $S$ un schéma localement noethérien et universellement caténaire.
Soit $\delta : S \to \mathbf{Z}$ une fonction de dimension.
Soit $X$ un espace algébrique décent sur $S$ tel que le morphisme structural
$X \to S$ soit localement de type fini. Soit
$\delta_X : |X| \to \mathbf{Z}$ l'application qui envoie $x$ sur
$\delta(f(x))$ augmenté du degré de transcendance de $x/f(x)$.
Alors $\delta_X$ est une fonction de dimension sur $|X|$.
\end{lemma}

\begin{proof}
Soit $\varphi : U \to X$ un morphisme étale surjectif, où $U$ est un schéma.
Alors la fonction $\delta_U$, définie de manière analogue, est une fonction
de dimension sur $U$ d'après Morphismes, lemme
\ref{morphisms-lemma-dimension-function-propagates}.
D'autre part, par définition du degré de transcendance relatif dans
Morphismes d'espaces, définition
\ref{spaces-morphisms-definition-dimension-fibre}, on a
$\delta_U(u) = \delta_X(\varphi(u))$.

\medskip\noindent
Soit $x \leadsto x'$ une spécialisation de points de $|X|$.
D'après le Lemme \ref{lemma-decent-specialization}, on peut trouver une
spécialisation $u \leadsto u'$ de points de $U$ telle que
$\varphi(u) = x$ et $\varphi(u') = x'$. De plus, $x = x'$ si et seulement si
$u = u'$ ; voir le
Lemme \ref{lemma-decent-no-specializations-map-to-same-point}.
Le fait que $\delta_U$ soit une fonction de dimension implique donc que
$\delta_X$ en est une ; voir Topologie, définition
\ref{topology-definition-dimension-function}.
\end{proof}

\begin{lemma}
\label{lemma-universally-catenary-scheme}
Soit $S$ un schéma localement noethérien et universellement caténaire.
Soit $X$ un espace algébrique sur $S$ tel que $X$ soit décent et que le
morphisme structural $X \to S$ soit localement de type fini. Alors $X$ est
caténaire.
\end{lemma}

\begin{proof}
La question est locale sur $S$ (utiliser Topologie, lemme
\ref{topology-lemma-catenary}).
On peut donc supposer que $S$ possède une fonction de dimension ; voir
Topologie, lemme \ref{topology-lemma-locally-dimension-function}.
D'après le Lemme \ref{lemma-scheme-with-dimension-function},
$|X|$ possède alors une fonction de dimension.
Comme $|X|$ est sobre (Proposition \ref{proposition-reasonable-sober}),
on conclut que $|X|$ est caténaire d'après Topologie, lemme
\ref{topology-lemma-dimension-function-catenary}.
\end{proof}

\noindent
D'après le Lemme \ref{lemma-universally-catenary-scheme}, la définition
suivante est compatible avec la notion déjà existante pour les espaces
algébriques représentables.

\begin{definition}
\label{definition-universally-catenary}
Soit $S$ un schéma. Soit $X$ un espace algébrique décent et localement
noethérien sur $S$. On dit que $X$ est {\it universellement caténaire} si,
pour tout morphisme $Y \to X$ d'espaces algébriques qui est localement de
type fini, avec $Y$ décent, l'espace algébrique $Y$ est caténaire.
\end{definition}

\noindent
Si $X$ est un espace algébrique, la condition « $X$ est décent et localement
noethérien » équivaut à la condition « $X$ est quasi-séparé et localement
noethérien ». C'est le
Lemme \ref{lemma-locally-Noetherian-decent-quasi-separated}.
Une autre manière de comprendre la définition précédente est donc que $X$
est universellement caténaire si et seulement si $Y$ est caténaire pour tout
morphisme $Y \to X$ quasi-séparé et localement de type fini.

\begin{lemma}
\label{lemma-universally-catenary}
Soit $S$ un schéma. Soit $X$ un espace algébrique décent, localement
noethérien et universellement caténaire sur $S$. Alors tout espace algébrique
décent localement de type fini sur $X$ est universellement caténaire.
\end{lemma}

\begin{proof}
Cela résulte formellement des définitions et du fait que les composés de
morphismes localement de type fini sont localement de type fini
(Morphismes d'espaces, lemme
\ref{spaces-morphisms-lemma-composition-finite-type}).
\end{proof}

\begin{lemma}
\label{lemma-check-dimension-function-finite-cover}
Soit $S$ un schéma. Soit $f : Y \to X$ un morphisme fini surjectif d'espaces
algébriques décents et localement noethériens. Soit
$\delta : |X| \to \mathbf{Z}$ une fonction. Si $\delta \circ |f|$ est une
fonction de dimension, alors $\delta$ est une fonction de dimension.
\end{lemma}

\begin{proof}
Soit $x \leadsto x'$, avec $x \not = x'$, une spécialisation dans $|X|$.
Choisissons $y \in |Y|$ tel que $|f|(y) = x$. Comme $|f|$ est fermée
(Morphismes d'espaces, lemme \ref{spaces-morphisms-lemma-finite-proper}),
il existe une spécialisation $y \leadsto y'$ telle que $|f|(y') = x'$.
On conclut donc que
$\delta(x) = \delta(|f|(y)) > \delta(|f|(y')) = \delta(x')$
(voir Topologie, définition \ref{topology-definition-dimension-function}).
Si $x \leadsto x'$ est une spécialisation immédiate, alors
$y \leadsto y'$ en est aussi une : en effet, si
$y \leadsto y'' \leadsto y'$, alors $|f|(y'')$ est nécessairement égal à
$x$ ou à $x'$, et il n'existe aucune spécialisation non triviale entre des
points d'une même fibre de $|f|$, d'après le
Lemme \ref{lemma-conditions-on-fibre-and-qf}.
\end{proof}

\noindent
La discussion se poursuivra dans Compléments sur les morphismes d'espaces,
section \ref{spaces-more-morphisms-section-catenary}.










\begin{multicols}{2}[\section{Autres chapitres}]
\noindent
Préliminaires
\begin{enumerate}
\item \hyperref[introduction-section-phantom]{Introduction}
\item \hyperref[conventions-section-phantom]{Conventions}
\item \hyperref[sets-section-phantom]{Théorie des ensembles}
\item \hyperref[categories-section-phantom]{Catégories}
\item \hyperref[topology-section-phantom]{Topologie}
\item \hyperref[sheaves-section-phantom]{Faisceaux sur les espaces}
\item \hyperref[sites-section-phantom]{Sites et faisceaux}
\item \hyperref[stacks-section-phantom]{Champs}
\item \hyperref[fields-section-phantom]{Corps}
\item \hyperref[algebra-section-phantom]{Algèbre commutative}
\item \hyperref[brauer-section-phantom]{Groupes de Brauer}
\item \hyperref[homology-section-phantom]{Algèbre homologique}
\item \hyperref[derived-section-phantom]{Catégories dérivées}
\item \hyperref[simplicial-section-phantom]{Méthodes simpliciales}
\item \hyperref[more-algebra-section-phantom]{Compléments d'algèbre}
\item \hyperref[smoothing-section-phantom]{Lissage des morphismes d'anneaux}
\item \hyperref[modules-section-phantom]{Faisceaux de modules}
\item \hyperref[sites-modules-section-phantom]{Modules sur les sites}
\item \hyperref[injectives-section-phantom]{Objets injectifs}
\item \hyperref[cohomology-section-phantom]{Cohomologie des faisceaux}
\item \hyperref[sites-cohomology-section-phantom]{Cohomologie sur les sites}
\item \hyperref[dga-section-phantom]{Algèbre différentielle graduée}
\item \hyperref[dpa-section-phantom]{Algèbre à puissances divisées}
\item \hyperref[sdga-section-phantom]{Faisceaux différentiels gradués}
\item \hyperref[hypercovering-section-phantom]{Hyperrecouvrements}
\end{enumerate}
Schémas
\begin{enumerate}
\setcounter{enumi}{25}
\item \hyperref[schemes-section-phantom]{Schémas}
\item \hyperref[constructions-section-phantom]{Constructions de schémas}
\item \hyperref[properties-section-phantom]{Propriétés des schémas}
\item \hyperref[morphisms-section-phantom]{Morphismes de schémas}
\item \hyperref[coherent-section-phantom]{Cohomologie des schémas}
\item \hyperref[divisors-section-phantom]{Diviseurs}
\item \hyperref[limits-section-phantom]{Limites de schémas}
\item \hyperref[varieties-section-phantom]{Variétés}
\item \hyperref[topologies-section-phantom]{Topologies sur les schémas}
\item \hyperref[descent-section-phantom]{Descente}
\item \hyperref[perfect-section-phantom]{Catégories dérivées de schémas}
\item \hyperref[more-morphisms-section-phantom]{Compléments sur les morphismes}
\item \hyperref[flat-section-phantom]{Compléments sur la platitude}
\item \hyperref[groupoids-section-phantom]{Schémas en groupoïdes}
\item \hyperref[more-groupoids-section-phantom]{Compléments sur les schémas en groupoïdes}
\item \hyperref[etale-section-phantom]{Morphismes étales de schémas}
\end{enumerate}
Thèmes de théorie des schémas
\begin{enumerate}
\setcounter{enumi}{41}
\item \hyperref[chow-section-phantom]{Homologie de Chow}
\item \hyperref[intersection-section-phantom]{Théorie de l'intersection}
\item \hyperref[pic-section-phantom]{Schémas de Picard des courbes}
\item \hyperref[weil-section-phantom]{Théories cohomologiques de Weil}
\item \hyperref[adequate-section-phantom]{Modules adéquats}
\item \hyperref[dualizing-section-phantom]{Complexes dualisants}
\item \hyperref[duality-section-phantom]{Dualité pour les schémas}
\item \hyperref[discriminant-section-phantom]{Discriminants et différentes}
\item \hyperref[derham-section-phantom]{Cohomologie de de Rham}
\item \hyperref[local-cohomology-section-phantom]{Cohomologie locale}
\item \hyperref[algebraization-section-phantom]{Géométrie algébrique et formelle}
\item \hyperref[curves-section-phantom]{Courbes algébriques}
\item \hyperref[resolve-section-phantom]{Résolution des surfaces}
\item \hyperref[models-section-phantom]{Réduction semi-stable}
\item \hyperref[functors-section-phantom]{Foncteurs et morphismes}
\item \hyperref[equiv-section-phantom]{Catégories dérivées de variétés}
\item \hyperref[pione-section-phantom]{Groupes fondamentaux de schémas}
\item \hyperref[etale-cohomology-section-phantom]{Cohomologie étale}
\item \hyperref[crystalline-section-phantom]{Cohomologie cristalline}
\item \hyperref[proetale-section-phantom]{Cohomologie pro-étale}
\item \hyperref[relative-cycles-section-phantom]{Cycles relatifs}
\item \hyperref[more-etale-section-phantom]{Compléments de cohomologie étale}
\item \hyperref[trace-section-phantom]{La formule des traces}
\end{enumerate}
Espaces algébriques
\begin{enumerate}
\setcounter{enumi}{64}
\item \hyperref[spaces-section-phantom]{Espaces algébriques}
\item \hyperref[spaces-properties-section-phantom]{Propriétés des espaces algébriques}
\item \hyperref[spaces-morphisms-section-phantom]{Morphismes d'espaces algébriques}
\item \hyperref[decent-spaces-section-phantom]{Espaces algébriques décents}
\item \hyperref[spaces-cohomology-section-phantom]{Cohomologie des espaces algébriques}
\item \hyperref[spaces-limits-section-phantom]{Limites d'espaces algébriques}
\item \hyperref[spaces-divisors-section-phantom]{Diviseurs sur les espaces algébriques}
\item \hyperref[spaces-over-fields-section-phantom]{Espaces algébriques sur des corps}
\item \hyperref[spaces-topologies-section-phantom]{Topologies sur les espaces algébriques}
\item \hyperref[spaces-descent-section-phantom]{Descente et espaces algébriques}
\item \hyperref[spaces-perfect-section-phantom]{Catégories dérivées d'espaces}
\item \hyperref[spaces-more-morphisms-section-phantom]{Compléments sur les morphismes d'espaces}
\item \hyperref[spaces-flat-section-phantom]{Platitude sur les espaces algébriques}
\item \hyperref[spaces-groupoids-section-phantom]{Groupoïdes dans les espaces algébriques}
\item \hyperref[spaces-more-groupoids-section-phantom]{Compléments sur les groupoïdes dans les espaces}
\item \hyperref[bootstrap-section-phantom]{Amorçage}
\item \hyperref[spaces-pushouts-section-phantom]{Sommes amalgamées d'espaces algébriques}
\end{enumerate}
Thèmes de géométrie
\begin{enumerate}
\setcounter{enumi}{81}
\item \hyperref[spaces-chow-section-phantom]{Groupes de Chow des espaces}
\item \hyperref[groupoids-quotients-section-phantom]{Quotients de groupoïdes}
\item \hyperref[spaces-more-cohomology-section-phantom]{Compléments sur la cohomologie des espaces}
\item \hyperref[spaces-simplicial-section-phantom]{Espaces simpliciaux}
\item \hyperref[spaces-duality-section-phantom]{Dualité pour les espaces}
\item \hyperref[formal-spaces-section-phantom]{Espaces algébriques formels}
\item \hyperref[restricted-section-phantom]{Algébrisation des espaces formels}
\item \hyperref[spaces-resolve-section-phantom]{Résolution des surfaces revisitée}
\end{enumerate}
Théorie des déformations
\begin{enumerate}
\setcounter{enumi}{89}
\item \hyperref[formal-defos-section-phantom]{Théorie formelle des déformations}
\item \hyperref[defos-section-phantom]{Théorie des déformations}
\item \hyperref[cotangent-section-phantom]{Le complexe cotangent}
\item \hyperref[examples-defos-section-phantom]{Problèmes de déformation}
\end{enumerate}
Champs algébriques
\begin{enumerate}
\setcounter{enumi}{93}
\item \hyperref[algebraic-section-phantom]{Champs algébriques}
\item \hyperref[examples-stacks-section-phantom]{Exemples de champs}
\item \hyperref[stacks-sheaves-section-phantom]{Faisceaux sur les champs algébriques}
\item \hyperref[criteria-section-phantom]{Critères de représentabilité}
\item \hyperref[artin-section-phantom]{Axiomes d'Artin}
\item \hyperref[quot-section-phantom]{Espaces de Quot et de Hilbert}
\item \hyperref[stacks-properties-section-phantom]{Propriétés des champs algébriques}
\item \hyperref[stacks-morphisms-section-phantom]{Morphismes de champs algébriques}
\item \hyperref[stacks-limits-section-phantom]{Limites de champs algébriques}
\item \hyperref[stacks-cohomology-section-phantom]{Cohomologie des champs algébriques}
\item \hyperref[stacks-perfect-section-phantom]{Catégories dérivées de champs}
\item \hyperref[stacks-introduction-section-phantom]{Introduction aux champs algébriques}
\item \hyperref[stacks-more-morphisms-section-phantom]{Compléments sur les morphismes de champs}
\item \hyperref[stacks-geometry-section-phantom]{La géométrie des champs}
\end{enumerate}
Thèmes de théorie des espaces de modules
\begin{enumerate}
\setcounter{enumi}{107}
\item \hyperref[moduli-section-phantom]{Champs de modules}
\item \hyperref[moduli-curves-section-phantom]{Espaces de modules de courbes}
\end{enumerate}
Divers
\begin{enumerate}
\setcounter{enumi}{109}
\item \hyperref[examples-section-phantom]{Exemples}
\item \hyperref[exercises-section-phantom]{Exercices}
\item \hyperref[guide-section-phantom]{Guide bibliographique}
\item \hyperref[desirables-section-phantom]{Desiderata}
\item \hyperref[coding-section-phantom]{Style de programmation}
\item \hyperref[obsolete-section-phantom]{Obsolète}
\item \hyperref[fdl-section-phantom]{Licence de documentation libre GNU}
\item \hyperref[index-section-phantom]{Index généré automatiquement}
\end{enumerate}
\end{multicols}


\bibliography{my}
\bibliographystyle{amsalpha}

\end{document}
